% !TeX root = main-vnmm-mlpg-eigenvalues.tex
\documentclass[VANCOUVER]{USG}

\graphicspath{{./}{./images/}}

\articletype{Research Article}

\received{00}
\revised{00}
\accepted{00}
\journal{International Journal for Numerical Methods in Engineering}
\volume{00}
\copyyear{2026}
\startpage{1}
\articledoi{10.1002/0000}

\raggedbottom



\begin{document}

\title{A Truly Cell-Free Vector Nodal Meshless Method via Meshless Local Petrov--Galerkin (VNMM-MLPG) for 2D Electromagnetic Eigenvalue Problems}

\author[1]{Renato C. Mesquita*}
\author[2]{Williams L. de Nicomedes}
\author[3]{B\protect\'{a}rbara M. F. Gon\protect\c{c}alves}
\author[4]{Luilly A. G. Ortiz}

\authormark{MESQUITA \textsc{et al.}}
\titlemark{CELL-FREE VNMM-MLPG FOR 2D MAXWELL EIGENPROBLEMS}

\address[1]{\orgdiv{Department of Electrical Engineering}, \orgname{Universidade Federal de Minas Gerais}, \orgaddress{\city{Belo Horizonte}, \state{MG}, \country{Brazil}}}
\address[2]{\orgdiv{Department of Engineering and Computing}, \orgname{Instituto Federal de Minas Gerais, Campus Bambu\'{i}}, \orgaddress{\city{Bambu\'{i}}, \state{MG}, \country{Brazil}}}
\address[3]{\orgdiv{Mechanical Department}, \orgname{Instituto Federal de Minas Gerais, Campus Congonhas}, \orgaddress{\city{Congonhas}, \state{MG}, \country{Brazil}}}
\address[4]{\orgdiv{Industrial Control and Processes Department}, \orgname{Instituto Federal de Minas Gerais, Campus Sabar\'{a}}, \orgaddress{\city{Sabar\'{a}}, \state{MG}, \country{Brazil}}}

\corres{Renato C. Mesquita, Department of Electrical Engineering, Universidade Federal de Minas Gerais, Av. Ant\^{o}nio Carlos, 6627, Belo Horizonte, MG, 31270-901, Brazil. \email{renato@ufmg.br}}

\fundingInfo{Brazilian National Council for Scientific and Technological Development (CNPq); Minas Gerais State Research Foundation (FAPEMIG); Coordination for the Improvement of Higher Education Personnel (CAPES).}

\abstract[ABSTRACT]{Source-free electromagnetic eigenvalue problems in $H(\mathbf{curl})$ govern the resonant frequencies and modal fields of microwave cavities, waveguide filters, and RF resonators. However, meshless vector discretizations of these problems remain challenging due to spurious gradient modes and background quadrature overhead. While the Vector Nodal Meshless Method (VNMM) constructs node-based vector approximations on unstructured point clouds via directional nodal degrees of freedom, prior implementations relied on Cartesian background integration cells, introducing intra-cell stencil discontinuities and limiting meshless autonomy. This paper presents a truly cell-free VNMM via Meshless Local Petrov--Galerkin (VNMM-MLPG) for 2D electromagnetic eigenvalue problems. Combining linear ($\mathcal{P}^1$) vector point interpolation with weak forms over autonomous circular subdomains eliminates background grids. Because fixed linear stencils yield constant curl and divergence per disk, Stokes' and Gauss' theorems show that boundary-vanishing weights (MLPG1) produce zero stiffness, motivating a Local Galerkin (MLPG6) scheme. To avoid the asymmetry and complex spectra of row-by-row MLPG without lens-intersection quadratures, we formulate a closed-form overlapping \textit{Symmetric Subdomain Galerkin} assembly yielding a real symmetric positive-semidefinite stiffness $\mathbf{K}$ and positive-definite mass $\mathbf{M}$, whose continuous multiplicity-weighted Rayleigh quotient homogenizes at second order on regular rectangular lattices. Modal decomposition further reveals two harmonic linear gradient modes in the local div-curl kernel, which are controlled via a VMS-inspired orthogonal energy stabilization paired with a resolution-scaled divergence penalty. Following auxiliary local-equilibrium verification tests, 2D PEC cavity benchmarks show zero spurious modes across the tested 10-mode band, monotonic convergence down to $0.78\%\text{--}0.79\%$ error with empirical cutoff rates of $p = 1.72$ (regular) and $p = 1.61$ ($25\%$ jitter) under linear stabilization, and $20\%\text{--}24\%$ sparser operators assembled 2.0$\times$--3.7$\times$ faster than calibrated cell-based VNMM-Gauss.}

\keywords{Vector Nodal Meshless Method (VNMM) | Meshless Local Petrov--Galerkin (MLPG) | Point Interpolation Method (PIM) | Truly cell-free methods | Maxwell eigenvalue problems | Div-curl regularization | Orthogonal energy stabilization | Curl-curl vector discretization}

\maketitle



%\linenumbers

%% =========================================================================
%% SECTION 1: INTRODUCTION
%% =========================================================================
\section{Introduction}
\label{sec:introduction}

Source-free electromagnetic eigenvalue problems governed by Maxwell's curl-curl equations in $H(\mathbf{curl})$ play a foundational role in computational electromagnetics (CEM) \cite{monk2003finite, jin2014finite, boffi2006finite}. Accurate computation of resonant frequencies, cutoff wavenumbers, and modal field profiles is essential across a wide range of engineering applications, including microwave and radio-frequency (RF) resonant cavities, waveguide filters and multiplexers, modal analysis of metallic and dielectric waveguides, particle-accelerator RF cavities, photonic resonators, and cavity-based material characterization. In many of these applications---particularly during parametric sweeps, shape sensitivity analysis, and structural or topology optimization of resonant devices---the cavity or waveguide geometry evolves iteratively, making repeated mesh regeneration a dominant computational bottleneck.

Traditionally, Maxwell eigenvalue problems have been solved using the finite element method (FEM) with conforming edge elements \cite{nedelec1980mixed, monk2003finite, jin2014finite}. Edge elements assign tangential field components to elemental edges, enforcing tangential trace continuity while naturally accommodating normal jumps across dielectric interfaces and satisfying the exact sequence of the discrete de Rham complex. Consequently, irrotational gradient fields reside in the exact discrete nullspace of the curl operator, $\mathrm{grad} \, W_h \subset \ker(\mathrm{curl}_h)$, preventing spectral pollution by non-physical spurious eigenmodes \cite{bossavit1998computational, boffi2006finite, boffi2013mixed}. Nevertheless, edge elements strictly depend on generating and maintaining conforming simplex meshes (triangles in 2D, tetrahedra in 3D), which motivates the development of truly meshless vector formulations capable of resolving electromagnetic spectra directly on point clouds.

To circumvent mesh generation bottlenecks, meshless and meshfree methods have garnered sustained interest across solid mechanics, fluid dynamics, and wave physics \cite{belytschko2023meshfree, liu2009mesh, fasshauer2007meshfree, atluri2004meshless}. In computational electromagnetics, early meshless adaptations employed scalar formulations for electric scalar potentials, magnetic vector potentials, or out-of-plane field components governed by scalar Helmholtz equations \cite{viana1999moving, ho2001moving, parreira2006element, nicomedes2012meshless}. Extending meshless concepts to vector Maxwell eigenvalue problems proved substantially more challenging. Early vector adaptations attempted to mimic finite element edge conformity by assigning degrees of freedom to pairs of nodes forming virtual edges, as in the Edge Meshless Method (EMM) \cite{lima2017edge}. While effective at suppressing spurious eigenmodes, such schemes reintroduce topological connectivity between node pairs, partially compromising the autonomy of the point cloud. Conversely, purely nodal collocated schemes with vector degrees of freedom typically induce severe spurious-mode pollution and derivative aliasing due to the lack of discrete de Rham complex exactness on arbitrary point clouds \cite{bossavit1990spurious, boffi2006finite}.

The Vector Nodal Meshless Method (VNMM), introduced by Ortiz et al.~\cite{ortiz2021nodal, ortiz2023tese} and Gon\c{c}alves et al.~\cite{goncalves2023vector, goncalves2025tese}, provides a node-based vector interpolation strategy. Fundamentally, the VNMM can be recognized as the asymptotic limit of the EMM as the virtual edge length tends to zero. Collapsing the node pair into a single evaluation point $\mathbf{x}_I$ reduces the edge circulation to a local directional projection along a unit director $\vec{t}_I$, restoring point autonomy without pairwise connectivity. From an interpolation standpoint, the VNMM constructs shape functions via local polynomial point collocation on support stencils, conceptually generalizing the Point Interpolation Method (PIM) \cite{liu2001point, liu2005introduction} to vector fields endowed with directional degrees of freedom $e_I = \vec{E}(\mathbf{x}_I) \cdot \vec{t}_I$.

In a recent study \cite{stable_vnmm2d_2026}, the mathematical foundations of the 2D VNMM were systematically examined on unconstrained irregular point clouds. That study resolved several fundamental questions:
\begin{enumerate}
    \item \textbf{Incomplete Basis Stagnation:} The classical incomplete solenoidal vector basis ($\mathcal{L}^1$, consisting of 3 terms per stencil) was shown to suffer from a persistent $\mathcal{O}(1)$ error plateau in curl approximation on irregular clouds. This plateau arose from derivative aliasing between curl and divergence components (divergence modes of the field gradient tensor \cite{stable_vnmm2d_2026}). This was resolved by establishing the complete linear vector space $\mathcal{P}_1 \times \mathcal{P}_1$ (requiring a 6-node stencil, denoted $\mathcal{P}^1$), which restored full linear completeness and $\mathcal{O}(h)$ curl convergence.
    \item \textbf{Geometric Singularity Mitigation:} A scale-invariant determinant qualification criterion on a dimensionless placement matrix $\widehat{\mathbf{A}}$, coupled with adaptive $k$-d tree searches, was developed to automatically reject nearly singular node and director configurations without empirical parameter tuning.
    \item \textbf{Variational Div-Curl Stabilization:} Because dynamic support stencils across point clouds lack global trace conformity, discrete gradients leak into the positive curl spectrum. A regularized cell-based Galerkin formulation with a volumetric div-curl penalty ($s_{\mathrm{div}}$) was introduced, successfully shifting spurious modes to higher frequencies while preserving physical cavity resonances.
\end{enumerate}

Despite these advances, the formulation reported in \cite{stable_vnmm2d_2026} retained an important limitation: \textbf{the reliance on a Cartesian background integration grid}. As classified by Belytschko et al.~\cite{belytschko1994element} and Atluri and Shen \cite{atluri2002meshless}, methods like the Element-Free Galerkin (EFG) are not \textit{truly cell-free} because evaluation of the global weak form requires an auxiliary background mesh of quadrature cells covering the domain $\Omega$. In the VNMM context, this background scaffolding introduces two main numerical challenges:
\begin{itemize}
    \item \textbf{Intra-Cell Stencil Discontinuities:} In cell-based integration, the 6-node support stencil is selected point-by-point at each Gauss integration point $P_g$. Because nearest-neighbor stencils switch abruptly within the cell interior, the resulting shape functions and derivatives are piecewise discontinuous across intra-cell boundaries. As demonstrated in \cite{stable_vnmm2d_2026}, under-integrating these cells (e.g., using $2 \times 2$ Gauss points) produces artificial numerical softening that conceals rank deficiencies in the discrete curl operator. Suppressing this artifact required a dense, resolution-matched quadrature ($3 \times 3$ Gauss points per cell, exceeding 3,600 quadrature evaluations for a 441-node grid).
    \item \textbf{Loss of True Meshless Autonomy:} Background integration grids do not demand the strict shape-regularity or aspect-ratio metrics of finite element meshes. However, constructing background cells and trimming cut cells along complex, curved, or 3D boundaries still introduces geometric pre-processing overhead and quadrature complexity, limiting meshless autonomy.
\end{itemize}

Several alternative domain integration strategies have been investigated in the broader meshfree literature \cite{hillman2018performance}, but each presents specific drawbacks when applied to unstructured vector problems:
\begin{itemize}
    \item \textbf{Direct Nodal Integration (DNI) and Rank Instabilities:} The simplest meshless strategy evaluates the global weak form directly at the nodes. However, as demonstrated by Babu{\v{s}}ka et al.~\cite{babuska2008quadrature} and Hillman and Chen~\cite{hillman2018performance}, DNI underestimates the stored electromagnetic energy of low-order modes, leading to rank deficiencies (sawtooth zero-energy modes) and degraded convergence rates without extensive ad-hoc stabilization.
    \item \textbf{Conforming vs. Non-Conforming Nodal Smoothing (SCNI and SNNI):} Stabilized Conforming Nodal Integration (SCNI) \cite{chen2001stabilized} restores stability and patch test exactness by computing smoothed gradients via boundary integrals over conforming Voronoi dual cells, preserving matrix symmetry \cite{hillman2018performance}. However, constructing Voronoi diagrams on irregular point clouds is geometrically brittle, sensitive to boundary slivers, and computationally demanding in 3D. While Stabilized Non-Conforming Nodal Integration (SNNI) dispenses with Voronoi meshing by using non-conforming local boxes \cite{hillman2018performance}, the loss of partition conformity fails the patch test, necessitating Petrov--Galerkin variational consistency (VC) corrections \cite{chen2013arbitrary, hillman2018performance} that require auxiliary local equation solves.
    \item \textbf{Strong-Form Nodal Collocation:} Enforcing the governing PDE directly at the nodes avoids numerical quadrature \cite{hillman2018performance}. However, for second-order PDEs, collocation requires computing high-order derivatives of shape functions and demands at least quadratic polynomial completeness for convergence \cite{hillman2018performance}. More fundamentally, for vector electromagnetic problems governed by Maxwell's curl-curl equations, strong-form collocation lacks a discrete de Rham complex structure, triggering spectral corruption and unphysical gradient spurious modes \cite{bossavit1990spurious, boffi2006finite}.
    \item \textbf{Discontinuous Galerkin with Trace Jump Penalties:} One could view shifting stencils as non-conforming elements and introduce interface integrals with consistent fluxes and tangential trace jump penalties, as in interior penalty discontinuous Galerkin (IPDG) formulations for electromagnetics \cite{houston2004mixed, hesthaven2004high}. For pointwise-selected stencils evaluated independently at each quadrature point (as in Gauss-VNMM), however, the resulting sub-cell discontinuity skeleton is implicit and geometrically ambiguous to track; conversely, confining stencil transitions to explicit cell faces to enable IPDG fluxes retains a background cell partition rather than achieving a truly cell-free formulation.
    \item \textbf{Autonomous Local Weak Forms (MLPG and LPIM):} Formulations based on local subdomains, such as the Meshless Local Petrov--Galerkin (MLPG) method of Atluri and Zhu \cite{atluri1998new, atluri2002meshless} and the Local Point Interpolation Method (LPIM) of Liu and Gu \cite{liu2001local, liu2005introduction}, avoid background cells by establishing weak forms over autonomous circular disks or spheres $\Omega_s^I$ centered at each node.
\end{itemize}

While MLPG and LPIM have been widely applied to scalar mechanics, acoustics, and electromagnetics \cite{sladek2001local, atluri2002meshless, liu2001local, nicomedes2012meshless}, extending them to directional $\mathcal{P}^1$ vector formulations raises four challenges. First, boundary-vanishing test functions (MLPG1/LRPIM \cite{atluri1998new, atluri2002meshless, liu2005introduction}) collapse the local stiffness of constant $\mathcal{P}^1$ derivatives to zero. Second, classical row-by-row assembly ($\Omega_s^I \ne \Omega_s^J$) \cite{atluri2002meshless, liu2001local, nicomedes2012meshless} yields unsymmetric operators risking complex spectra \cite{saad2011numerical, boffi2013mixed}, whereas existing symmetric overlapping methods such as the Method of Finite Spheres \cite{de2000method, nicomedes2017meshfree, nicomedes2021acoustic} rely on pairwise lens-intersection cubatures. Third, boundary intersections truncate disks into circular segments $\Omega_s^B = \Omega_s^I \cap \Omega$ while requiring preservation of VNMM directional Kronecker-delta collocation \cite{ortiz2021nodal, goncalves2023vector, stable_vnmm2d_2026}. Fourth, single-disk $\mathcal{P}^1$ integration caps local stiffness rank at 2 (versus mass rank 6), leaving two harmonic gradient modes unpenalized.

To address these challenges, we formulate and validate \textbf{VNMM-MLPG}---a truly cell-free vector nodal meshless method for source-free 2D electromagnetic eigenvalue problems---combining directional $\mathcal{P}^1$ VNMM interpolation \cite{ortiz2021nodal, stable_vnmm2d_2026} with an overlapping \textbf{Symmetric Subdomain Galerkin (SSG)} Local Galerkin (MLPG6) assembly \cite{atluri2002meshless}. The main contributions are:
\begin{itemize}
    \item \textbf{Closed-Form Subdomain Operators:} Local curl, divergence, and mass matrices factor in exact closed form via $\boldsymbol{\beta}^I = \mathbf{A}_I^{-1}$ and analytical geometric moments over interior disks and boundary-truncated circular segments, avoiding numerical and lens-intersection quadratures.
    \item \textbf{Symmetric Subdomain Galerkin Assembly:} An additive subdomain energy assembly produces a symmetric positive-semidefinite stiffness $\mathbf{K}$ and a positive-definite mass $\mathbf{M}$, guarantees fixed-mesh PEC positive definiteness under connected stencils, and avoids the refinement instability observed in unsymmetric row-by-row MLPG/LPIM.
    \item \textbf{Orthogonal Energy Stabilization:} Mapping the harmonic linear gradient modes ($\mathbf{v}_{\mathrm{shear}}^I, \mathbf{v}_{\mathrm{hyp}}^I$) onto nodal DoFs and pairing a VMS-inspired linear stabilization ($\alpha_{\mathrm{stab}} = c_{\mathrm{stab}}(h_{\mathrm{char}}/L)$, baseline $c_{\mathrm{stab}} = 4.0$) with superlinear divergence scaling ($s_{\mathrm{div}} \propto (h_{\mathrm{char}}/L)^{-1.5}$) restores rank 4 over the gradient subspace. This prevents fine-mesh overshoot and yields monotonic decay to $0.78\%\text{--}0.79\%$ error across $15\times 15$ to $81\times 81$ grids ($p = 1.72$ regular, $p = 1.61$ at $25\%$ jitter).
    \item \textbf{Unified Truncated-Disk Boundary Treatment:} Directional Kronecker-delta collocation enforces tangential Dirichlet/PEC constraints via direct row/column elimination without Lagrange multipliers or Nitsche penalties, while natural Neumann (PMC) data enter via 1D segment integrals.
    \item \textbf{Verification and Computational Efficiency:} Local equilibrium tests and subdomain sweeps ($\alpha_s = 0.80$ \cite{atluri2002meshless, liu2009mesh} within the operational band $\alpha_s \in [0.70, 1.20]$) establish baseline reliability. On PEC cavity benchmarks, VNMM-MLPG achieves 2.0$\times$--3.7$\times$ faster assembly, $20\%\text{--}24\%$ sparser operators, and lower cutoff errors on most perturbed clouds (five of six matched grids) relative to calibrated VNMM-Gauss, while sustaining monotonic fine-mesh convergence to sub-$0.8\%$ error at $81\times 81$ with clean modal field reconstructions.
\end{itemize}

The paper is organized as follows. Section~\ref{sec:formulation} develops the VNMM-MLPG formulation, analytical MLPG6 operators, SSG assembly, boundary treatment, and orthogonal energy stabilization (with variational consistency and error decomposition deferred to Appendix~\ref{subsec:variational_consistency_strang}). Section~\ref{sec:results} presents the six-stage numerical verification, spectral cavity validation, and VNMM-Gauss benchmark. Sections~\ref{sec:discussion} and~\ref{sec:conclusions} summarize advantages, open challenges, and conclusions.

%% =========================================================================
%% SECTION 2: THEORETICAL FORMULATION OF VNMM-MLPG
%% =========================================================================
\section{Theoretical Formulation of the VNMM-MLPG}
\label{sec:formulation}

\subsection{Continuous Boundary-Value Problem}
\label{subsec:continuous_bvp}

Consider a bounded two-dimensional domain $\Omega \subset \mathbb{R}^2$ with a Lipschitz boundary $\partial\Omega$, occupied by a homogeneous, linear, isotropic, non-dispersive medium characterized by constant relative permeability $\mu_r$ and permittivity $\epsilon_r$. The time-harmonic electric field $\vec{E}(x, y) = [E_x, E_y]^T$ satisfies Maxwell's curl-curl equation. To shift gradient-type spurious modes outside the operational frequency band on non-conforming meshless spaces, we adopt the div-curl regularized vector Helmholtz equation \cite{stable_vnmm2d_2026}:
\begin{equation}
    \nabla \times \left(\frac{1}{\mu_r} \nabla \times \vec{E}\right) - \frac{s_{\mathrm{div}}}{\mu_r} \nabla(\nabla \cdot \vec{E}) - k_0^2 \epsilon_r \vec{E} = \vec{J}_s \quad \text{in } \Omega,
    \label{eq:helmholtz_pde}
\end{equation}
where $k_0 = \omega\sqrt{\mu_0\epsilon_0}$ is the free-space wavenumber, $s_{\mathrm{div}} \ge 0$ is the dimensionless divergence regularization parameter, and $\vec{J}_s$ is an applied electric current density. The primary target of the present cell-free formulation is the source-free electromagnetic eigenvalue problem ($\vec{J}_s \equiv \mathbf{0}$ with homogeneous PEC/PMC walls); the source term $\vec{J}_s$ and inhomogeneous boundary data are retained in \eqref{eq:helmholtz_pde} for the controlled local-equilibrium verification tests in Stages~2--3 and to state the general natural boundary structure. In two dimensions with in-plane vector fields $\vec{E} = [E_x, E_y]^T$, the differential operators are defined by:
\begin{equation}
    \nabla \times \vec{E} = \left(\frac{\partial E_y}{\partial x} - \frac{\partial E_x}{\partial y}\right) \hat{\mathbf{z}}, \qquad \nabla \cdot \vec{E} = \frac{\partial E_x}{\partial x} + \frac{\partial E_y}{\partial y}.
\end{equation}

To establish a well-posed boundary-value problem, the physical boundary is partitioned into pairwise disjoint segments $\partial\Omega = \Gamma_D \cup \Gamma_N$ (with $\Gamma_D \cap \Gamma_N = \emptyset$), where $\vec{n}$ and $\vec{t}$ denote the outward unit normal and unit tangent to $\partial\Omega$, respectively:
\begin{enumerate}
    \item \textbf{Dirichlet Boundary Conditions ($\Gamma_D$):} The tangential trace of the electric field is prescribed:
          \begin{equation}
              \vec{n} \times \vec{E} = \vec{g}_D \quad \iff \quad \vec{E} \cdot \vec{t} = g_D \quad \text{on } \Gamma_D.
              \label{eq:pec_bc}
          \end{equation}
          When $g_D = 0$, this represents the homogeneous Dirichlet condition on perfectly conducting (PEC) electric walls ($\Gamma_{\mathrm{PEC}}$). When $g_D \ne 0$, it models inhomogeneous Dirichlet excitations used in the verification tests of numerical stages~2--3.
    \item \textbf{Neumann Boundary Conditions ($\Gamma_N$):} The natural boundary flux (tangential magnetic field excitation) is specified:
          \begin{equation}
              -\frac{1}{\mu_r} \left(\vec{n} \times (\nabla \times \vec{E})\right) = \vec{g}_N \quad \text{on } \Gamma_N.
              \label{eq:neumann_bvp_pde}
          \end{equation}
          When $\vec{g}_N = \mathbf{0}$, this represents the homogeneous Neumann condition on perfect magnetic conductors (PMC) or electromagnetic symmetry planes. When $\vec{g}_N \ne \mathbf{0}$, under the $e^{j\omega t}$ time-harmonic convention where $\nabla \times \vec{E} = -j\omega\mu_0\mu_r\vec{H}$, it specifies the impressed tangential magnetic flux $\vec{g}_N = j\omega\mu_0 (\vec{n} \times \vec{H})$.
\end{enumerate}

\subsection{Nodal Discretization and Local Support Geometries}
\label{subsec:nodal_discretization}

The domain $\overline{\Omega} = \Omega \cup \partial\Omega$ is discretized by a cloud of $N$ nodes $\mathcal{X}_N = \{P_I\}_{I=1}^N$ with Cartesian coordinates $\mathbf{x}_I = [x_I, y_I]^T \in \overline{\Omega}$. Each node $P_I$ is endowed with an assigned unit director vector $\vec{t}_I = [t_{x,I}, t_{y,I}]^T \in \mathbb{S}^1$, defining a scalar degree of freedom (DoF):
\begin{equation}
    e_I = \vec{E}(\mathbf{x}_I) \cdot \vec{t}_I.
\end{equation}
Along physical boundaries with prescribed tangential electric fields, directors are oriented tangent to the boundary ($\vec{t}_I \parallel \partial\Omega$). On perfectly conducting walls ($\mathbf{x}_I \in \Gamma_{\mathrm{PEC}}$), condition \eqref{eq:pec_bc} then translates directly into the homogeneous constraint $e_I = 0$, whereas on inhomogeneous Dirichlet boundaries ($\Gamma_D$), it permits direct collocation of the prescribed tangential profile ($e_I = g_D(\mathbf{x}_I)$). Likewise, along material interfaces between distinct media, aligning directors parallel to the interface can, in principle, preserve tangential electric-field continuity while accommodating physical jump discontinuities in the normal component (an extension left for future work; Section~\ref{subsec:deficiencies_challenges}).

In the MLPG framework, two distinct local domains are defined for each node $P_I$ (Fig.~\ref{fig:schematic_mlpg}):
\begin{enumerate}
    \item \textbf{Trial Support Domain ($\Omega_u^I$):} The spatial neighborhood containing an admissible stencil of $K = 6$ nodes used to construct the local trial vector field $\vec{E}_{I,h}(\mathbf{x})$.
    \item \textbf{Local Integration Subdomain ($\Omega_s^I$):} A circular disk of radius $r_s^I$ centered at $\mathbf{x}_I$:
          \begin{equation}
              \Omega_s^I = \left\{ \mathbf{x} \in \mathbb{R}^2 : \|\mathbf{x} - \mathbf{x}_I\| \le r_s^I \right\}, \qquad r_s^I = \alpha_s h_I,
              \label{eq:subdomain_disk}
          \end{equation}
          where $h_I \equiv \min_{J \ne I} \|\mathbf{x}_J - \mathbf{x}_I\|_2$ denotes the local nearest-neighbor spacing (matching $h_{\mathrm{local}}$ in \cite{stable_vnmm2d_2026}), and $\alpha_s > 0$ is a dimensionless radius scaling factor. For a rectangular domain of dimensions $L_x \times L_y$ (reference length $L = \max(L_x, L_y)$) discretized on an $N_x \times N_y$ grid, we distinguish three spacing metrics: local spacing $h_I$, nominal reference spacing $h_{\mathrm{nom}} = L/(N_x - 1)$, and domain-averaged spacing $\overline{h} = \frac{1}{N}\sum_{I=1}^N h_I$ on perturbed clouds. These are unified into a single characteristic spacing $h_{\mathrm{char}}$, equal to $h_{\mathrm{nom}}$ on regular grids and $\overline{h}$ on disordered clouds. We adopt the baseline value $\alpha_s = 0.80$, whose coverage ($100\%$ on regular lattices and $94.4\%$ under $25\%$ jitter) and stencil-containment bounds are analyzed in Section~\ref{subsec:stage4_alpha_s}.
\end{enumerate}

A representative nodal distribution with $25\%$ jitter is depicted in Fig.~\ref{fig:schematic_mlpg}, where interior directors $\vec{t}_k$ are randomly oriented and boundary directors satisfy $\vec{t}_k \parallel \partial\Omega$, alongside the 6-node trial stencils $\Omega_u$ (Section~\ref{subsec:stencil_selection}). Depending on node position relative to $\partial\Omega$, three subdomain configurations arise:
\begin{itemize}
    \item \textbf{Case 1 (Fully Interior, $\Omega_s^{I_1} \subset \Omega$):} When $\mathrm{dist}(\mathbf{x}_{I_1}, \partial\Omega) \ge r_s^{I_1}$, the integration domain is a full circular disk.
    \item \textbf{Case 2 (Truncated Interior, $\Omega_s^{I_2} \cap \Omega$):} When $\mathrm{dist}(\mathbf{x}_{I_2}, \partial\Omega) < r_s^{I_2}$, the disk crosses $\partial\Omega$ and is truncated to a circular segment bounded by an internal arc $L_{sI_2} = \partial\Omega_s^{I_2} \cap \Omega$ and a straight boundary segment $\Gamma_s^{I_2} = \Omega_s^{I_2} \cap \partial\Omega$.
    \item \textbf{Case 3 (Truncated Boundary, $\Omega_s^B \cap \Omega$):} For $\mathbf{x}_B \in \partial\Omega$, the disk is truncated into a half-disk (or corner sector) bounded by the internal arc $L_{sB} = \partial\Omega_s^B \cap \Omega$ and the boundary segment $\Gamma_s^B = \Omega_s^B \cap \partial\Omega$ where boundary conditions are enforced.
\end{itemize}

\begin{figure}[!t]
    \centering
    \includegraphics[width=0.72\textwidth]{fig_mlpg_subdomains_schematic.pdf}
    \caption{Schematic of the VNMM-MLPG nodal cloud, local integration subdomains $\Omega_s$ (drawn with an enlarged radius factor $\alpha_s = 1.10$ for clear visual illustration of near-boundary Case~2 truncation), and 6-node trial support stencils $\Omega_u$.}
    \label{fig:schematic_mlpg}
\end{figure}

\subsection{VNMM Vector Shape Functions and Field Interpolation via $\mathcal{P}^1$}
\label{subsec:trial_space}

Following \cite{stable_vnmm2d_2026}, within the local evaluation domain of node $P_I$, the vector field $\vec{E}(\mathbf{x})$ is approximated in the complete linear polynomial vector space $\mathcal{P}^1 = \mathcal{P}_1 \times \mathcal{P}_1$. This space is spanned by the basis matrix $\mathbf{P}(\Delta\mathbf{x}) \in \mathbb{R}^{2 \times 6}$:
\begin{equation}
    \mathbf{P}(\Delta\mathbf{x}) = \begin{bmatrix}
        1 & 0 & \Delta x & \Delta y & 0        & 0        \\
        0 & 1 & 0        & 0        & \Delta x & \Delta y
    \end{bmatrix},
    \label{eq:P_basis_matrix}
\end{equation}
with local shifted coordinates $\Delta\mathbf{x} = [\Delta x, \Delta y]^T \equiv \mathbf{x} - \mathbf{x}_I$ centered at node $P_I$, so that $\mathbf{P}(\Delta\mathbf{x}) \equiv \mathbf{P}(\mathbf{x} - \mathbf{x}_I)$.
Each local vector shape function $\vec{N}_a^I(\mathbf{x})$ associated with the $a$-th node of the support stencil $\Omega_u^I$ ($a = 1, \dots, 6$) is expressed directly as a linear combination of these basis terms:
\begin{equation}
    \vec{N}_a^I(\mathbf{x}) = \mathbf{P}(\mathbf{x} - \mathbf{x}_I) \boldsymbol{\beta}_a^I,
    \label{eq:P1_shape_function_expansion_mlpg}
\end{equation}
where $\boldsymbol{\beta}_a^I = [\beta_{1a}^I, \dots, \beta_{6a}^I]^T \in \mathbb{R}^6$ collects the polynomial coefficients for shape function $a$. Enforcing the directional Kronecker-delta projection property \cite{stable_vnmm2d_2026}:
\begin{equation}
    \vec{N}_a^I(\mathbf{x}_k) \cdot \vec{t}_k = \delta_{ak}, \qquad k = 1, \dots, 6,
    \label{eq:kronecker_projection}
\end{equation}
across all 6 stencil nodes yields the local linear placement system $\mathbf{A}_I \boldsymbol{\beta}_a^I = \boldsymbol{\delta}_a$, where $\boldsymbol{\delta}_a = [\delta_{1a}, \dots, \delta_{6a}]^T$ is the $a$-th canonical basis vector of $\mathbb{R}^6$, $(\Delta x_k, \Delta y_k) \equiv (x_k - x_I, y_k - y_I)$, and the $k$-th row of the moment placement matrix $\mathbf{A}_I \in \mathbb{R}^{6 \times 6}$ is:
\begin{equation}
    (\mathbf{A}_I)_{k, :} = \vec{t}_k^T \mathbf{P}(\mathbf{x}_k - \mathbf{x}_I) = \begin{bmatrix} t_{kx} & t_{ky} & \Delta x_k t_{kx} & \Delta y_k t_{kx} & \Delta x_k t_{ky} & \Delta y_k t_{ky} \end{bmatrix}.
    \label{eq:matrix_A_row}
\end{equation}
Once an admissible stencil $\Omega_u^I$ with non-singular $\mathbf{A}_I$ is selected (Section~\ref{subsec:stencil_selection}), grouping all six systems yields:
\begin{equation}
    \mathbf{A}_I \boldsymbol{\beta}^I = \mathbf{I}_{6 \times 6} \implies \boldsymbol{\beta}^I = \mathbf{A}_I^{-1},
    \label{eq:beta_matrix_inversion}
\end{equation}
where $\boldsymbol{\beta}^I = [\boldsymbol{\beta}_1^I, \dots, \boldsymbol{\beta}_6^I] \in \mathbb{R}^{6 \times 6}$. Gathering the local stencil DoFs into $\mathbf{e}_I \equiv \mathbf{L}_I \mathbf{e} = [e_1, \dots, e_6]^T \in \mathbb{R}^6$ via the Boolean restriction matrix $\mathbf{L}_I \in \mathbb{R}^{6 \times N}$ (where $\mathbf{e} = [e_1, \dots, e_N]^T \in \mathbb{R}^N$ is the global DoF vector), the local interpolated electric field is:
\begin{equation}
    \vec{E}_{I,h}(\mathbf{x}) = \sum_{a=1}^6 \vec{N}_a^I(\mathbf{x}) (\mathbf{e}_I)_a = \mathbf{P}(\mathbf{x} - \mathbf{x}_I) \boldsymbol{\beta}^I \mathbf{e}_I.
    \label{eq:field_interpolation_mlpg}
\end{equation}
Kinematically, the VNMM acts as a directional Vector Point Interpolation Method (Vector PIM) generalizing scalar PIM \cite{liu2001point, liu2005introduction}, preserving the Kronecker-delta property \eqref{eq:kronecker_projection} absent in MLS or RBF schemes \cite{atluri2002meshless}. Because $\vec{N}_a^I$ are linear in $(x, y)$, their curl and divergence are spatially constant across $\Omega_s^I$:
\begin{align}
    (\nabla \times \vec{N}_a^I)_z & = \frac{\partial N_{ay}^I}{\partial x} - \frac{\partial N_{ax}^I}{\partial y} = \beta_{5a}^I - \beta_{4a}^I \equiv c_{\mathrm{curl}, a}^I \in \mathbb{R}, \label{eq:curl_constant} \\
    \nabla \cdot \vec{N}_a^I      & = \frac{\partial N_{ax}^I}{\partial x} + \frac{\partial N_{ay}^I}{\partial y} = \beta_{3a}^I + \beta_{6a}^I \equiv c_{\mathrm{div}, a}^I \in \mathbb{R}. \label{eq:div_constant}
\end{align}
Defining the constant operator vectors $\mathbf{c}_{\mathrm{curl}}^I = (\boldsymbol{\beta}_{5, :}^I - \boldsymbol{\beta}_{4, :}^I)^T \in \mathbb{R}^6$ and $\mathbf{c}_{\mathrm{div}}^I = (\boldsymbol{\beta}_{3, :}^I + \boldsymbol{\beta}_{6, :}^I)^T \in \mathbb{R}^6$, the local curl and divergence evaluate directly as:
\begin{align}
    (\nabla \times \vec{E}_{I,h})_z(\mathbf{x}) & = \sum_{a=1}^6 c_{\mathrm{curl}, a}^I (\mathbf{e}_I)_a = (\mathbf{c}_{\mathrm{curl}}^I)^T \mathbf{e}_I \in \mathbb{R}, \label{eq:field_curl_constant} \\
    \nabla \cdot \vec{E}_{I,h}(\mathbf{x})      & = \sum_{a=1}^6 c_{\mathrm{div}, a}^I (\mathbf{e}_I)_a = (\mathbf{c}_{\mathrm{div}}^I)^T \mathbf{e}_I \in \mathbb{R}. \label{eq:field_div_constant}
\end{align}

\subsubsection{Formal Definition of the Global Discrete Approximation Space}
\label{subsubsec:discrete_space_definition}

On each active subdomain $D_I \equiv \Omega_s^I \cap \Omega$, let $\Pi_I: \mathbb{R}^N \to \mathcal{P}^1(D_I)$ denote the local polynomial reconstruction operator $\Pi_I(\mathbf{e})(\mathbf{x}) \equiv \vec{E}_{I,h}(\mathbf{x})|_{D_I}$ defined by \eqref{eq:field_interpolation_mlpg}.
Because adjacent subdomains overlap ($\Omega_s^I \cap \Omega_s^J \ne \emptyset$) and stencils $\Omega_u^I \ne \Omega_u^J$ are selected autonomously, $\vec{E}_{I,h}$ and $\vec{E}_{J,h}$ do not coincide pointwise on overlaps. Thus, the global approximation space $\mathcal{V}_h$ is a non-conforming subspace of the \emph{broken polynomial product space} $\prod_{I=1}^N \mathcal{P}^1(D_I)$:
\begin{equation}
    \mathcal{V}_h \equiv \left\{ \vec{v}_h = (\vec{v}_{1,h}, \dots, \vec{v}_{N,h}) \in \prod_{I=1}^N \mathcal{P}^1(D_I) \;:\; \exists\, \mathbf{e} \in \mathbb{R}^N \text{ such that } \vec{v}_{I,h} = \Pi_I(\mathbf{e}) \quad \forall I \in \{1, \dots, N\} \right\}.
    \label{eq:global_discrete_space_def}
\end{equation}

Let the union of active integration subdomains $D_I \equiv \Omega_s^I \cap \Omega$ define the \emph{covered domain region}:
\begin{equation}
    \Omega_{\mathrm{cov}} \equiv \bigcup_{I=1}^N D_I \subseteq \Omega.
    \label{eq:covered_region_def}
\end{equation}
On regular lattices with $\alpha_s \ge \sqrt{2}/2$, the subdomains form a complete cover ($\Omega_{\mathrm{cov}} = \Omega$), whereas on perturbed clouds with fixed $\alpha_s = 0.80$ they form a partial cover ($\Omega_{\mathrm{cov}} \subset \Omega$, Section~\ref{subsec:stage4_alpha_s}). Denoting by $M(\mathbf{x}) \equiv \sum_{I=1}^N \chi_{D_I}(\mathbf{x})$ the pointwise subdomain overlap multiplicity (with domain mean $\bar{M} = \frac{1}{|\Omega|}\iint_\Omega M(\mathbf{x})\,d\Omega \approx \pi \alpha_s^2$), a single-valued piecewise field for post-processing and visualization on $\Omega_{\mathrm{cov}}$ (where $M(\mathbf{x}) \ge 1$) is obtained via multiplicity-weighted averaging:
\begin{equation}
    \vec{E}_{\mathrm{global}}(\mathbf{x}) = \frac{1}{M(\mathbf{x})} \sum_{I=1, \mathbf{x} \in D_I}^N \vec{E}_{I,h}(\mathbf{x}), \qquad \forall \mathbf{x} \in \Omega_{\mathrm{cov}}.
    \label{eq:global_field_reconstruction}
\end{equation}

Even when mapped onto the single-valued piecewise field $\vec{u}_h \equiv \vec{E}_{\mathrm{global}}$ on $\Omega_{\mathrm{cov}}$ via Eq.~\eqref{eq:global_field_reconstruction}, $\vec{u}_h \notin H(\mathbf{curl}, \Omega)$ in general: across the circular boundary arcs $L_{sI} = \partial\Omega_s^I \cap \Omega$, both the overlap multiplicity $M(\mathbf{x})$ and the active local interpolants undergo jump discontinuities whenever $\mathbf{e} \ne \mathbf{0}$. Because tangential continuity ($[\hat{\mathbf{n}} \times \vec{u}_h] = \mathbf{0}$) across internal interfaces is necessary and sufficient for a piecewise-smooth vector field to belong to $H(\mathbf{curl}, \Omega)$ \cite{monk2003finite, boffi2013mixed}, any nonzero tangential jump along an arc $L_{sI}$ contributes a Dirac surface distribution $[\hat{\mathbf{n}} \times \vec{u}_h]_z\,\delta_{L_{sI}} \notin L^2(\Omega)$ to the distributional curl $\nabla \times \vec{u}_h$. On smooth solutions, however, these inter-stencil discrepancies decay at $\mathcal{O}(h^2)$ for the field and $\mathcal{O}(h)$ for its curl (Appendix~\ref{subsec:variational_consistency_strang}, Eqs.~\eqref{eq:stencil_jump_L2}--\eqref{eq:stencil_jump_curl}).

\subsection{Scale-Invariant Stencil Selection Algorithm}
\label{subsec:stencil_selection}

On unstructured or perturbed point clouds, the local moment matrix $\mathbf{A}_I \in \mathbb{R}^{6 \times 6}$ defined in \eqref{eq:matrix_A_row} can suffer from geometric or directional collinearity singularities ($\det(\mathbf{A}_I) = 0$). To ensure that every node $P_I$ is assigned an algebraically well-conditioned trial support stencil $\Omega_u^I$, we adopt the scale-invariant selection framework established in \cite{stable_vnmm2d_2026}.
Because columns 1--2 of $\mathbf{A}_I$ are dimensionless ($\mathcal{O}(1)$) while columns 3--6 scale with physical distances ($\mathcal{O}(h_I)$), the determinant exhibits a physical quartic scaling: $\det(\mathbf{A}_I) \sim \mathcal{O}(h_I^4)$. Normalizing local coordinates by $h_I$ yields the scale-invariant admissibility criterion:
\begin{equation}
    |\det(\widehat{\mathbf{A}}_I)| = \frac{|\det(\mathbf{A}_I)|}{h_I^4} \ge \widehat{\mathrm{Tol}} = 0.80,
    \label{eq:scale_invariant_det_criterion}
\end{equation}
which bounds the spectral 2-norm condition number of $\widehat{\mathbf{A}}_I$ to $\kappa_2(\widehat{\mathbf{A}}_I) \le 5.9$ and preserves compact nearest-neighbor supports ($K_{\mathrm{avg}} \approx 6.5\text{--}7.0$) without threshold decay under refinement \cite{stable_vnmm2d_2026}. Candidate stencils are assembled via a $k$-d tree query around $\mathbf{x}_I$, anchoring $\mathbf{x}_{(1)} = \mathbf{x}_I$ and selecting the nearest 5-combination satisfying \eqref{eq:scale_invariant_det_criterion}. In VNMM-MLPG, stencil selection is executed \textbf{once per node} ($N$ evaluations in total), avoiding the $\approx 10 N_{\mathrm{DoF}}$ Gauss-point queries required in cell-based VNMM-Gauss \cite{stable_vnmm2d_2026}.

\subsection{Local Weak Form on Overlapping Subdomains}
\label{subsec:local_weak_form}

Taking the inner product of the regularized Helmholtz equation \eqref{eq:helmholtz_pde} with a local vector test function $\vec{W}_I(\mathbf{x})$ over each active subdomain $D_I = \Omega_s^I \cap \Omega$ and applying Green's first vector identity yields the local weak identity:
\begin{equation}
    \begin{aligned}
        \iint_{D_I} \bigg[ \frac{1}{\mu_r} (\nabla \times \vec{W}_I)_z (\nabla \times \vec{E})_z & + \frac{s_{\mathrm{div}}}{\mu_r} (\nabla \cdot \vec{W}_I) (\nabla \cdot \vec{E}) - k_0^2 \epsilon_r \vec{W}_I \cdot \vec{E} \bigg] d\Omega \\
                                                                                                    & = \iint_{D_I} \vec{W}_I \cdot \vec{J}_s \, d\Omega + \int_{L_{sI}} \vec{W}_I \cdot \vec{\tau}(\vec{E}) \, ds \\
                                                                                                    & \quad + \int_{\Gamma_s^I} \vec{W}_I \cdot \vec{\tau}_{\mathrm{ext}} \, ds,
    \end{aligned}
    \label{eq:local_weak_form}
\end{equation}
where $\partial D_I = L_{sI} \cup \Gamma_s^I$ comprises the internal circular arc $L_{sI} = \partial\Omega_s^I \cap \Omega$ and the physical boundary segment $\Gamma_s^I = \Omega_s^I \cap \partial\Omega$ carrying prescribed natural fluxes $\vec{\tau}_{\mathrm{ext}}$ (Section~\ref{subsec:boundary_treatment}), and $\vec{\tau}(\vec{E}) = -\mu_r^{-1}(\vec{n} \times (\nabla \times \vec{E})) + s_{\mathrm{div}}\mu_r^{-1}(\nabla \cdot \vec{E})\vec{n}$ is the natural boundary flux along the outward unit normal $\vec{n}$. Grouping all volumetric and boundary contributions on the right-hand side into a local linear functional $\ell_{D_I}(\vec{W}_I)$ recasts \eqref{eq:local_weak_form} in compact volume energy form:
\begin{equation}
    a_{D_I}(\vec{W}_I, \vec{E}) - k_0^2 m_{D_I}(\vec{W}_I, \vec{E}) = \ell_{D_I}(\vec{W}_I),
    \label{eq:volume_energy_form}
\end{equation}
with symmetric local bilinear forms $a_{D_I}(\vec{W}_I, \vec{E}) \equiv \iint_{D_I} \mu_r^{-1}[(\nabla \times \vec{W}_I)_z (\nabla \times \vec{E})_z + s_{\mathrm{div}}(\nabla \cdot \vec{W}_I)(\nabla \cdot \vec{E})]\,d\Omega$ and $m_{D_I}(\vec{W}_I, \vec{E}) \equiv \iint_{D_I} \epsilon_r \vec{W}_I \cdot \vec{E}\,d\Omega$.
The specification of the local test space $\vec{W}_I$, the closed-form integration of $a_{D_I}$ and $m_{D_I}$, and the global overlapping assembly are presented in Sections~\ref{subsec:mlpg6_selection}--\ref{subsec:global_assembly_stability}.

\subsection{Selection of the MLPG6 (Local Galerkin) Test Space}
\label{subsec:mlpg6_selection}

A central question in local meshless methods concerns the choice of the vector test space $\vec{W}_I(\mathbf{x})$ in the local weak form \eqref{eq:local_weak_form}. In classical MLPG1 formulations based on Moving Least Squares (MLS) \cite{atluri1998new, atluri2002meshless} or radial-basis point interpolation (LRPIM) \cite{liu2005introduction}, the test function is constructed from a compactly supported weight function chosen to vanish identically along the internal subdomain boundary $L_{sI} = \partial\Omega_s^I \cap \Omega$ (i.e., $\vec{W}_I(\mathbf{x}) = \mathbf{0}$ for all $\mathbf{x} \in \partial\Omega_s^I$ on interior disks $D_I = \Omega_s^I$). This choice eliminates the contour flux integral $\int_{L_{sI}} \vec{W}_I \cdot \vec{\tau}(\vec{E}) \, ds$ in \eqref{eq:local_weak_form} and remains viable in MLS/RBF schemes because their non-polynomial trial derivatives vary spatially across $\Omega_s^I$. In the present Vector PIM setting with a fixed linear polynomial stencil ($\mathcal{P}^1$) per subdomain \cite{liu2001local, stable_vnmm2d_2026}, however, a boundary-vanishing test function causes the volumetric stiffness operator to vanish identically (analogous to the vanishing of constant-gradient domain integrals for compactly supported weights in scalar LPIM \cite{liu2001local}). Specifically, by Eqs.~\eqref{eq:field_curl_constant}--\eqref{eq:field_div_constant}, $(\nabla \times \vec{E}_{I,h})_z$ and $\nabla \cdot \vec{E}_{I,h}$ are spatially constant across $\Omega_s^I$. Factoring these constants out of the domain integrals in \eqref{eq:local_weak_form}--\eqref{eq:volume_energy_form} over an interior disk $D_I = \Omega_s^I$ (for homogeneous $\mu_r$) and applying Stokes' circulation theorem (Green's theorem in the plane) to the curl integral and Gauss' divergence theorem to the divergence integral yields:
\begin{equation}
    \begin{aligned}
        a_{\Omega_s^I}(\vec{W}_I, \vec{E}_{I,h}) & = \frac{(\mathbf{c}_{\mathrm{curl}}^I)^T \mathbf{e}_I}{\mu_r} \iint_{\Omega_s^I} (\nabla \times \vec{W}_I)_z \, d\Omega + \frac{s_{\mathrm{div}} (\mathbf{c}_{\mathrm{div}}^I)^T \mathbf{e}_I}{\mu_r} \iint_{\Omega_s^I} (\nabla \cdot \vec{W}_I) \, d\Omega \\
                                                 & = \frac{(\mathbf{c}_{\mathrm{curl}}^I)^T \mathbf{e}_I}{\mu_r} \underbrace{\oint_{\partial\Omega_s^I} \left(\vec{W}_I \cdot \vec{t}\right) ds}_{=\,0 \;\text{(by Stokes' theorem)}} + \frac{s_{\mathrm{div}} (\mathbf{c}_{\mathrm{div}}^I)^T \mathbf{e}_I}{\mu_r} \underbrace{\oint_{\partial\Omega_s^I} \left(\vec{W}_I \cdot \vec{n}\right) ds}_{=\,0 \;\text{(by Gauss' theorem)}} \equiv 0, \qquad \forall \mathbf{e}_I \in \mathbb{R}^6,
    \end{aligned}
    \label{eq:mlpg1_stiffness_collapse}
\end{equation}
where $\vec{t}$ and $\vec{n}$ denote the unit counterclockwise tangent and outward unit normal along $\partial\Omega_s^I$. Because $\vec{W}_I|_{\partial\Omega_s^I} = \mathbf{0}$, both closed contour integrals vanish identically for every trial basis function $\vec{N}_a^I$ ($a = 1, \dots, 6$). Consequently, all interior subdomain stiffness matrices vanish ($\mathbf{K}_{I,\mathrm{curl}}^{\mathrm{MLPG1}} \equiv \mathbf{0}$ and $\mathbf{K}_{I,\mathrm{div}}^{\mathrm{MLPG1}} \equiv \mathbf{0}$, so $\mathbf{K}_I^{\mathrm{MLPG1}} \equiv \mathbf{0}$). Weight-based MLPG1 testing is therefore mathematically incompatible with fixed-stencil first-order ($\mathcal{P}^1$) vector approximations.

To resolve this limitation and preserve the mathematical structure of the continuous Maxwell operators, we adopt the \textbf{MLPG6 (Local Bubnov--Galerkin)} framework \cite{atluri2002meshless}, choosing the vector test functions directly from the space of \textbf{trial shape functions themselves} within each local subdomain $\Omega_s^I$:
\begin{equation}
    \vec{W}_a^I(\mathbf{x}) = \vec{N}_a^I(\mathbf{x}) \quad \forall \mathbf{x} \in \Omega_s^I, \quad a = 1, \dots, 6.
    \label{eq:mlpg6_choice}
\end{equation}
This Local Galerkin choice confers three key advantages:
\begin{enumerate}
    \item \textbf{Exact Symmetry and Positive-Definiteness:} Testing with the trial basis ($\vec{W}_a^I = \vec{N}_a^I$) ensures that the local stiffness matrix $\mathbf{K}_I^{\mathrm{local}}$ and mass matrix $\mathbf{M}_I^{\mathrm{local}}$ are formed by inner products of the basis functions with themselves. This ensures that every local matrix is symmetric and positive-(semi)definite, naturally preserving the self-adjointness of the vector curl-curl operator and preventing unphysical complex eigenvalues.
    \item \textbf{Non-Trivial Stiffness Transmission:} Because $(\nabla \times \vec{N}_a^I)_z$ and $\nabla \cdot \vec{N}_a^I$ are non-zero constants, their inner products against themselves are positive, avoiding the null-stiffness behavior of MLPG1.
    \item \textbf{Exact Closed-Form Integration:} Integrands reduce to low-order Cartesian polynomials, enabling exact closed-form analytical integration over both circular and truncated subdomains without numerical quadrature errors (derived in Section~\ref{subsec:analytic_formulas} and Appendix~\ref{app:truncated_geometry}).
\end{enumerate}

\subsection{Closed-Form Analytical Integration of Local Subdomain Operators}
\label{subsec:analytic_formulas}

With $\vec{W}_a^I = \vec{N}_a^I$, we derive exact closed-form integrals for \eqref{eq:volume_energy_form} across the three subdomain configurations illustrated in Fig.~\ref{fig:schematic_mlpg}: first over interior disks (Case~1, Section~\ref{subsubsec:interior_disk_analytic}) and then over near-boundary truncated subdomains (Cases~2 and~3, Section~\ref{subsubsec:truncated_subdomains_analytic}).

\subsubsection{Interior Circular Disk Subdomains (Case 1)}
\label{subsubsec:interior_disk_analytic}

Let $\Omega_s^I$ be a circular disk of radius $r_s^I$. Its area and polar second moment of area are:
\begin{equation}
    \begin{aligned}
        S_0^{\mathrm{full}} & = \iint_{\Omega_s^I} d\Omega = \pi (r_s^I)^2,                                                                \\
        S_2^{\mathrm{full}} & = \iint_{\Omega_s^I} (x - x_I)^2 d\Omega = \iint_{\Omega_s^I} (y - y_I)^2 d\Omega = \frac{\pi (r_s^I)^4}{4}.
    \end{aligned}
    \label{eq:disk_moments}
\end{equation}
All odd moments vanish identically by radial symmetry:
\begin{equation}
    \iint_{\Omega_s^I} (x - x_I) \, d\Omega = \iint_{\Omega_s^I} (y - y_I) \, d\Omega = \iint_{\Omega_s^I} (x - x_I)(y - y_I) \, d\Omega \equiv 0.
\end{equation}

Within each subdomain $\Omega_s^I$, the complete local trial basis consists of the 6 vector shape functions $\vec{N}_a^I(\mathbf{x})$ ($a=1,\dots,6$). Substituting \eqref{eq:P1_shape_function_expansion_mlpg} into the local weak form \eqref{eq:local_weak_form}, the local stiffness matrix $\mathbf{K}_I^{\mathrm{local}} \in \mathbb{R}^{6 \times 6}$ is partitioned into curl and divergence contributions:
\begin{equation}
    \mathbf{K}_I^{\mathrm{local}} = \mathbf{K}_{I, \mathrm{curl}}^{\mathrm{local}} + s_{\mathrm{div}} \mathbf{K}_{I, \mathrm{div}}^{\mathrm{local}},
\end{equation}
where, using the constant vectors $\mathbf{c}_{\mathrm{curl}}^I, \mathbf{c}_{\mathrm{div}}^I \in \mathbb{R}^6$ defined in \eqref{eq:curl_constant}--\eqref{eq:div_constant}:
\begin{align}
    \mathbf{K}_{I, \mathrm{curl}}^{\mathrm{local}} & = \frac{1}{\mu_r} \left( \mathbf{c}_{\mathrm{curl}}^I (\mathbf{c}_{\mathrm{curl}}^I)^T \right) \pi (r_s^I)^2, \label{eq:k_curl_analytic} \\
    \mathbf{K}_{I, \mathrm{div}}^{\mathrm{local}}  & = \frac{1}{\mu_r} \left( \mathbf{c}_{\mathrm{div}}^I (\mathbf{c}_{\mathrm{div}}^I)^T \right) \pi (r_s^I)^2. \label{eq:k_div_analytic}
\end{align}
Both matrices are outer products of vectors with themselves, and are therefore symmetric and positive-semidefinite.

For the local mass matrix $M_{I, ab}^{\mathrm{local}} = \iint_{\Omega_s^I} \epsilon_r (\vec{N}_a^I \cdot \vec{N}_b^I) \, d\Omega$, expanding the dot product via \eqref{eq:P1_shape_function_expansion_mlpg} with $\boldsymbol{\beta}_{1:2, a}^I = [\beta_{1a}^I, \beta_{2a}^I]^T$ and $\boldsymbol{\beta}_{3:6, a}^I = [\beta_{3a}^I, \beta_{4a}^I, \beta_{5a}^I, \beta_{6a}^I]^T$ and integrating over the disk yields the closed-form expression:
\begin{equation}
    \mathbf{M}_I^{\mathrm{local}} = \epsilon_r \left( \mathbf{C}_0^I \pi (r_s^I)^2 + \mathbf{C}_1^I \frac{\pi (r_s^I)^4}{4} \right),
    \label{eq:m_local_analytic}
\end{equation}
where $\mathbf{C}_0^I = (\boldsymbol{\beta}_{1:2}^I)^T \boldsymbol{\beta}_{1:2}^I \in \mathbb{R}^{6 \times 6}$ and $\mathbf{C}_1^I = (\boldsymbol{\beta}_{3:6}^I)^T \boldsymbol{\beta}_{3:6}^I \in \mathbb{R}^{6 \times 6}$ are Gram matrices, ensuring that $\mathbf{M}_I^{\mathrm{local}}$ is symmetric and positive-definite.

\subsubsection{Truncated Subdomains Near Physical Boundaries (Cases 2 and 3)}
\label{subsubsec:truncated_subdomains_analytic}

When a local subdomain $\Omega_s^I$ intersects $\partial\Omega$ (Cases~2 and 3 in Fig.~\ref{fig:schematic_mlpg}), the active integration domain is $D_I \equiv \Omega_s^I \cap \Omega$. Because the curl and divergence of $\mathcal{P}^1$ shape functions are spatially constant (Eqs.~\eqref{eq:field_curl_constant}--\eqref{eq:field_div_constant}), the stiffness integrand factors outside the integral:
\begin{equation}
    \mathbf{K}_{I, \mathrm{curl}}^{\mathrm{trunc}} = \frac{S_0^I}{\mu_r} \mathbf{c}_{\mathrm{curl}}^I (\mathbf{c}_{\mathrm{curl}}^I)^T, \qquad
    \mathbf{K}_{I, \mathrm{div}}^{\mathrm{trunc}} = \frac{S_0^I}{\mu_r} \mathbf{c}_{\mathrm{div}}^I (\mathbf{c}_{\mathrm{div}}^I)^T,
    \label{eq:k_curl_truncated}
\end{equation}
where $S_0^I = \iint_{D_I} d\Omega$ denotes the active subdomain area. Defining the dimensionless active area fraction $\eta_I \equiv S_0^I / S_0^{\mathrm{full}} \in (0, 1]$ and denoting the untrimmed full-disk operators from \eqref{eq:k_curl_analytic}--\eqref{eq:m_local_analytic} by $\mathbf{K}_I^{\mathrm{full}}$ and $\mathbf{M}_I^{\mathrm{full}}$, \textbf{the exact truncated stiffness matrix is obtained by simple area scaling}:
\begin{equation}
    \mathbf{K}_I^{\mathrm{trunc}} = \eta_I \mathbf{K}_I^{\mathrm{full}}.
    \label{eq:k_truncated_scaled}
\end{equation}
Unlike classical MLPG strategies that shrink near-boundary radii tangentially ($r_s^I = \min(r_{s,0}^I, \mathrm{dist}(\mathbf{x}_I, \partial\Omega))$ \cite{atluri2002meshless}) or map truncated sectors onto reference circles for Gauss quadrature \cite{atluri2004meshless}, Eq.~\eqref{eq:k_truncated_scaled} requires only the scalar area fraction $\eta_I$. Because each disk $\Omega_s^I$ is centered at a physical node $\mathbf{x}_I \in \overline{\Omega}$, on any convex polygonal domain with interior corner angles $\theta_{\mathrm{corner}}$ the area fraction satisfies $\eta_I \ge \theta_{\mathrm{corner}}/(2\pi)$ ($\eta_I \ge 1/2$ on flat walls and $\eta_I \ge 1/4$ at $90^\circ$ corners). This lower bound prevents the small-cut ill-conditioning ($\eta \to 0$) typical of unfitted cell methods such as CutFEM \cite{hansbo2002unfitted}.

For the local mass matrix $\mathbf{M}_I^{\mathrm{local}}$ on a truncated subdomain, broken reflection symmetry introduces first- and second-order geometric moments relative to $\mathbf{x}_I$:
\begin{equation}
    \mathbf{M}_I^{\mathrm{trunc}} = \epsilon_r \left( \mathbf{P}_0^I S_0^I + \mathbf{P}_x^I S_{1x}^I + \mathbf{P}_y^I S_{1y}^I + \mathbf{P}_{xx}^I S_{2xx}^I + \mathbf{P}_{yy}^I S_{2yy}^I + \mathbf{P}_{xy}^I S_{2xy}^I \right),
    \label{eq:m_truncated_moments}
\end{equation}
where $S_{1x}^I = \iint_{D_I} (x - x_I)\,d\Omega$, $S_{1y}^I = \iint_{D_I} (y - y_I)\,d\Omega$, $S_{2xx}^I = \iint_{D_I} (x - x_I)^2\,d\Omega$, $S_{2yy}^I = \iint_{D_I} (y - y_I)^2\,d\Omega$, and $S_{2xy}^I = \iint_{D_I} (x - x_I)(y - y_I)\,d\Omega$ are the Cartesian first- and second-order geometric moments. Denoting by $\mathbf{b}_r^I \equiv (\boldsymbol{\beta}_{r,:}^I)^T \in \mathbb{R}^6$ the $r$-th row of $\boldsymbol{\beta}^I$ from \eqref{eq:beta_matrix_inversion}, expanding $\vec{N}_a^I \cdot \vec{N}_b^I$ via \eqref{eq:P_basis_matrix}--\eqref{eq:P1_shape_function_expansion_mlpg} yields the constant symmetric outer-product matrices $\mathbf{P}_0^I, \dots, \mathbf{P}_{xy}^I \in \mathbb{R}^{6 \times 6}$ explicitly:
\begin{equation}
    \begin{aligned}
        \mathbf{P}_0^I    & = \mathbf{b}_1^I (\mathbf{b}_1^I)^T + \mathbf{b}_2^I (\mathbf{b}_2^I)^T, &
        \mathbf{P}_{xx}^I & = \mathbf{b}_3^I (\mathbf{b}_3^I)^T + \mathbf{b}_5^I (\mathbf{b}_5^I)^T, \\
        \mathbf{P}_{yy}^I & = \mathbf{b}_4^I (\mathbf{b}_4^I)^T + \mathbf{b}_6^I (\mathbf{b}_6^I)^T, &
        \mathbf{P}_x^I    & = \mathbf{b}_1^I (\mathbf{b}_3^I)^T + \mathbf{b}_3^I (\mathbf{b}_1^I)^T + \mathbf{b}_2^I (\mathbf{b}_5^I)^T + \mathbf{b}_5^I (\mathbf{b}_2^I)^T, \\
        \mathbf{P}_y^I    & = \mathbf{b}_1^I (\mathbf{b}_4^I)^T + \mathbf{b}_4^I (\mathbf{b}_1^I)^T + \mathbf{b}_2^I (\mathbf{b}_6^I)^T + \mathbf{b}_6^I (\mathbf{b}_2^I)^T, &
        \mathbf{P}_{xy}^I & = \mathbf{b}_3^I (\mathbf{b}_4^I)^T + \mathbf{b}_4^I (\mathbf{b}_3^I)^T + \mathbf{b}_5^I (\mathbf{b}_6^I)^T + \mathbf{b}_6^I (\mathbf{b}_5^I)^T.
    \end{aligned}
    \label{eq:P_outer_products}
\end{equation}
Closed-form analytical expressions for $\eta_I$ and the Cartesian geometric moments across rectilinear walls and orthogonal corners are derived in Appendix~\ref{app:truncated_geometry}.

\subsection{Global Matrix Assembly, Overlapping Energy Formulation, and Algebraic Properties}
\label{subsec:global_assembly_stability}

Once the local MLPG6 subdomain operators have been integrated in closed form (Sections~\ref{subsec:mlpg6_selection}--\ref{subsec:analytic_formulas}), the algebraic structure of the global discrete system depends on how local weighted-residual equations are assembled across overlapping nodal neighborhoods. Here we first review classical row-by-row MLPG/LPIM assembly and explain why it loses symmetry (Section~\ref{subsubsec:row_mlpg_symmetry}), and then present the proposed Symmetric Subdomain Galerkin (SSG) assembly together with its discrete action functional, algebraic definiteness properties, and multiplicity-weighted Rayleigh quotient (Section~\ref{subsubsec:symmetric_assembly}).

\subsubsection{Classical Row-by-Row MLPG Assembly and Matrix Symmetry Analysis}
\label{subsubsec:row_mlpg_symmetry}

In classical row-by-row MLPG/LPIM \cite{atluri1998new, atluri2002meshless, liu2001local}, each node $P_I$ is tested with a single function $\vec{W}_I(\mathbf{x}) = \vec{N}_{i_0}^I(\mathbf{x})$ (where $\mathbf{x}_{i_0} = \mathbf{x}_I$ in $\Omega_u^I$). Note that if the contour integral $-\int_{L_{sI}} \vec{W}_I \cdot \vec{\tau}(\vec{E}_{I,h})\,ds$ in \eqref{eq:local_weak_form} is evaluated using the same fixed linear stencil $\Omega_u^I$, Green's identity causes the volume and boundary stiffness terms to cancel identically (requiring numerical quadrature across shifting neighbor stencils along $L_{sI}$ in classical LPIM \cite{liu2001local}). Testing the volume form \eqref{eq:volume_energy_form} with $\vec{W}_I = \vec{N}_{i_0}^I$ yields the $I$-th row equation $\sum_{k=1}^6 K_{I, k}^{\mathrm{row}} (\mathbf{e}_I)_k = F_I$ in closed form:
\begin{equation}
    K_{I, k}^{\mathrm{row}} = \left[ \frac{1}{\mu_r} c_{\mathrm{curl}, i_0}^I c_{\mathrm{curl}, k}^I + \frac{s_{\mathrm{div}}}{\mu_r} c_{\mathrm{div}, i_0}^I c_{\mathrm{div}, k}^I \right] \pi (r_s^I)^2.
    \label{eq:k_row_analytic}
\end{equation}

\begin{remark}[Inherent Matrix Unsymmetry in Classical MLPG and LPIM]
    The global matrices $\mathbf{K}$ and $\mathbf{M}$ assembled row-by-row via \eqref{eq:k_row_analytic} are intrinsically unsymmetric ($K_{Ij} \ne K_{Ji}$) \cite{atluri2002meshless, liu2001local} due to distinct integration domains ($\Omega_s^I \ne \Omega_s^J$) and autonomous stencils ($\boldsymbol{\beta}^I \ne \boldsymbol{\beta}^J$). Consequently, row-by-row assembly loses positive-definiteness on non-polynomial fields under refinement (Section~\ref{subsubsec:mms_assembly_comparison}).
\end{remark}

\subsubsection{Symmetric Subdomain Galerkin Assembly and Overlapping Rayleigh Quotient}
\label{subsubsec:symmetric_assembly}

To preserve self-adjointness without the curved arc integrals over $L_{sI}$ of row-by-row LPIM or the pairwise lens-intersection quadratures of symmetric overlapping methods such as the Method of Finite Spheres \cite{de2000method, nicomedes2017meshfree, nicomedes2021acoustic}, we introduce the \textbf{Symmetric Subdomain Galerkin (SSG)} scheme. Using the local symmetric bilinear forms \eqref{eq:volume_energy_form}, we define the discrete quadratic action functional $J_h: \mathbb{R}^N \to \mathbb{R}$ summed over all active subdomains $\{D_I\}_{I=1}^N$:
\begin{equation}
    J_h(\mathbf{e}) = \sum_{I=1}^N \left\{ \frac{1}{2} a_{D_I}(\vec{E}_{I,h}, \vec{E}_{I,h}) - \frac{1}{2} k_0^2 m_{D_I}(\vec{E}_{I,h}, \vec{E}_{I,h}) - \ell_{D_I}(\vec{E}_{I,h}) \right\} = \frac{1}{2} \mathbf{e}^T (\mathbf{K} - k_0^2 \mathbf{M}) \mathbf{e} - \mathbf{e}^T \mathbf{F}.
    \label{eq:discrete_action_functional}
\end{equation}
Equivalently, each disk $\Omega_s^I$ is tested against \textit{all six} local MLPG6 trial shape functions $\vec{W}_a^I = \vec{N}_a^I$ ($a=1,\dots,6$), and the resulting $6 \times 6$ closed-form local matrices \eqref{eq:k_curl_analytic}--\eqref{eq:m_local_analytic} (and \eqref{eq:k_truncated_scaled}--\eqref{eq:m_truncated_moments}) are accumulated via Boolean scatter operations:
\begin{equation}
    \mathbf{K} = \sum_{I=1}^N \mathbf{L}_I^T \mathbf{K}_I^{\mathrm{local}} \mathbf{L}_I, \qquad \mathbf{M} = \sum_{I=1}^N \mathbf{L}_I^T \mathbf{M}_I^{\mathrm{local}} \mathbf{L}_I, \qquad \mathbf{F} = \sum_{I=1}^N \mathbf{L}_I^T \mathbf{f}_I^{\mathrm{local}},
    \label{eq:global_assembly}
\end{equation}
where $\mathbf{f}_I^{\mathrm{local}} = [\ell_{D_I}(\vec{N}_1^I), \dots, \ell_{D_I}(\vec{N}_6^I)]^T \in \mathbb{R}^6$. Directly from the congruence structure \eqref{eq:global_assembly}, two algebraic and spectral properties hold:
\begin{enumerate}[(i)]
    \item \textbf{Discrete Symmetry and Real Spectrum ($\mathbf{K} = \mathbf{K}^T$, $\mathbf{M} = \mathbf{M}^T$):}
          Because each local block $\mathbf{K}_I^{\mathrm{local}} \ge 0$ and $\mathbf{M}_I^{\mathrm{local}} > 0$ is a symmetric Gram/outer-product matrix, the congruence sum \eqref{eq:global_assembly} is symmetric. Eliminating Dirichlet (PEC) boundary DoFs (Section~\ref{subsubsec:dirichlet_bc}) preserves symmetry on the reduced generalized eigenproblem, ensuring real eigenvalues ($\operatorname{Im}(\lambda) = 0$). While symmetry precludes complex-conjugate eigenvalues, real irrotational gradient modes ($\nabla H^1(\Omega)$) are shifted outside the operational band by the divergence penalty $s_{\mathrm{div}}\mathbf{K}_{\mathrm{div}}$ ($\lambda_{\mathrm{spur}} \sim \mathcal{O}(s_{\mathrm{div}})$).
    \item \textbf{Global Positive Semidefiniteness of $\mathbf{K}$ and Strict Positive Definiteness of $\mathbf{M}$:}
          For any $\mathbf{e} \in \mathbb{R}^N$, $\mathbf{e}^T \mathbf{K} \mathbf{e} = \sum_{I=1}^N (\mathbf{L}_I \mathbf{e})^T \mathbf{K}_I^{\mathrm{local}} (\mathbf{L}_I \mathbf{e}) \ge 0$. Moreover, because every node $I$ belongs to its own stencil $\Omega_u^I$, any non-zero global vector $\mathbf{e} \ne \mathbf{0}$ yields at least one non-zero local vector $\mathbf{L}_I \mathbf{e} \ne \mathbf{0}$. Since $\mathbf{M}_I^{\mathrm{local}} > 0$ on every subdomain of positive area, $\mathbf{e}^T \mathbf{M} \mathbf{e} = \sum_{I=1}^N (\mathbf{L}_I \mathbf{e})^T \mathbf{M}_I^{\mathrm{local}} (\mathbf{L}_I \mathbf{e}) > 0$, proving $\mathbf{M} > 0$ algebraically without relying on global geometric coverage.
\end{enumerate}

For conservative cavity eigenproblems ($\ell_{D_I} \equiv 0$), stationarity of the discrete action functional ($\delta J_h(\mathbf{e}) = \mathbf{0}$) yields the symmetric generalized eigenvalue problem:
\begin{equation}
    \mathbf{K} \mathbf{e} = \lambda_h \mathbf{M} \mathbf{e}, \qquad \lambda_h = \frac{\mathbf{e}^T \mathbf{K} \mathbf{e}}{\mathbf{e}^T \mathbf{M} \mathbf{e}} = \frac{a_h(\vec{E}_h, \vec{E}_h)}{m_h(\vec{E}_h, \vec{E}_h)},
    \label{eq:rayleigh_quotient_overlap}
\end{equation}
with broken bilinear forms $a_h(\vec{v}_h, \vec{w}_h) = \sum_{I=1}^N a_{D_I}(\vec{v}_{I,h}, \vec{w}_{I,h})$ and $m_h(\vec{v}_h, \vec{w}_h) = \sum_{I=1}^N m_{D_I}(\vec{v}_{I,h}, \vec{w}_{I,h})$. When evaluated on a globally continuous solenoidal field $\vec{E}$ ($\vec{E}_{I,h} = \vec{E}|_{D_I}$, $\nabla \cdot \vec{E} = 0$), this overlapping sum yields the continuous multiplicity-weighted Rayleigh quotient:
\begin{equation}
    \lambda^{(M)}(\vec{E}) = \frac{\iint_{\Omega} M(\mathbf{x}) \mu_r^{-1} |\nabla \times \vec{E}|^2 \, d\Omega}{\iint_{\Omega} M(\mathbf{x}) \epsilon_r |\vec{E}|^2 \, d\Omega}, \qquad M(\mathbf{x}) = \bar{M} + \widetilde{M}(\mathbf{x}).
    \label{eq:multiplicity_rayleigh_quotient}
\end{equation}
While the mean $\bar{M}$ cancels, the fluctuation $\widetilde{M}(\mathbf{x})$ averages out to $\mathcal{O}(h_{\mathrm{nom}}^2)$ on regular rectangular lattices via cell-wise moment cancellation (Appendix~\ref{subsec:conformity_discussion}; quantified in Section~\ref{subsubsec:multiplicity_rayleigh_effect}), whereas its interplay with $\mathcal{P}^1$ interpolation, non-conformity, and stabilization in \eqref{eq:rayleigh_quotient_overlap} is analyzed in Appendix~\ref{subsec:variational_consistency_strang}. Because each unstabilized local stiffness block $\mathbf{K}_I^{\mathrm{local}}$ has rank at most 2, achieving strict positive definiteness of $\mathbf{K}$ after PEC boundary condensation requires both Dirichlet elimination (Section~\ref{subsubsec:dirichlet_bc}) and orthogonal energy stabilization (Section~\ref{subsec:rank_deficiency_stabilization}).

\subsection{Analysis and Treatment of Boundary Conditions}
\label{subsec:boundary_treatment}

For any node $\mathbf{x}_I$ whose subdomain $\Omega_s^I$ intersects $\partial\Omega = \Gamma_D \cup \Gamma_N$, the boundary segment $\Gamma_s^I = \Omega_s^I \cap \partial\Omega$ enters the right-hand side of the local weak form \eqref{eq:local_weak_form} via $I_{\Gamma_s^I} = \int_{\Gamma_s^I} \vec{W}_I \cdot \vec{\tau}(\vec{E}) \, ds$. From \eqref{eq:helmholtz_pde}, the complete natural traction vector $\vec{\tau}(\vec{E})$ decomposes into tangential curl and normal divergence components:
\begin{equation}
    \vec{\tau}(\vec{E}) \equiv \vec{\tau}_t(\vec{E}) + \tau_n(\vec{E})\vec{n} = -\frac{1}{\mu_r} \left( \vec{n} \times (\nabla \times \vec{E}) \right) + \frac{s_{\mathrm{div}}}{\mu_r} (\nabla \cdot \vec{E}) \vec{n}.
    \label{eq:traction_vector}
\end{equation}

\subsubsection{Inhomogeneous and Homogeneous Dirichlet Boundary Conditions}
\label{subsubsec:dirichlet_bc}

On $\Gamma_D$, where $\vec{E} \cdot \vec{t} = g_D(\mathbf{x})$ is prescribed, aligning boundary directors parallel to the wall ($\vec{t}_k \parallel \Gamma_D$) and invoking the directional Kronecker-delta property \eqref{eq:kronecker_projection} enforces $e_k = \vec{E}(\mathbf{x}_k) \cdot \vec{t}_k = g_D(\mathbf{x}_k)$ directly by nodal collocation. Let $\mathbf{A} \equiv \mathbf{K} - k_0^2 \mathbf{M}$ denote the assembled global system matrix. Partitioning the global DoFs into active interior ($\mathbf{e}_{\mathrm{int}}$) and prescribed boundary ($\mathbf{e}_{\mathrm{bnd}}$) unknowns yields the condensed symmetric system $\mathbf{A}\mathbf{e} = \mathbf{F}$ in reduced form matching \cite{stable_vnmm2d_2026}:
\begin{equation}
    \mathbf{A}_{\mathrm{int}, \mathrm{int}} \mathbf{e}_{\mathrm{int}} = \mathbf{F}_{\mathrm{int}} - \mathbf{A}_{\mathrm{int}, \mathrm{bnd}} \mathbf{e}_{\mathrm{bnd}},
    \label{eq:dirichlet_condensed_system}
\end{equation}
On homogeneous PEC walls ($\mathbf{e}_{\mathrm{bnd}} = \mathbf{0}$), Eq.~\eqref{eq:dirichlet_condensed_system} simplifies to $\mathbf{A}_{\mathrm{int}, \mathrm{int}} \mathbf{e}_{\mathrm{int}} = \mathbf{F}_{\mathrm{int}}$. In the source-free, lossless cavity limit ($\mathbf{F}_{\mathrm{int}} = \mathbf{0}$), it reduces further to the generalized eigenvalue problem $\mathbf{K}_{\mathrm{red}} \mathbf{e}_{\mathrm{red}} = \lambda \mathbf{M}_{\mathrm{red}} \mathbf{e}_{\mathrm{red}}$ without boundary integrals along $\Gamma_D$, where $\mathbf{K}_{\mathrm{red}} \equiv \mathbf{K}_{\mathrm{int},\mathrm{int}}$, $\mathbf{M}_{\mathrm{red}} \equiv \mathbf{M}_{\mathrm{int},\mathrm{int}}$, $\mathbf{e}_{\mathrm{red}} \equiv \mathbf{e}_{\mathrm{int}}$, and $\lambda = k_0^2$.

\subsubsection{Natural Boundary Conditions (Neumann / PMC) and Solenoidality Assumption}
\label{subsubsec:neumann_divergence_treatment}

On natural boundaries $\Gamma_N$, physical electromagnetic conditions prescribe the tangential magnetic traction $\vec{\tau}_t(\vec{E}) = -\mu_r^{-1}(\vec{n} \times (\nabla \times \vec{E})) = \vec{g}_N$ \eqref{eq:neumann_bvp_pde}, whereas the regularized traction \eqref{eq:traction_vector} also contains the normal divergence flux $\tau_n(\vec{E}) = s_{\mathrm{div}}\mu_r^{-1}(\nabla \cdot \vec{E})$.
In the present formulation, natural boundary conditions are formulated under the explicit assumption that the exact solution is \textbf{solenoidal} ($\nabla \cdot \vec{E} \equiv 0$ in $\overline{\Omega}$), which holds in source-free homogeneous media and for the manufactured and cavity problems investigated herein. Under this solenoidal condition, $\tau_n(\vec{E}) \equiv 0$ on $\Gamma_N$, so substituting $\vec{\tau}_t(\vec{E}) = \vec{g}_N$ into the boundary integral $I_{\Gamma_s^I}$ yields the consistent local Neumann load contribution:
\begin{equation}
    F_a^{\mathrm{Neumann}, I} = \int_{\Omega_s^I \cap \Gamma_N} \vec{g}_N(\mathbf{x}) \cdot \vec{N}_a^I(\mathbf{x}) \, ds,
    \label{eq:neumann_load_integral}
\end{equation}
which vanishes identically on homogeneous PMC walls ($\vec{g}_N = \mathbf{0}$). For general problems where the exact field is non-solenoidal ($\nabla \cdot \vec{E} \ne 0$, e.g., due to free space charges, non-solenoidal current sources, or material inhomogeneities), the normal regularization flux $\int_{\Gamma_s^I} s_{\mathrm{div}}\mu_r^{-1}(\nabla \cdot \vec{E})(\vec{N}_a^I \cdot \vec{n})\,ds$ must be explicitly included in the boundary formulation.
Because $\vec{N}_a^I$ are linear along straight segments $\Gamma_s^I$, 2-point Gauss--Legendre quadrature evaluates $F_a^{\mathrm{Neumann}, I}$ for affine boundary data $\vec{g}_N$ to machine precision (Table~\ref{tab:stage1_element_verification}), and provides fourth-order accuracy for general smooth profiles $\vec{g}_N$.

\subsection{Discrete Operator Rank Deficiency and Orthogonal Energy Stabilization}
\label{subsec:rank_deficiency_stabilization}

Although the SSG scheme preserves global symmetry and yields full-rank local mass matrices, evaluating the curl--divergence stiffness integral over a single subdomain $D_I$ with a $\mathcal{P}^1$ vector stencil produces an enlarged local kernel that can soften discrete modes under refinement. To diagnose and suppress this mechanism without sacrificing cell-free analytical integration, we first establish the rank deficiency of the unstabilized local stiffness matrix (Section~\ref{subsubsec:local_rank_analysis}) and then introduce an orthogonal energy stabilization term derived from the canonical decomposition of the 2D gradient tensor $\nabla\vec{E}$ (Section~\ref{subsubsec:canonical_decomposition}).

\subsubsection{Local Operator Rank Analysis under Linear Vector Basis}
\label{subsubsec:local_rank_analysis}

On each stencil $\Omega_u^I$, the linear field $\vec{E}_{I,h}(\mathbf{x})$ has 6 degrees of freedom $\mathbf{e}_I = \mathbf{L}_I \mathbf{e} \in \mathbb{R}^6$ and polynomial coefficients $\mathbf{a} = \boldsymbol{\beta}^I \mathbf{e}_I \in \mathbb{R}^6$. By \eqref{eq:field_curl_constant}--\eqref{eq:field_div_constant} and \eqref{eq:k_curl_analytic}--\eqref{eq:k_div_analytic}, the local curl and divergence stiffness matrices on $D_I$ (active area $S_0^I = \eta_I \pi (r_s^I)^2$) are rank-1 dyadic outer products of $\mathbf{c}_{\mathrm{curl}}^I$ and $\mathbf{c}_{\mathrm{div}}^I$. Consequently, the unstabilized local stiffness matrix has rank at most 2:
\begin{equation}
    \mathrm{rank}(\mathbf{K}_I^{\mathrm{local}}) = \mathrm{rank}\left(\mathbf{K}_{I, \mathrm{curl}}^{\mathrm{local}} + s_{\mathrm{div}}\mathbf{K}_{I, \mathrm{div}}^{\mathrm{local}}\right) \le 1 + 1 = 2,
    \label{eq:stiffness_rank_deficiency}
\end{equation}
whereas $\mathbf{M}_I^{\mathrm{local}}$ \eqref{eq:m_local_analytic} has full rank 6. In cell-based VNMM-Gauss \cite{stable_vnmm2d_2026}, the $3 \times 3$ Gauss points inside each cell sample shifting neighbor stencils across an enlarged cell nodal patch, eliminating the unpenalized local harmonic gradient modes within each cell. By contrast, single-disk integration leaves a 4-dimensional local stiffness kernel consisting of two constant vector translations and two curl- and divergence-free harmonic linear gradient modes. This unpenalized nullspace allows discrete modes to soften, leading to sublinear eigenvalue convergence ($p = 0.53\text{--}0.84$ on $15^2\text{--}49^2$, Section~\ref{subsubsec:cavity_jitter_error}).

\subsubsection{Canonical Orthogonal Decomposition and VMS-Inspired Stabilization}
\label{subsubsec:canonical_decomposition}

To restore full rank over the 4-dimensional linear gradient subspace, we decompose the 2D field gradient tensor $\nabla\vec{E} \in \mathbb{R}^{2 \times 2}$ into four mutually orthogonal components:
\begin{equation}
    \nabla \vec{E} = \underbrace{\frac{1}{2} (\nabla \times \vec{E})_z \mathbf{J}}_{\text{curl}} + \underbrace{\frac{1}{2} (\nabla \cdot \vec{E}) \mathbf{I}_2}_{\text{divergence}} + \underbrace{\nabla_{\mathrm{sym},0}\vec{E}}_{\text{harmonic gradient}},
    \label{eq:gradient_orthogonal_decomposition}
\end{equation}
where $\mathbf{J} = \begin{bmatrix} 0 & -1 \\ 1 & 0 \end{bmatrix}$ and $\nabla_{\mathrm{sym},0}\vec{E} = \frac{1}{2}(\nabla\vec{E} + (\nabla\vec{E})^T) - \frac{1}{2}(\nabla\cdot\vec{E})\mathbf{I}_2$ is the trace-free symmetric gradient tensor (corresponding to fields satisfying $\nabla\times\vec{E}=0$ and $\nabla\cdot\vec{E}=0$). Alongside $\mathbf{c}_{\mathrm{curl}}^I$ and $\mathbf{c}_{\mathrm{div}}^I$ from Section~\ref{subsec:trial_space}, the two trace-free harmonic gradient modes on stencil $\Omega_u^I$ map to discrete nodal vectors:
\begin{align}
    \mathbf{v}_{\mathrm{shear}}^I & = (\boldsymbol{\beta}_{5, :}^I + \boldsymbol{\beta}_{4, :}^I)^T \quad \left(\text{cross-derivative harmonic mode } \frac{\partial E_y}{\partial x} + \frac{\partial E_x}{\partial y}\right), \label{eq:v_shear_def}                   \\
    \mathbf{v}_{\mathrm{hyp}}^I   & = (\boldsymbol{\beta}_{3, :}^I - \boldsymbol{\beta}_{6, :}^I)^T \quad \left(\text{hyperbolic quadrupole harmonic mode } \frac{\partial E_x}{\partial x} - \frac{\partial E_y}{\partial y}\right). \label{eq:v_hyp_def}
\end{align}
Because the two curl- and divergence-free harmonic modes $\{\mathbf{v}_{\mathrm{shear}}^I, \mathbf{v}_{\mathrm{hyp}}^I\}$ reside in $\ker(\mathbf{K}_I^{\mathrm{local}})$ and the matrix trace inner product $\operatorname{tr}(\mathbf{U}^T\mathbf{V}) = \sum_{i,j \in \{x,y\}} U_{ij} V_{ij}$ applied to $\vec{W}_I = [W_x, W_y]^T$ and $\vec{E} = [E_x, E_y]^T$ satisfies
\begin{equation}
    \operatorname{tr}\!\left((\nabla_{\mathrm{sym},0}\vec{W}_I)^T \nabla_{\mathrm{sym},0}\vec{E}\right) = \frac{1}{2}\left[\left(\frac{\partial W_y}{\partial x} + \frac{\partial W_x}{\partial y}\right)\left(\frac{\partial E_y}{\partial x} + \frac{\partial E_x}{\partial y}\right) + \left(\frac{\partial W_x}{\partial x} - \frac{\partial W_y}{\partial y}\right)\left(\frac{\partial E_x}{\partial x} - \frac{\partial E_y}{\partial y}\right)\right],
    \label{eq:dev_trace_identity}
\end{equation}
we introduce the \textbf{Orthogonal Energy Stabilization Operator} to penalize these modes. It is defined by the local bilinear form $a_{D_I}^{\mathrm{stab}}(\vec{W}_I, \vec{E}) = \iint_{D_I} (2\alpha_{\mathrm{stab}}/\mu_r) \operatorname{tr}\!\left((\nabla_{\mathrm{sym},0}\vec{W}_I)^T \nabla_{\mathrm{sym},0}\vec{E}\right) d\Omega = \mathbf{w}_I^T \mathbf{K}_{I, \mathrm{stab}}^{\mathrm{local}} \mathbf{e}_I$ for local test and trial DoF vectors $\mathbf{w}_I, \mathbf{e}_I \in \mathbb{R}^6$:
\begin{equation}
    \mathbf{K}_{I, \mathrm{stab}}^{\mathrm{local}} = \frac{\alpha_{\mathrm{stab}}}{\mu_r} S_0^I \left[ \mathbf{v}_{\mathrm{shear}}^I (\mathbf{v}_{\mathrm{shear}}^I)^T + \mathbf{v}_{\mathrm{hyp}}^I (\mathbf{v}_{\mathrm{hyp}}^I)^T \right], \qquad
    \mathbf{K}_I^{\mathrm{local, stab}} = \mathbf{K}_{I, \mathrm{curl}}^{\mathrm{local}} + s_{\mathrm{div}}\mathbf{K}_{I, \mathrm{div}}^{\mathrm{local}} + \mathbf{K}_{I, \mathrm{stab}}^{\mathrm{local}}.
    \label{eq:k_stab_analytic}
\end{equation}
By construction, $\mathbf{K}_{I, \mathrm{stab}}^{\mathrm{local}}$ vanishes identically on constant vector fields, pure-curl fields, and pure-divergence fields, while penalizing the two harmonic linear gradient modes to raise $\mathrm{rank}(\mathbf{K}_I^{\mathrm{local, stab}})$ from 2 to 4 (analogous to orthogonal zero-energy mode control \cite{flanagan1981uniform}). Note that because $\mathbf{K}_{I, \mathrm{stab}}^{\mathrm{local}}$ penalizes the trace-free symmetric gradient, it is not neutral on harmonic affine fields unless balanced in the right-hand functional (see Remark~\ref{rem:patch_test_stab} in Section~\ref{subsec:stage2_patch_test}), and it introduces a finite-mesh perturbation in homogeneous eigenvalue problems. Motivated heuristically by Variational Multiscale (VMS) subgrid-scale concepts \cite{hughes1998variational, codina2000stabilization} (rather than a formal multiscale variational splitting), we adopt the dimensionless power-law scaling:
\begin{equation}
    \alpha_{\mathrm{stab}}(h_{\mathrm{char}}) = c_{\mathrm{stab}} \left(\frac{h_{\mathrm{char}}}{L}\right)^{\gamma},
    \label{eq:alpha_stab_vms_scaling}
\end{equation}
where $c_{\mathrm{stab}} \ge 0$ is a dimensionless calibration constant, $\gamma \ge 0$ is the scaling exponent, and $L$ and $h_{\mathrm{char}}$ are the reference length and characteristic spacing defined in Section~\ref{subsec:nodal_discretization}. As calibrated across the extended $15\times 15\text{--}81\times 81$ refinement study in Section~\ref{subsubsec:resolution_stagnation_stabilization}, linear scaling ($\gamma = 1$, $c_{\mathrm{stab}} = 4.0$) preserves $\mathcal{O}(h_{\mathrm{char}})$ bias decay while preventing the fine-mesh rank re-softening exhibited by quadratic decay ($\gamma = 2$).

\begin{remark}[Variational Analysis and Convergence Hypotheses]
\label{rem:variational_summary}
    Appendix~\ref{subsec:variational_consistency_strang} formalizes the variational properties of the non-conforming SSG operators over $\mathcal{V}_h \not\subset H(\mathbf{curl}, \Omega)$, including multiplicity homogenization, local Green consistency, fixed-mesh definiteness ($\mathbf{K}_{\mathrm{red}} > 0$), and the spectral error decomposition.
\end{remark}

\section{Numerical Validation and Results}
\label{sec:results}

The six-stage numerical validation on $\Omega = [0, \pi]^2\text{ m}^2$ (matching the benchmark suite of \cite{stable_vnmm2d_2026}) progresses from element-level analytical checks (Stage~1), affine patch tests (Stage~2), and auxiliary local-equilibrium MMS stability checks (Stage~3) to subdomain coverage sensitivity (Stage~4), PEC cavity spectral validation (Stage~5), and performance comparisons against VNMM-Gauss (Stage~6).

\subsection{Stage 1: Element-Level Closed-Form Verification}
\label{subsec:stage1_results}

To establish the mathematical fidelity and cell-free precision of the local operators before undertaking global system assembly, the element-level stiffness matrix $\mathbf{K}_I^{\mathrm{local}}$, mass matrix $\mathbf{M}_I^{\mathrm{local}}$, and boundary terms were verified by comparing closed-form analytical evaluations against high-order numerical quadratures. We investigate three distinct geometric regimes:
\begin{enumerate}
    \item Full interior circular disks ($\eta = 1.0$),
    \item Truncated subdomains along or near physical boundaries ($\eta \in (0, 1)$),
    \item 1D natural boundary integrals entering Neumann boundary conditions along $\Gamma_s^I = \Omega_s^I \cap \partial\Omega$.
\end{enumerate}

\begin{table}[!t]
    \centering
    \small
    \caption{Element-level verification of stiffness $\mathbf{K}$, mass $\mathbf{M}$, and boundary terms across full and truncated integration domains, comparing analytical closed-form formulas against high-precision numerical quadrature.}
    \label{tab:stage1_element_verification}

    \par\medskip
    \textbf{(a) 2D Subdomain Integrals: Stiffness $\mathbf{K}$ and Mass $\mathbf{M}$ on Full and Truncated Domains}\par\smallskip

    \resizebox{\textwidth}{!}{
        \begin{tabular}{llcccc}
            \toprule
            \textbf{Domain}    & \textbf{Geometric}        & \textbf{Fraction}     & \textbf{Stiffness Error}         & \textbf{Mass Error}                               & \textbf{Scaling Error}                                            \\
            \textbf{Type}      & \textbf{Configuration}    & $\boldsymbol{\eta_I}$ & $\boldsymbol{\Delta \mathbf{K}}$ & $\boldsymbol{\Delta \mathbf{M}}$ \textbf{(Exact)} & $\boldsymbol{\eta \mathbf{M}^{\mathrm{full}}}$ \textbf{vs.~Exact} \\
            \midrule
            Full Circular Disk & Interior node             & $1.0$                 & $2.14 \times 10^{-14}$           & $1.87 \times 10^{-14}$                            & $0.00$                                                            \\
            Flat Boundary Wall & Half-disk ($\eta=1/2$)    & $1/2$                 & $6.91 \times 10^{-17}$           & $2.93 \times 10^{-15}$                            & $1.56 \times 10^{0}$                                              \\
            Corner Vertex      & Quarter-disk ($\eta=1/4$) & $1/4$                 & $1.85 \times 10^{-16}$           & $3.70 \times 10^{-16}$                            & $4.22 \times 10^{-1}$                                             \\
            Near-Wall Cut      & Chord ($d = 0.5 r_s$)     & $0.8045$              & $2.09 \times 10^{-7}$            & $1.69 \times 10^{-7}$                             & $1.55 \times 10^{-1}$                                             \\
            \bottomrule
        \end{tabular}
    }

    \par\medskip
    \textbf{(b) 1D Natural Boundary Line Integrals along $\Gamma_s^I = \Omega_s^I \cap \partial\Omega$}\par\smallskip

    \resizebox{\textwidth}{!}{
        \begin{tabular}{llccc}
            \toprule
            \textbf{Boundary}                                 & \textbf{Intersection}         & \textbf{1-Point Gauss} & \textbf{2-Point Gauss} & \textbf{3-Point Gauss} \\
            \textbf{Term}                                     & \textbf{Segment $\Gamma_s^I$} & \textbf{Error}         & \textbf{Error}         & \textbf{Error}         \\
            \midrule
            Inhomogeneous Neumann $F_a^{\mathrm{Neumann}, I}$ & Straight wall $[-r_s, r_s]$   & $1.01 \times 10^{-1}$  & $1.92 \times 10^{-16}$ & $4.81 \times 10^{-17}$ \\
            \bottomrule
        \end{tabular}
    }
\end{table}

\subsubsection{Full Circular Disks (Interior Subdomains)}
For an uncut circular disk $\Omega_s^I$ ($r_s = 0.35\text{ m}$, $s_{\mathrm{div}} = 4.0$), polar Gauss--Legendre quadrature with $n_r = 6$ radial and $n_\theta = 16$ circumferential points integrates the constant stiffness integrand and quadratic mass integrand exactly. Measuring the relative difference in the entrywise maximum norm ($\|\mathbf{X}\|_{\max} \equiv \max_{a,b} |X_{ab}|$), Table~\ref{tab:stage1_element_verification}(a) shows that the closed-form formulas \eqref{eq:k_curl_analytic}--\eqref{eq:m_local_analytic} match numerical quadrature to double-precision machine epsilon ($\le 2.14 \times 10^{-14}$).

\subsubsection{Truncated Subdomains Near Boundaries}
For the boundary-truncated subdomains of Section~\ref{subsubsec:truncated_subdomains_analytic} and Appendix~\ref{app:truncated_geometry}, Table~\ref{tab:stage1_element_verification}(a) confirms exact agreement for area-scaled stiffness \eqref{eq:k_truncated_scaled} on flat walls and corners ($\le 1.85 \times 10^{-16}$). For the mass matrix $\mathbf{M}_I^{\mathrm{trunc}}$, naive area scaling $\eta_I \mathbf{M}_I^{\mathrm{full}}$ incurs $15.5\%\text{--}156\%$ error due to broken reflection symmetry, whereas the full moment expansion \eqref{eq:m_truncated_moments} restores machine-precision accuracy ($\le 2.93 \times 10^{-15}$). For interior chord cuts ($d = 0.5 r_s$), the $\sim 10^{-7}$ discrepancy reflects the $\mathcal{O}(h_{\mathrm{quad}})$ boundary-clipping error of polar quadrature across a non-aligned straight chord $r(\theta) = d/\cos\theta$, which the analytical cap and moment formulas avoid completely.

\subsubsection{Boundary Condition Integrals (Neumann)}
Along straight natural-boundary segments $\Gamma_s^I = \Omega_s^I \cap \partial\Omega$ of length $2r_s$ with affine boundary flux $\vec{g}_N(s)$, the Neumann integrand \eqref{eq:neumann_load_integral} is quadratic in $s$. As shown in Table~\ref{tab:stage1_element_verification}(b), a 1-point Gauss rule exhibits a $10.1\%$ truncation error, whereas 2-point Gauss--Legendre quadrature matches the closed-form line integral to machine precision ($\le 1.92 \times 10^{-16}$) while offering fourth-order accuracy for smooth $\vec{g}_N$.

\subsection{Stage 2: Galerkin Patch Test for Linear Consistency}
\label{subsec:stage2_patch_test}

The Galerkin patch test was evaluated across the seven canonical affine vector states:
\begin{enumerate}
    \item Constant Field $X$: $\vec{E}^* = [1, 0]^T$ ($\nabla \times \vec{E}^* = 0, \nabla \cdot \vec{E}^* = 0$).
    \item Constant Field $Y$: $\vec{E}^* = [0, 1]^T$ ($\nabla \times \vec{E}^* = 0, \nabla \cdot \vec{E}^* = 0$).
    \item Hyperbolic Quadrupole (Harmonic): $\vec{E}^* = [x, -y]^T$ ($\nabla \times \vec{E}^* = 0, \nabla \cdot \vec{E}^* = 0$).
    \item Cross-Derivative Quadrupole (Harmonic): $\vec{E}^* = [y, x]^T$ ($\nabla \times \vec{E}^* = 0, \nabla \cdot \vec{E}^* = 0$).
    \item Pure Solenoidal Curl: $\vec{E}^* = [-y, x]^T$ ($(\nabla \times \vec{E}^*)_z = 2, \nabla \cdot \vec{E}^* = 0$).
    \item Pure Irrotational Divergence: $\vec{E}^* = [x, y]^T$ ($(\nabla \times \vec{E}^*)_z = 0, \nabla \cdot \vec{E}^* = 2$).
    \item General Affine Field: $\vec{E}^* = [1.5 + 0.8x - 0.3y, -0.7 + 0.4x + 1.2y]^T$ ($(\nabla \times \vec{E}^*)_z = 0.7, \nabla \cdot \vec{E}^* = 2.0$).
\end{enumerate}

In each mode, the boundary-value problem is formulated on $\Omega$ in the static / zero-frequency limit ($k_0 = 0$) of \eqref{eq:helmholtz_pde} with exact essential Dirichlet conditions prescribed along the entire physical boundary: $\vec{E} \cdot \vec{t} = \vec{E}^* \cdot \vec{t}$ on $\partial\Omega$. These states span the complete 6-dimensional linear basis $\mathcal{P}^1$, representing canonical electromagnetic configurations: uniform fields (Modes~1--2), electrostatic quadrupoles (Modes~3--4), Faraday induction (Mode~5), and uniform space-charge distributions (Mode~6). Because the curl and divergence of any affine field are spatial constants and $k_0 = 0$, the continuous volumetric source in \eqref{eq:helmholtz_pde} vanishes identically ($\vec{J}_s = \nabla \times (\mu_r^{-1}\nabla\times\vec{E}^*) - s_{\mathrm{div}}\mu_r^{-1}\nabla(\nabla\cdot\vec{E}^*) \equiv \mathbf{0}$), while the local weak boundary fluxes along $\partial D_I$ produce the discrete local load vectors $\mathbf{f}_I^{\mathrm{local}}$ assembled into $\mathbf{F}$.

We compared the classical unsymmetric row-by-row MLPG assembly \eqref{eq:k_row_analytic} against the Symmetric Subdomain Galerkin assembly \eqref{eq:global_assembly} on an $11 \times 11$ discretization ($N = 121$ nodes, 81 free interior DoFs). Two nodal configurations were tested: (a) a uniform Cartesian grid with alternating interior directors ($45^\circ / 135^\circ$), and (b) an irregular cloud with $25\%$ random nodal position jitter and isotropic random directors. Exact Dirichlet values $e_k = \vec{E}^*(\mathbf{x}_k) \cdot \vec{t}_k$ were imposed on all boundary nodes, and the partitioned linear system $\mathbf{K}_{\mathrm{int}, \mathrm{int}} \mathbf{e}_{\mathrm{int}} = \mathbf{F}_{\mathrm{int}} - \mathbf{K}_{\mathrm{int}, \mathrm{bnd}} \mathbf{e}_{\mathrm{bnd}}$ was solved for the interior DoFs via sparse factorization.

\begin{table}[!t]
    \centering
    \small
    \caption{Relative infinity-norm error on free interior degrees of freedom ($\|\mathbf{e}_{\mathrm{int}} - \mathbf{e}^*\|_{\infty} / \|\mathbf{e}^*\|_{\infty}$) for the Galerkin Patch Test across 7 affine states, comparing unsymmetric row-by-row and Symmetric Subdomain Galerkin assemblies.}
    \label{tab:patch_test_results}
    \resizebox{\textwidth}{!}{
        \begin{tabular}{lcccccc}
            \toprule
                                       &                                          &                                     & \multicolumn{2}{c}{\textbf{Regular Grid ($11 \times 11$)}} & \multicolumn{2}{c}{\textbf{Irregular Cloud (25\% Jitter)}}                                                   \\
            \cmidrule(lr){4-5} \cmidrule(lr){6-7}
            \textbf{Affine Field Mode} & $\boldsymbol{(\nabla\times\vec{E}^*)_z}$ & $\boldsymbol{\nabla\cdot\vec{E}^*}$ & \textbf{Unsymmetric}                                       & \textbf{Symmetric}                                         & \textbf{Unsymmetric}   & \textbf{Symmetric}     \\
            \midrule
            1. Constant Field $X$      & 0.0                                      & 0.0                                 & $1.41 \times 10^{-14}$                                     & $2.45 \times 10^{-14}$                                     & $1.04 \times 10^{-14}$ & $9.67 \times 10^{-14}$ \\
            2. Constant Field $Y$      & 0.0                                      & 0.0                                 & $1.77 \times 10^{-14}$                                     & $1.60 \times 10^{-14}$                                     & $3.80 \times 10^{-14}$ & $7.48 \times 10^{-14}$ \\
            3. Hyperbolic Quadrupole   & 0.0                                      & 0.0                                 & $4.47 \times 10^{-15}$                                     & $8.47 \times 10^{-15}$                                     & $9.62 \times 10^{-15}$ & $7.20 \times 10^{-14}$ \\
            4. Cross-Derivative Quad.  & 0.0                                      & 0.0                                 & $1.66 \times 10^{-14}$                                     & $1.23 \times 10^{-14}$                                     & $2.40 \times 10^{-14}$ & $6.03 \times 10^{-14}$ \\
            5. Pure Solenoidal Curl    & 2.0                                      & 0.0                                 & $8.70 \times 10^{-15}$                                     & $1.26 \times 10^{-14}$                                     & $1.71 \times 10^{-14}$ & $1.91 \times 10^{-14}$ \\
            6. Pure Divergence         & 0.0                                      & 2.0                                 & $1.11 \times 10^{-14}$                                     & $1.21 \times 10^{-14}$                                     & $2.40 \times 10^{-14}$ & $1.29 \times 10^{-13}$ \\
            7. General Affine Field    & 0.7                                      & 2.0                                 & $2.03 \times 10^{-14}$                                     & $1.26 \times 10^{-14}$                                     & $2.31 \times 10^{-14}$ & $1.21 \times 10^{-13}$ \\
            \bottomrule
        \end{tabular}
    }
\end{table}

As reported in Table~\ref{tab:patch_test_results}, the relative error across all seven affine modes remains on the order of $10^{-14}\text{--}10^{-15}$ on regular grids and $\le 1.29 \times 10^{-13}$ on perturbed point clouds for both formulations. Because the local boundary fluxes on the circular arcs $L_{sI}$ are evaluated using the analytical field, this test verifies the local Green identity, $\mathcal{P}^1$ linear completeness, and operator assembly rather than global consistency for general source-driven BVPs (Section~\ref{subsec:deficiencies_challenges} and Appendix~\ref{subsubsec:consistency_residual}).

\begin{remark}[Patch Test Consistency under Orthogonal Energy Stabilization]
\label{rem:patch_test_stab}
    Activating the local orthogonal energy stabilization $\mathbf{K}_{I,\mathrm{stab}}^{\mathrm{local}}$ \eqref{eq:k_stab_analytic} with parameter $\alpha_{\mathrm{stab}} \in [0.1, 5.0]$ preserves consistency in two distinct ways. First, without any modification to the right-hand load vector, $\mathbf{K}_{I, \mathrm{stab}}^{\mathrm{local}}$ annihilates all uniform constant states (Modes~1 and 2, $\vec{E}^* = [a_1, a_2]^T$), pure-curl states ($\vec{E}^* = [-y, x]^T$), and pure-divergence states ($\vec{E}^* = [x, y]^T$) identically ($\mathbf{K}_{I, \mathrm{stab}}^{\mathrm{local}} \mathbf{L}_I \mathbf{e}^* \equiv \mathbf{0}$). Consequently, their relative errors remain below $7.0 \times 10^{-14}$ across all meshes. Second, for affine modes possessing a non-zero trace-free symmetric gradient (Modes~3, 4, and 7), because $\nabla_{\mathrm{sym},0}\vec{E}^*$ is spatially constant ($\nabla \cdot \nabla_{\mathrm{sym},0}\vec{E}^* \equiv \mathbf{0}$), integration by parts gives $a_{D_I}^{\mathrm{stab}}(\vec{W}_I, \vec{E}^*) = \oint_{\partial D_I} \vec{W}_I \cdot \vec{\tau}_{\mathrm{stab}}(\vec{E}^*) \, ds$, with explicit stabilization boundary flux $\vec{\tau}_{\mathrm{stab}}(\vec{E}^*) = (2\alpha_{\mathrm{stab}}/\mu_r)\nabla_{\mathrm{sym},0}\vec{E}^*\,\vec{n}$. Including this consistent boundary flux in $\ell_{D_I}$ reproduces all seven patch test states to machine precision ($\le 1.3 \times 10^{-13}$). When applied without auxiliary stabilization boundary fluxes under the adopted linear scaling law \eqref{eq:alpha_stab_vms_scaling} ($\gamma = 1$, $c_{\mathrm{stab}} = 4.0$), the stabilization perturbation decays linearly as $\mathcal{O}(h_{\mathrm{char}})$ under refinement.
\end{remark}

\subsection{Stage 3: Local Approximation and Algebraic Verification under Exact Local Equilibrium}
\label{subsec:stage3_mms}

As an auxiliary verification stage prior to the primary spectral validation in Stage~5, Stage~3 employs the Method of Manufactured Solutions (MMS) under exact local weak equilibrium to test non-polynomial local approximation, truncated-subdomain boundary fluxes, and algebraic stability independently of an eigensolver (rather than validating the eigenproblem or source-driven BVPs). It compares: (i) Pure Dirichlet vs.\ Mixed Dirichlet--Neumann conditions under SSG assembly (Section~\ref{subsubsec:mms_dirichlet_neumann}), and (ii) SSG vs.\ unsymmetric row-by-row assembly (Section~\ref{subsubsec:mms_assembly_comparison}).

Both studies solve the MMS problem on $\Omega = [0, \pi]^2\text{ m}^2$ ($L = \pi\text{ m}$) with exact analytical field:
\begin{equation}
    \vec{E}^*(x, y) = \begin{bmatrix} \sin(x) \cos(y) \\ -\cos(x) \sin(y) \end{bmatrix}, \qquad (\nabla \times \vec{E}^*)_z = 2\sin(x)\sin(y), \quad \nabla \cdot \vec{E}^* = 0.
\end{equation}
To obtain a strictly coercive resolvent operator free of internal resonances, we replace $-k_0^2 \mapsto +\kappa$ ($\kappa = +1.0\text{ m}^{-2}$) in \eqref{eq:helmholtz_pde}, using $s_{\mathrm{div}} = 6.0$, $\alpha_s = 0.80$, $\mu_r = \epsilon_r = 1$, $\alpha_{\mathrm{stab}} = 0$ (since $+\kappa\mathbf{M}$ already guarantees strict definiteness), and $\vec{J}_s = (2 + \kappa) \vec{E}^*$. In accordance with \eqref{eq:local_weak_form} (Appendix~\ref{subsubsec:consistency_residual}, Eq.~\eqref{eq:local_green_flux_balance}), the load vector evaluates $f_{I, a}^{\mathrm{local}} = \ell_{D_I}(\vec{N}_a^I) = a_{D_I}(\vec{N}_a^I, \vec{E}^*) + \kappa m_{D_I}(\vec{N}_a^I, \vec{E}^*)$, including the analytical fluxes on $L_{sI}$. Refinement spans the six grids matching \cite{stable_vnmm2d_2026} ($15\times 15$ to $49\times 49$, $h_{\mathrm{nom}} \in [0.2244, 0.0654]\text{ m}$) with isotropic random directors ($\text{seed}=42$, $|\det(\widehat{\mathbf{A}}_I)| \ge 0.80$).

\paragraph*{Explicit Definition of Error Metrics in Multiplicity-Weighted Broken Sobolev Norms:}
In accordance with Appendix~\ref{subsec:conformity_discussion}, the numerical errors reported throughout this section are evaluated in the \textbf{multiplicity-weighted relative broken Sobolev norms} over the active subdomains $\{D_I\}_{I=1}^N$:
\begin{align}
    \text{Relative Broken Field Error: } \quad & \varepsilon_{\boldsymbol{L}^2, h} = \left( \frac{\sum_{I=1}^N \iint_{D_I} |\vec{E}_{I,h} - \vec{E}^*|^2 \, d\Omega}{\sum_{I=1}^N \iint_{D_I} |\vec{E}^*|^2 \, d\Omega} \right)^{1/2}, \label{eq:mms_broken_L2} \\
    \text{Relative Broken Rotational Error: } \quad & \varepsilon_{\mathbf{curl}, h} = \left( \frac{\sum_{I=1}^N \iint_{D_I} |(\nabla \times \vec{E}_{I,h})_z - (\nabla \times \vec{E}^*)_z|^2 \, d\Omega}{\sum_{I=1}^N \iint_{D_I} |(\nabla \times \vec{E}^*)_z|^2 \, d\Omega} \right)^{1/2}, \label{eq:mms_broken_curl}
\end{align}
where the volume integrals on each disk $\Omega_s^I$ are evaluated numerically via polar Gauss quadrature ($n_r=4, n_\theta=12$) to sample the trigonometric modes of $\vec{E}^*$ accurately, and summing over subdomains weights each point $\mathbf{x} \in \Omega_{\mathrm{cov}}$ by its overlap multiplicity $M(\mathbf{x})$.

\subsubsection{Refinement Convergence for Pure Dirichlet and Mixed Neumann Conditions}
\label{subsubsec:mms_dirichlet_neumann}

To systematically test the boundary formulation developed in Section~\ref{subsec:boundary_treatment}, two distinct boundary-value problem configurations were investigated:
\begin{enumerate}
    \item \textbf{Case A (Pure Inhomogeneous Dirichlet):} On all four boundary walls $\partial\Omega$, the tangential trace is prescribed by the exact manufactured field $e_k = \vec{E}^*(\mathbf{x}_k) \cdot \vec{t}_k$, and boundary nodes are eliminated via condensed system \eqref{eq:dirichlet_condensed_system}.
    \item \textbf{Case B (Mixed Dirichlet--Neumann / PMC):} Inhomogeneous Dirichlet conditions are imposed on the vertical walls $x \in \{0, \pi\}$, while the horizontal walls $y \in \{0, \pi\}$ are treated as homogeneous natural Neumann (PMC) boundaries. Because $(\nabla \times \vec{E}^*)_z = 2\sin(x)\sin(y)$ vanishes identically along $y=0$ and $y=\pi$ and $\nabla \cdot \vec{E}^* \equiv 0$, the exact electromagnetic boundary flux satisfies $\vec{\tau}(\vec{E}^*) \equiv \mathbf{0}$. Therefore, the boundary nodes along $y \in \{0, \pi\}$ are retained as active, unconstrained degrees of freedom and determined directly by the global system.
\end{enumerate}

\begin{table}[!t]
    \centering
    \small
    \caption{MMS refinement convergence comparison between Pure Dirichlet and Mixed Dirichlet--Neumann (PMC) boundary conditions for VNMM-MLPG 2D ($s_{\mathrm{div}} = 6.0, \alpha_s = 0.80$, isotropic random directors) across the 6-mesh sequence from \cite{stable_vnmm2d_2026}. Errors are evaluated in the multiplicity-weighted relative broken Sobolev norms ($\varepsilon_{\boldsymbol{L}^2, h}$ and $\varepsilon_{\mathbf{curl}, h}$, Eqs.~\eqref{eq:mms_broken_L2}--\eqref{eq:mms_broken_curl}) over the non-conforming broken space $\mathcal{V}_h$.}
    \label{tab:mms_convergence}
    \resizebox{\textwidth}{!}{
        \begin{tabular}{cc|ccc|ccc}
            \toprule
                                                                 &                                              & \multicolumn{3}{c|}{\textbf{Case A: Pure Dirichlet}} & \multicolumn{3}{c}{\textbf{Case B: Mixed Dirichlet--Neumann (PMC)}}                                                                                                                                                                           \\
            \textbf{Mesh}                                        & $\boldsymbol{h_{\mathrm{nom}}}$ \textbf{[m]} & \textbf{Active DoFs}                                 & $\boldsymbol{\varepsilon_{\boldsymbol{L}^2, h}}$                     & $\boldsymbol{\varepsilon_{\mathbf{curl}, h}}$        & \textbf{Active DoFs}                                & $\boldsymbol{\varepsilon_{\boldsymbol{L}^2, h}}$    & $\boldsymbol{\varepsilon_{\mathbf{curl}, h}}$        \\
            \midrule
            $15 \times 15$                                       & 0.2244                                       & 169                                                  & $6.529\%$                                                           & $21.817\%$                                          & 197                                                 & $7.330\%$                                           & $20.865\%$                                          \\
            $22 \times 22$                                       & 0.1496                                       & 400                                                  & $3.398\%$                                                           & $13.991\%$                                          & 442                                                 & $3.943\%$                                           & $13.661\%$                                          \\
            $29 \times 29$                                       & 0.1122                                       & 729                                                  & $1.922\%$                                                           & $10.556\%$                                          & 785                                                 & $2.187\%$                                           & $10.282\%$                                          \\
            $36 \times 36$                                       & 0.0898                                       & 1156                                                 & $1.184\%$                                                           & $8.336\%$                                           & 1226                                                & $1.352\%$                                           & $8.250\%$                                           \\
            $43 \times 43$                                       & 0.0748                                       & 1681                                                 & $0.804\%$                                                           & $6.835\%$                                           & 1765                                                & $1.022\%$                                           & $6.741\%$                                           \\
            $49 \times 49$                                       & 0.0654                                       & 2209                                                 & $0.637\%$                                                           & $5.941\%$                                           & 2305                                                & $0.780\%$                                           & $5.892\%$                                           \\
            \midrule
            \multicolumn{2}{c|}{\textbf{Fitted Asymptotic Rate}} & ---                                          & $\boldsymbol{\mathcal{O}(h_{\mathrm{nom}}^{1.93})}$  & $\boldsymbol{\mathcal{O}(h_{\mathrm{nom}}^{1.05})}$                 & ---                                                 & $\boldsymbol{\mathcal{O}(h_{\mathrm{nom}}^{1.85})}$ & $\boldsymbol{\mathcal{O}(h_{\mathrm{nom}}^{1.02})}$ \\
            \multicolumn{2}{c|}{\textbf{Expected Approx.\ Rate}} & ---                                          & $\boldsymbol{\mathcal{O}(h_{\mathrm{nom}}^{2.00})}$  & $\boldsymbol{\mathcal{O}(h_{\mathrm{nom}}^{1.00})}$                 & ---                                                 & $\boldsymbol{\mathcal{O}(h_{\mathrm{nom}}^{2.00})}$ & $\boldsymbol{\mathcal{O}(h_{\mathrm{nom}}^{1.00})}$ \\
            \bottomrule
        \end{tabular}
    }
\end{table}

\begin{figure}[!t]
    \centering
    \includegraphics[width=0.78\textwidth]{fig_mlpg_mms_convergence.pdf}
    \caption{Log--log refinement curves corresponding to Table~\ref{tab:mms_convergence}, comparing Pure Dirichlet (solid) and Mixed Dirichlet--Neumann (dashed) broken errors against the theoretical $\mathcal{O}(h_{\mathrm{nom}}^2)$ and $\mathcal{O}(h_{\mathrm{nom}})$ reference slopes.}
    \label{fig:mms_convergence}
\end{figure}

The numerical results in Table~\ref{tab:mms_convergence} and Fig.~\ref{fig:mms_convergence} confirm four findings:
\begin{enumerate}
    \item \textbf{Field Convergence in Broken $L^2$ Norm:} In both Case~A (Dirichlet) and Case~B (Mixed Neumann), $\varepsilon_{\boldsymbol{L}^2, h}$ decreases monotonically across all six refinement levels at empirical rates $\mathcal{O}(h_{\mathrm{nom}}^{1.93})$ and $\mathcal{O}(h_{\mathrm{nom}}^{1.85})$, closely matching the theoretical $\mathcal{O}(h_{\mathrm{nom}}^2)$ order.
    \item \textbf{Rotational Convergence in Broken $H(\mathbf{curl})$ Semi-Norm:} Both formulations achieve monotonic rates of $\mathcal{O}(h_{\mathrm{nom}}^{1.05})$ and $\mathcal{O}(h_{\mathrm{nom}}^{1.02})$ in $\varepsilon_{\mathbf{curl}, h}$, matching the expected linear order $\mathcal{O}(h_{\mathrm{nom}})$.
    \item \textbf{Solvability Without Penalty Parameters:} The mixed Dirichlet--Neumann problem was solved without adding boundary penalty terms or boundary element integrations, matching Section~\ref{subsec:boundary_treatment}.
    \item \textbf{Boundary Accuracy and Absence of Boundary Layer Degradation:} Releasing the Neumann nodes on $y \in \{0, \pi\}$ produces slightly lower broken rotational errors than fully clamped Dirichlet walls ($20.865\%$ vs.\ $21.817\%$ on $15 \times 15$ and $5.892\%$ vs.\ $5.941\%$ on $49 \times 49$). Partitioning subdomains into interior ($\Omega_s^I \subset \Omega$) and boundary-truncated ($\Omega_s^B$, $\eta \le 1/2$) disks on the $49 \times 49$ grid yields a lower $L^2$ error at the boundary ($\mathbf{0.53\%}$, rate $\mathcal{O}(h_{\mathrm{nom}}^{2.03})$) than in the interior ($\mathbf{0.64\%}$, rate $\mathcal{O}(h_{\mathrm{nom}}^{1.93})$), while the 1D trace error along $\partial\Omega$ decays at $\mathcal{O}(h_{\mathrm{nom}}^{2.04})$ ($6.79\% \to 0.57\%$). Thus, directional collocation and natural weak boundary fluxes enforce boundary conditions cleanly without boundary-layer oscillations.
\end{enumerate}

\subsubsection{Comparison of Symmetric Subdomain Galerkin and Unsymmetric Row Assemblies}
\label{subsubsec:mms_assembly_comparison}

To inspect matrix symmetry and the smallest eigenvalue ($\min\mathrm{Re}(\lambda)$) of the reduced shifted operator $\mathbf{A}_{\mathrm{int},\mathrm{int}} = \mathbf{K}_{\mathrm{int},\mathrm{int}} + \kappa\mathbf{M}_{\mathrm{int},\mathrm{int}}$ ($\kappa = +1.0 > 0$), the Pure Dirichlet MMS benchmark was solved under both SSG assembly \eqref{eq:global_assembly} ($\mathbf{A}_{\mathrm{sym}} = \mathbf{A}_{\mathrm{sym}}^T > 0$) and classical unsymmetric row-by-row MLPG assembly \eqref{eq:k_row_analytic} ($\mathbf{A}_{\mathrm{unsym}} \ne \mathbf{A}_{\mathrm{unsym}}^T$). Table~\ref{tab:mms_assembly_comparison} reports their comparative performance across coarse-to-medium grids ($9 \times 9$ to $17 \times 17$, stopping at $17 \times 17$ because row-by-row assembly already diverges sharply).

\begin{table}[!t]
    \centering
    \small
    \caption{MMS broken-norm errors ($\varepsilon_{\boldsymbol{L}^2, h}$, $\varepsilon_{\mathbf{curl}, h}$, Eqs.~\eqref{eq:mms_broken_L2}--\eqref{eq:mms_broken_curl}), conditioning, and minimum real eigenvalue comparing Symmetric Subdomain Galerkin and Unsymmetric Row-by-Row Petrov--Galerkin assemblies ($s_{\mathrm{div}} = 6.0, \alpha_s = 0.80$, Pure Dirichlet conditions).}
    \label{tab:mms_assembly_comparison}
    \resizebox{\textwidth}{!}{
        \begin{tabular}{cccccc}
            \toprule
            \textbf{Mesh}  & \textbf{Formulation}           & $\boldsymbol{\varepsilon_{\boldsymbol{L}^2, h}}$ \textbf{(Broken $L^2$)} & $\boldsymbol{\varepsilon_{\mathbf{curl}, h}}$ \textbf{(Broken $|H(\mathbf{curl})|$)} & $\mathrm{cond}(\mathbf{A}_{\mathrm{int}, \mathrm{int}})$ & $\min \mathrm{Re}(\lambda)$ \\
            \midrule
            $9 \times 9$   & Symmetric (Subdomain Galerkin) & $\mathbf{12.31\%}$                         & $\mathbf{47.86\%}$                             & $1.85 \times 10^2$                                       & $+1.64$                     \\
            $9 \times 9$   & Unsymmetric (Row Petrov--Gal.)  & $75.09\%$                                  & $104.31\%$                                     & $1.82 \times 10^2$                                       & $-14.19$ (10 neg.)          \\
            \midrule
            $13 \times 13$ & Symmetric (Subdomain Galerkin) & $\mathbf{7.97\%}$                          & $\mathbf{29.13\%}$                             & $8.25 \times 10^2$                                       & $+0.36$                     \\
            $13 \times 13$ & Unsymmetric (Row Petrov--Gal.)  & $895.71\%$                                 & $2063.16\%$                                    & $2.80 \times 10^4$                                       & $-13.25$ (26 neg.)          \\
            \midrule
            $17 \times 17$ & Symmetric (Subdomain Galerkin) & $\mathbf{3.10\%}$                          & $\mathbf{19.17\%}$                             & $1.26 \times 10^3$                                       & $+0.24$                     \\
            $17 \times 17$ & Unsymmetric (Row Petrov--Gal.)  & $3310.09\%$                                & $7870.43\%$                                    & $2.18 \times 10^5$                                       & $-12.73$ (48 neg.)          \\
            \bottomrule
        \end{tabular}
    }
\end{table}

The comparative results highlight a sharp contrast in algebraic stability:
\begin{itemize}
    \item \textbf{Patch Test vs.\ Non-Polynomial Fields:} Although both assemblies pass the affine patch test (Table~\ref{tab:patch_test_results}) because affine fields have zero local residual in $\mathcal{P}^1$, classical row-by-row assembly (testing each disk with a single function) loses positive-definiteness on non-polynomial fields. Even on $9 \times 9$, it exhibits $\min\mathrm{Re}(\lambda) = -14.19$ with 10 negative real and 26 complex eigenvalues, and conditioning deteriorates to $2.18 \times 10^5$ by $17 \times 17$ with diverging errors.
    \item \textbf{Positive-Definiteness under Subdomain Galerkin Assembly:} Testing against all six local basis functions on each disk yields symmetric positive-semidefinite stiffness and positive-definite mass blocks, ensuring $\mathbf{A} = \mathbf{K} + \kappa\mathbf{M} > 0$ ($\min \lambda_i > 0$) on each fixed mesh and monotonic $\mathcal{O}(h_{\mathrm{nom}}^2)$ refinement under exact local equilibrium.
\end{itemize}

\subsection{Stage 4: Subdomain Scaling Bounds and Parametric Sensitivity of $\alpha_s$}
\label{subsec:stage4_alpha_s}

The local integration radius is defined in \eqref{eq:subdomain_disk} as $r_s^I = \alpha_s h_I$. In scalar MLPG and LPIM \cite{atluri2002meshless, liu2009mesh}, $\alpha_s \in [0.75, 0.85]$ (commonly $\alpha_s = 0.80$) balances coverage against overlap. For 6-node directional vector stencils ($\mathcal{P}^1$), two geometric limits govern $\alpha_s$:
\begin{enumerate}
    \item \textbf{Lower Bound (Domain Coverage $\Omega_{\mathrm{cov}}$ \cite{atluri2002meshless, liu2009mesh}):} On a regular Cartesian lattice ($h_I = h_{\mathrm{nom}}$), the maximum distance from any interior point to its nearest node (the cell circumradius) is
          \begin{equation}
              r_{\mathrm{cov}} = \frac{\sqrt{2}}{2} h_{\mathrm{nom}} \approx 0.7071 h_{\mathrm{nom}},
              \label{eq:alpha_s_lower_bound}
          \end{equation}
          requiring $\alpha_s \ge \sqrt{2}/2$ for full ($100\%$) domain coverage. As shown in Fig.~\ref{fig:alpha_s_sensitivity}(a), $100\%$ coverage is preserved on regular grids for all $\alpha_s \ge 0.75$ (leaving a $13\%$ buffer at $\alpha_s = 0.80$), whereas below $0.7071$ coverage drops to $95.7\%$ ($\alpha_s = 0.60$), $78.9\%$ ($0.50$), and $29.4\%$ ($0.30$).
          
          Under $25\%$ nodal jitter, pairwise clustering shrinks nearest-neighbor $h_I$ in dense pairs while leaving neighboring zones rarefied, reducing coverage at $\alpha_s = 0.80$ to \textbf{$94.4\%$} ($5.6\%$ interstitial voids with $M(\mathbf{x}) = 0$). Crucially, the method remains well-posed and computationally stable at $94.4\%$ coverage because voids occur strictly between nodal clusters rather than isolating degrees of freedom. Every node $P_I$ retains a non-empty disk $D_I$, and the interpolation stencils $\Omega_u^I$ ($r_u^I \approx 1.15\text{--}1.20\,h_{\mathrm{nom}} > r_s^I$) remain strongly overlapping with non-collinear shared directors (second singular value $\sigma_2 \ge 1.0 \times 10^{-3}$ of the shared-director matrix across all 16 benchmark clouds up to $81\times 81$; Appendix~\ref{subsec:variational_consistency_strang}, Eq.~\eqref{eq:condensed_positive_definite_proof}). This guarantees $\mathbf{K}_{\mathrm{red}} > 0$, at the cost of the $\sim 0.6\%\text{--}1.3\%$ multiplicity-induced geometric floor quantified in Section~\ref{subsubsec:multiplicity_rayleigh_effect}. Full coverage on disordered clouds can be restored either via an adaptive $k$-nearest-neighbor spacing $h_I^{(k)} = \frac{1}{k}\sum_{j=1}^k \|\mathbf{x}_{(j)} - \mathbf{x}_I\|$ or moderate dilation $\alpha_s \in [0.85, 1.00]$ ($\ge 95\%\text{--}98\%$ coverage; Fig.~\ref{fig:alpha_s_sensitivity}(a)).
    \item \textbf{Upper Bound (6-Node Stencil Containment):} With nearest and diagonal neighbors, a 6-node stencil has mean support radius $r_u \approx 1.15\text{--}1.20\,h_{\mathrm{nom}}$. Whenever $r_s^I > r_u^I$, $\Omega_s^I$ spills outside its interpolation cloud, causing polynomial extrapolation error. Direct containment checks show zero boundary violations ($r_s^I \le r_u^I$) for $\alpha_s \le 1.10$ under both $0\%$ and $25\%$ jitter, with extrapolation onset around $\alpha_s \approx 1.15\text{--}1.20$.
\end{enumerate}

\begin{figure}[!t]
    \centering
    \includegraphics[width=\textwidth]{fig_mlpg_alpha_s_sensitivity.pdf}
    \caption{Parametric sensitivity of the subdomain radius factor $\alpha_s$ ($21\times 21$ mesh, random directors): (a) Domain integration coverage percentage $|\Omega_{\mathrm{cov}}| / |\Omega|$; (b) Relative broken $L^2(\vec{E})$ error in the 2D MMS problem across regular and perturbed clouds.}
    \label{fig:alpha_s_sensitivity}
\end{figure}

Figure~\ref{fig:alpha_s_sensitivity}(b) plots the resulting relative $L^2(\vec{E})$ error on the Stage~3 MMS test ($21 \times 21$, random directors). Within the operational band $\alpha_s \in [0.70, 1.20]$, accuracy remains steady ($2.94\%$ regular and $4.26\%$ under $25\%$ jitter at $\alpha_s = 0.80$), whereas $\alpha_s < 0.70$ suffers severe unintegrated void plateaus (up to $70\%\text{--}85\%$ uncovered) and $\alpha_s > 1.20$ degrades through polynomial extrapolation ($4.71\%$ and $5.46\%$ at $\alpha_s = 2.00$). Hence $\alpha_s = 0.80$ ($r_s^I \approx 0.72\,r_u^I$) provides a balanced baseline mirroring $\widehat{\mathrm{Tol}} = 0.80$ \eqref{eq:scale_invariant_det_criterion} and classical MLPG/LPIM practice \cite{atluri2002meshless, liu2009mesh}.

\begin{remark}[Geometric Invariance of $\alpha_s$ under Energy Stabilization]
    Because the bounds \eqref{eq:alpha_s_lower_bound} and $r_s^I \le r_u^I$ depend solely on nodal geometry and stencil containment, repeating Fig.~\ref{fig:alpha_s_sensitivity}(b) with $\mathbf{K}_{\mathrm{stab}}$ active alters the error profile by less than $0.3\%$ across $\alpha_s \in [0.70, 1.20]$. This check was evaluated for $\alpha_{\mathrm{stab}} \in [0.175, 0.20]$ (covering the calibrated baseline $c_{\mathrm{stab}} = 4.0$ at $21\times 21$, where $h_{\mathrm{nom}}/L = 1/20$), leaving $\alpha_s = 0.80$ unchanged as the baseline.
\end{remark}

\subsection{Stage 5: Physical Spectral Validation and Divergence Penalty Scaling}
\label{subsec:stage5_cavity}

With $\alpha_s = 0.80$ selected as the baseline subdomain factor in Section~\ref{subsec:stage4_alpha_s}, the physical spectral fidelity of VNMM-MLPG is evaluated on the 2D PEC resonant cavity $\Omega = [0, \pi]^2\text{ m}^2$ ($L_x = L_y = \pi\text{ m}$) with $\mu_r = 1, \epsilon_r = 1$. The analytical cutoff wavenumbers for transverse electric modes $TE_{nm}$ evaluate to:
\begin{equation}
    k_{c, nm} = \sqrt{\lambda_{nm}} = \sqrt{n^2 + m^2}\text{ m}^{-1}, \qquad n, m \in \mathbb{N}_0, \quad (n, m) \ne (0, 0).
    \label{eq:analytical_cavity_kc}
\end{equation}
Global stiffness and mass operators $\mathbf{K}_{\mathrm{red}}$ and $\mathbf{M}_{\mathrm{red}}$ are assembled via the Symmetric Subdomain Galerkin scheme \eqref{eq:global_assembly}, ensuring real eigenvalues ($\lambda = k_c^2 \in \mathbb{R}$). Physical solenoidal modes are distinguished from spurious gradient modes via the divergence energy ratio matching \cite{stable_vnmm2d_2026}:
\begin{equation}
    R_{\mathrm{div}}(\mathbf{e}) = \frac{\mathbf{e}^T \mathbf{K}_{\mathrm{div}} \mathbf{e}}{\mathbf{e}^T \mathbf{K}_{\mathrm{curl}} \mathbf{e}} \times 100\%.
    \label{eq:R_div_definition}
\end{equation}
Physical solenoidal modes satisfy $R_{\mathrm{div}} \le 50\%$, whereas unregularized gradient modes exhibit $R_{\mathrm{div}} \gg 100\%$. Spectral accuracy across the first 10 physical modes is quantified by the mean relative cutoff wavenumber error $\overline{\varepsilon}_{k_c} \equiv \frac{1}{10}\sum_{m=1}^{10} |k_{c, \mathrm{num}}^{(m)} - k_{c, \mathrm{ref}}^{(m)}| / k_{c, \mathrm{ref}}^{(m)}$. In all simulations below, isotropic random vector directors ($\text{seed}=42$) and $\alpha_s = 0.80$ are employed.

\subsubsection{Parametric Analysis of $s_{\mathrm{div}}$ on the Baseline $21\times 21$ Grid}
\label{subsubsec:s_div_sweep}

We first investigate the sensitivity of the divergence regularization penalty $s_{\mathrm{div}}$ on the baseline $21 \times 21$ grid ($N = 441$ nodes, 361 active free interior DoFs, nominal spacing $h_{\mathrm{nom}} = \pi/20 \approx 0.1571\text{ m}$). In previous cell-based formulations \cite{stable_vnmm2d_2026}, a penalty of $s_{\mathrm{div}} = 6.0$ was adopted as standard. To systematically assess its efficacy in the cell-free MLPG setting, a parametric sweep across $s_{\mathrm{div}} \in [0.0, 50.0]$ was performed, as documented in Table~\ref{tab:s_div_sweep}.

In computational electromagnetics, non-physical spurious gradient modes originating from discrete nullspace leakage ($\nabla_h \times \nabla_h \phi_h \neq \mathbf{0}$) scale linearly with the divergence penalty ($d\lambda_{\mathrm{spur}} / ds_{\mathrm{div}} \gg 0$) and are shifted rapidly toward high frequencies \cite{jin2014finite, costabel2002weighted, bossavit1990spurious}. Conversely, physical cavity resonances stabilize asymptotically near their analytical cutoff values ($d\lambda / ds_{\mathrm{div}} \to 0$).

As reported in Table~\ref{tab:s_div_sweep}:
\begin{itemize}
    \item At $s_{\mathrm{div}} = 0.0$ (unregularized), the stiffness matrix $\mathbf{K}_{\mathrm{curl}}$ is numerically singular; all initial 10 numerical modes are spurious gradient fields with dominant divergence energy ($\max R_{\mathrm{div}} > 10^4\%$), giving an apparent error of $83.23\%$.
    \item For $s_{\mathrm{div}} \ge 4.0$, all 10 computed states correspond to physical cavity modes, with the lowest spurious mode shifted outside the operational band ($\lambda_{\mathrm{spur}} > 78$).
    \item Increasing the penalty from the cell-based standard $s_{\mathrm{div}} = 6.0$ ($8.83\%$ mean error) \cite{stable_vnmm2d_2026} to $s_{\mathrm{div}} = 15.0$ reduces the divergence residue of physical modes (down to $\max R_{\mathrm{div}} = 3.4\%$). Across the first 10 physical modes, it achieves a minimum mean error of $\mathbf{2.46\%}$, a maximum error below $4.95\%$, and stable conditioning ($\mathrm{cond} = 5.14 \times 10^4$). The slight non-monotonicity in ``Lowest Spur $\lambda$'' between $s_{\mathrm{div}} = 10.0$ ($\lambda = 252.60$) and $s_{\mathrm{div}} = 15.0$ ($\lambda = 237.24$) occurs far above the operational band ($\lambda \le 9$, near the 80th discrete mode). Near that Nyquist limit, ascending gradient modes hybridize with dispersive higher-order cavity modes across the fixed $R_{\mathrm{div}} = 50\%$ threshold.
    \item For $s_{\mathrm{div}} \ge 30.0$, over-penalization begins to slightly elevate the mean error ($4.83\%$ at $s_{\mathrm{div}} = 30.0$ and $8.65\%$ at $s_{\mathrm{div}} = 50.0$).
\end{itemize}
Consequently, $s_{\mathrm{div}} = 15.0$ is adopted as the calibrated baseline for the $21 \times 21$ discretization.

\begin{table}[!t]
    \centering
    \small
    \caption{Parametric sensitivity of the div-curl penalty parameter $s_{\mathrm{div}}$ on spurious mode shifting and cutoff accuracy ($21 \times 21$ grid, $\alpha_s = 0.80$, isotropic random directors, criterion $R_{\mathrm{div}} > 50\%$). The recommended calibration $s_{\mathrm{div}} = 15.0$ and the previous cell-based standard $s_{\mathrm{div}} = 6.0$ are highlighted.}
    \label{tab:s_div_sweep}
    \resizebox{\textwidth}{!}{
        \begin{tabular}{ccccccc}
            \toprule
            $\boldsymbol{s_{\mathrm{div}}}$ & \textbf{Spur in Band} & \textbf{Lowest Spur} $\boldsymbol{\lambda}$ & \textbf{Mean Error} $\boldsymbol{k_c}$ \textbf{(\%)} & \textbf{Max Error} $\boldsymbol{k_c}$ \textbf{(\%)} & $\boldsymbol{\mathrm{cond}(\mathbf{K}_{\mathrm{red}})}$ & \textbf{Max} $\boldsymbol{R_{\mathrm{div}}}$ \textbf{(\%)} \\
            \midrule
            0.0                             & 10                    & 0.05                                        & $83.23\%$                                            & $87.39\%$                                           & Singular                                                & $> 10^4\%$                                                 \\
            1.0                             & 3                     & 1.23                                        & $27.82\%$                                            & $40.83\%$                                           & $9.02 \times 10^3$                                      & $466.0\%$                                                  \\
            2.0                             & 3                     & 2.17                                        & $19.00\%$                                            & $28.74\%$                                           & $9.67 \times 10^3$                                      & $180.9\%$                                                  \\
            4.0                             & 0                     & 78.39                                       & $12.19\%$                                            & $25.31\%$                                           & $1.58 \times 10^4$                                      & $41.5\%$                                                   \\
            6.0 (\cite{stable_vnmm2d_2026}) & 0                     & 173.66                                      & $8.83\%$                                             & $21.40\%$                                           & $2.25 \times 10^4$                                      & $30.2\%$                                                   \\
            8.0                             & 0                     & 207.15                                      & $6.56\%$                                             & $13.70\%$                                           & $2.91 \times 10^4$                                      & $21.9\%$                                                   \\
            10.0                            & 0                     & 252.60                                      & $4.82\%$                                             & $7.51\%$                                            & $3.56 \times 10^4$                                      & $15.1\%$                                                   \\
            \textbf{15.0 (Recommended)}     & \textbf{0}            & \textbf{237.24}                             & \textbf{2.46\%}                                      & \textbf{4.95\%}                                     & $\mathbf{5.14 \times 10^4}$                             & \textbf{3.4\%}                                             \\
            20.0                            & 0                     & 471.77                                      & $2.85\%$                                             & $5.81\%$                                            & $6.69 \times 10^4$                                      & $2.2\%$                                                    \\
            30.0                            & 0                     & 693.46                                      & $4.83\%$                                             & $10.82\%$                                           & $9.68 \times 10^4$                                      & $1.0\%$                                                    \\
            50.0                            & 0                     & 1035.15                                     & $8.65\%$                                             & $18.03\%$                                           & $1.54 \times 10^5$                                      & $0.8\%$                                                    \\
            \bottomrule
        \end{tabular}
    }
\end{table}

\subsubsection{Cavity Spectrum and Mode-by-Mode Accuracy at $s_{\mathrm{div}} = 15.0$}
\label{subsubsec:cavity_modes_15}

Using the recommended penalty $s_{\mathrm{div}} = 15.0$, Table~\ref{tab:cavity_eigenvalues} presents the spectrum for the first 10 numerical transverse electric modes on the $21 \times 21$ mesh, compared against the exact analytical benchmark from Eq.~\eqref{eq:analytical_cavity_kc}.

\begin{table}[!t]
    \centering
    \small
    \caption{Comparison of the first 10 numerical cavity eigenvalues ($\lambda$ in $\mathrm{m}^{-2}$) and cutoff wavenumbers ($k_c$ in $\mathrm{m}^{-1}$) computed by VNMM-MLPG ($21 \times 21, s_{\mathrm{div}} = 15.0, \alpha_s = 0.80$, isotropic random directors) against analytical reference solutions from Eq.~\eqref{eq:analytical_cavity_kc} on a cavity with dimensions $\pi\text{ m} \times \pi\text{ m}$.}
    \label{tab:cavity_eigenvalues}
    \resizebox{\textwidth}{!}{
        \begin{tabular}{ccccccc}
            \toprule
            \textbf{Mode}                                                                    & $\boldsymbol{\lambda_{\mathrm{num}}}$ \textbf{[m$^{-2}$]} & $\boldsymbol{\lambda_{\mathrm{ref}}}$ \textbf{[m$^{-2}$]} & $\boldsymbol{k_{c, \mathrm{num}}}$ \textbf{[m$^{-1}$]} & $\boldsymbol{k_{c, \mathrm{ref}}}$ \textbf{[m$^{-1}$]} & \textbf{Error} $\boldsymbol{k_c}$ \textbf{(\%)} & $\boldsymbol{R_{\mathrm{div}}}$ \textbf{(\%)} \\
            \midrule
            $TE_{10}$                                                                        & 0.9035                                                    & 1.0000                                                    & 0.9505                                                 & 1.0000                                                 & \textbf{4.95\%}                                 & $0.63\%$                                      \\
            $TE_{01}$                                                                        & 0.9213                                                    & 1.0000                                                    & 0.9598                                                 & 1.0000                                                 & \textbf{4.02\%}                                 & $0.53\%$                                      \\
            $TE_{11}$                                                                        & 1.9366                                                    & 2.0000                                                    & 1.3916                                                 & 1.4142                                                 & \textbf{1.60\%}                                 & $0.67\%$                                      \\
            $TE_{20}$                                                                        & 3.7689                                                    & 4.0000                                                    & 1.9414                                                 & 2.0000                                                 & \textbf{2.93\%}                                 & $1.19\%$                                      \\
            $TE_{02}$                                                                        & 3.9924                                                    & 4.0000                                                    & 1.9981                                                 & 2.0000                                                 & \textbf{0.10\%}                                 & $0.91\%$                                      \\
            $TE_{21}$                                                                        & 4.6468                                                    & 5.0000                                                    & 2.1556                                                 & 2.2361                                                 & \textbf{3.60\%}                                 & $1.05\%$                                      \\
            $TE_{12}$                                                                        & 5.3363                                                    & 5.0000                                                    & 2.3101                                                 & 2.2361                                                 & \textbf{3.31\%}                                 & $1.19\%$                                      \\
            $TE_{22}$                                                                        & 8.2883                                                    & 8.0000                                                    & 2.8789                                                 & 2.8284                                                 & \textbf{1.79\%}                                 & $1.82\%$                                      \\
            $TE_{30}$                                                                        & 8.6164                                                    & 9.0000                                                    & 2.9354                                                 & 3.0000                                                 & \textbf{2.15\%}                                 & $3.43\%$                                      \\
            $TE_{03}$                                                                        & 9.0222                                                    & 9.0000                                                    & 3.0037                                                 & 3.0000                                                 & \textbf{0.12\%}                                 & $3.20\%$                                      \\
            \midrule
            \multicolumn{5}{l}{\textbf{Mean Relative Error across First 10 Physical Modes:}} & \textbf{2.46\%}                                           & ---                                                                                                                                                                                                                                                                           \\
            \multicolumn{5}{l}{\textbf{Spurious Modes in Operational Band ($\le 10$):}}      & \textbf{0}                                                & ---                                                                                                                                                                                                                                                                           \\
            \bottomrule
        \end{tabular}
    }
\end{table}

The mode-by-mode spectrum in Table~\ref{tab:cavity_eigenvalues} shows close agreement with theory:
\begin{itemize}
    \item The fundamental degenerate pair ($TE_{10}$ and $TE_{01}$, exact $k_c = 1.0\text{ m}^{-1}$) achieves cutoff errors of $4.95\%$ and $4.02\%$ ($R_{\mathrm{div}} \approx 0.5\%\text{--}0.6\%$). The slight doublet splitting arises because randomized interior directors ($\text{seed}=42$) break the $C_{4v}$ rotational symmetry of the square lattice and vanishes as $h_{\mathrm{nom}} \to 0$.
    \item Symmetric modes $TE_{02}$ and $TE_{03}$ achieve errors of $0.10\%$ and $0.12\%$, while $TE_{22}$ improves from $21.40\%$ ($R_{\mathrm{div}} = 30.18\%$) at $s_{\mathrm{div}} = 6.0$ to $\mathbf{1.79\%}$ ($R_{\mathrm{div}} = 1.82\%$).
\end{itemize}

To complement Table~\ref{tab:cavity_eigenvalues} with a qualitative field check under regular nodal distributions ($0\%$ jitter, $\alpha_s = 0.80$), Fig.~\ref{fig:modal_fields_21x21} compares the analytical electric-field profiles against the multiplicity-weighted global VNMM-MLPG reconstructions $\vec{E}_{\mathrm{global}}$ \eqref{eq:global_field_reconstruction} for three representative modes ($TE_{10}$, $TE_{11}$, and $TE_{21}$). Both the $21\times 21$ grid ($s_{\mathrm{div}} = 15.0$, $\alpha_{\mathrm{stab}} = 0$) and the refined $81\times 81$ grid (using the calibrated scalings of Sections~\ref{subsubsec:constraint_ratio_scaling} and~\ref{subsubsec:resolution_stagnation_stabilization}) are shown, with the degenerate doublets $TE_{10}/TE_{01}$ and $TE_{21}/TE_{12}$ aligned within their respective two-dimensional eigenspaces. Across all three modes, the cell-free reconstruction reproduces the analytical standing-wave lobes, nodal lines, tangential PEC boundary behavior, and vector field orientations cleanly without spurious checkerboard artifacts, with the finer $81\times 81$ discretization showing slightly sharper modal contours.

\begin{figure}[!t]
    \centering
    \includegraphics[width=\textwidth]{fig_mlpg_modal_fields.pdf}
    \caption{Visual comparison among exact analytical cavity modes (top row) and multiplicity-weighted global VNMM-MLPG reconstructions $\vec{E}_{\mathrm{global}}$ \eqref{eq:global_field_reconstruction} on regular grids ($0\%$ jitter, $\alpha_s = 0.80$): $21\times 21$ (middle row, $s_{\mathrm{div}} = 15.0$) and $81\times 81$ (bottom row, calibrated scaled parameters) for $TE_{10}$, $TE_{11}$, and $TE_{21}$.} 
    \label{fig:modal_fields_21x21}
\end{figure}

\subsubsection{Spectral Stagnation under Fixed Divergence Penalty}
\label{subsubsec:stagnation_fixed_sdiv}

Next, we examine grid refinement across the six benchmark resolutions from \cite{stable_vnmm2d_2026} ($15 \times 15$ to $49 \times 49$) with isotropic random directors ($\text{seed}=42$). Holding the penalty fixed at $s_{\mathrm{div}} = 15.0$, Table~\ref{tab:cavity_scaling_errors} reveals that after decreasing from $6.46\%$ ($15^2$) to $3.84\%$ ($22^2$), the mean cutoff error stagnates between $5.16\%$ and $7.32\%$ for $N_x \ge 29$ (mirroring the $\approx 11\%\text{--}13\%$ plateau at $s_{\mathrm{div}} = 6.0$).

\subsubsection{Constraint-to-DoF Ratios and Asymptotic Divergence Penalty Scaling}
\label{subsubsec:constraint_ratio_scaling}

To understand the origin of this spectral stagnation, we observe that in cell-based VNMM-Gauss \cite{stable_vnmm2d_2026}, $3 \times 3$ Gauss points per cell sample distinct neighboring stencils, providing an over-constraint ratio of $\approx 9:1$ that stiffens the curl operator. In contrast, cell-free MLPG evaluates exactly one disk per node, yielding an effective constraint-to-DoF ratio of only $\approx 1.1:1$. Consequently, keeping $s_{\mathrm{div}}$ fixed while refining the cloud allows discrete gradient modes in the expanding irrotational nullspace to soften and pollute the physical spectrum.

To maintain spectral separation across the tested resolutions without inducing volumetric locking, we adopt an empirical superlinear power-law scaling $s_{\mathrm{div}} \propto (h_{\mathrm{char}} / L)^{-\beta}$ with $\beta = 1.5$. This calibration ensures that the penalty grows sufficiently fast to push spurious gradient modes out of the lower operational spectrum while keeping penalty energy perturbations on physical solenoidal modes small. Across all benchmark experiments, $s_{\mathrm{div}}$ is evaluated in scale-invariant dimensionless form as a uniform global scalar:
\begin{equation}
    s_{\mathrm{div}} = s_{\mathrm{div}}^{(0)} \left( \frac{\hat{h}_0}{h_{\mathrm{char}} / L} \right)^{1.5} = c_{\mathrm{div}} \left( \frac{h_{\mathrm{char}}}{L} \right)^{-1.5},
    \label{eq:s_div_scaling_global}
\end{equation}
where $\hat{h}_0 = 1/20$ is the dimensionless reference resolution (matching $h_0 = L/20 = \pi/20\text{ m}$ on the $21\times 21$ baseline cavity with $L = \pi\text{ m}$), $s_{\mathrm{div}}^{(0)} = 15.0$, and $c_{\mathrm{div}} = s_{\mathrm{div}}^{(0)} \hat{h}_0^{1.5} \approx 0.1677$, ensuring geometric scale invariance across cavity dimensions $L$ consistent with \eqref{eq:alpha_stab_vms_scaling}. In the perturbed-lattice benchmarks of Sections~\ref{subsubsec:cavity_jitter_error}--\ref{subsec:stage6_benchmark} (constructed by perturbing $N_x \times N_y$ grids), both $s_{\mathrm{div}}$ and $\alpha_{\mathrm{stab}}$ are evaluated using the nominal grid spacing $h_{\mathrm{nom}} = L/(N_x - 1)$ so that regular and $25\%$-jitter configurations share identical global penalty parameters at each resolution. On general unstructured clouds without a parent lattice, the domain-averaged spacing $h_{\mathrm{char}} = \overline{h}$ applies instead.

\subsubsection{Spectral Error Reduction and Stagnation across Regular and Disordered Clouds}
\label{subsubsec:cavity_jitter_error}

To evaluate the baseline cell-free VNMM-MLPG formulation, the cavity eigenvalue problem was first solved across the six benchmark grid resolutions on regular nodal distributions with isotropic random directors ($\alpha_s = 0.80$, $\text{seed}=42$), applying the superlinear scaling law \eqref{eq:s_div_scaling_global} (yielding $s_{\mathrm{div}} \in [8.78, 55.77]$). Table~\ref{tab:cavity_scaling_errors} contrasts the resulting cutoff wavenumber errors against the fixed penalty $s_{\mathrm{div}} = 15.0$.

\begin{table}[!t]
    \centering
    \small
    \caption{Cavity eigenvalue cutoff wavenumber error comparison on the 2D PEC resonant cavity across the six grid resolutions from \cite{stable_vnmm2d_2026}: fixed penalty ($s_{\mathrm{div}} = 15.0$) vs. superlinear scaled penalty \eqref{eq:s_div_scaling_global} ($s_{\mathrm{div}} \in [8.78, 55.77]$) with isotropic random directors ($\alpha_s = 0.80$, $\text{seed}=42$).}
    \label{tab:cavity_scaling_errors}
    \resizebox{\textwidth}{!}{
        \begin{tabular}{ccccccc}
            \toprule
            \textbf{Grid}             & \textbf{Nominal Spacing}                     & \textbf{Active} & \multicolumn{2}{c}{\textbf{Fixed Penalty ($s_{\mathrm{div}} = 15.0$)}} & \multicolumn{2}{c}{\textbf{Superlinear Scaled Penalty ($s_{\mathrm{div}} \propto (h_{\mathrm{nom}}/L)^{-1.5}$)}}                                                                                             \\
            \cmidrule(lr){4-5} \cmidrule(lr){6-7}
            $\boldsymbol{N_x \times N_y}$ & $\boldsymbol{h_{\mathrm{nom}}}$ \textbf{[m]} & \textbf{DoFs}   & $\boldsymbol{s_{\mathrm{div}}}$                                        & \textbf{Mean Error} $\boldsymbol{k_c}$ \textbf{(\%)}                                          & $\boldsymbol{s_{\mathrm{div}}(h_{\mathrm{nom}})}$ & \textbf{Mean Error} $\boldsymbol{k_c}$ \textbf{(\%)} \\
            \midrule
            $15 \times 15$            & 0.2244                                       & 169             & 15.0                                                                   & $6.46\%$                                                                                      & 8.78                               & $\mathbf{3.87\%}$                                    \\
            $22 \times 22$            & 0.1496                                       & 400             & 15.0                                                                   & $3.84\%$                                                                                      & 16.14                              & $\mathbf{3.55\%}$                                    \\
            $29 \times 29$            & 0.1122                                       & 729             & 15.0                                                                   & $5.27\%$                                                                                      & 24.85                              & $\mathbf{2.58\%}$                                    \\
            $36 \times 36$            & 0.0898                                       & 1156            & 15.0                                                                   & $5.16\%$                                                                                      & 34.73                              & $\mathbf{2.28\%}$                                    \\
            $43 \times 43$            & 0.0748                                       & 1681            & 15.0                                                                   & $6.44\%$                                                                                      & 45.65                              & $\mathbf{2.03\%}$                                    \\
            $49 \times 49$            & 0.0654                                       & 2209            & 15.0                                                                   & $7.32\%$                                                                                      & 55.77                              & $\mathbf{2.32\%}$                                    \\
            \bottomrule
        \end{tabular}
    }
\end{table}

As shown in Table~\ref{tab:cavity_scaling_errors}, adopting the superlinear $h_{\mathrm{nom}}$-scaled penalty avoids coarse-grid over-stiffening ($3.87\%$ vs.\ $6.46\%$ on $15 \times 15$ at $s_{\mathrm{div}} = 8.78$) and lowers intermediate-to-fine errors from the $5.1\%\text{--}7.3\%$ fixed-penalty plateau down to $2.03\%\text{--}3.55\%$ across $22\times 22\text{--}49\times 49$. Nevertheless, the error still levels off around $\approx 2.0\%\text{--}2.3\%$ on dense grids, yielding a sublinear empirical rate of $p = 0.53$.

To evaluate the robustness of this unaugmented formulation under geometric irregularity, stochastic perturbations were next applied to all interior nodal coordinates:
\begin{equation}
    \mathbf{x}_i \longleftarrow \mathbf{x}_i + \delta\mathbf{x}_i, \quad \delta x_i, \delta y_i \sim \mathcal{U}(-j_{\mathrm{frac}} h_{\mathrm{nom}}, +j_{\mathrm{frac}} h_{\mathrm{nom}}),
\end{equation}
where $j_{\mathrm{frac}} \in \{0.0, 0.25\}$ denotes the jitter fraction relative to the nominal spacing $h_{\mathrm{nom}}$. Boundary nodes $\mathbf{x} \in \partial\Omega$ remain unperturbed, preserving geometric enforcement of the PEC boundary conditions ($E_t = 0$).

Fig.~\ref{fig:mlpg_cavity_jitter} compares the resulting mean cutoff errors under $0\%$ and $25\%$ jitter using the same superlinear penalty \eqref{eq:s_div_scaling_global} and baseline parameters of Table~\ref{tab:cavity_scaling_errors}.

\begin{figure}[!t]
    \centering
    \includegraphics[width=0.68\textwidth]{fig_mlpg_cavity_jitter.pdf}
    \caption{Log--log scaling of the unstabilized mean cutoff error $\overline{\varepsilon}_{k_c}$ versus $h_{\mathrm{nom}}$ across $15\times 15$ to $49\times 49$ under $0\%$ (Table~\ref{tab:cavity_scaling_errors}) and $25\%$ interior nodal jitter.}
    \label{fig:mlpg_cavity_jitter}
\end{figure}

The error trends in Fig.~\ref{fig:mlpg_cavity_jitter} reveal distinct scaling regimes:
\begin{itemize}
    \item \textbf{Regular Distribution ($0\%$ Jitter, $p = 0.53$):} As summarized in Table~\ref{tab:cavity_scaling_errors}, local harmonic softening limits the superlinearly penalized scaling to a sublinear plateau around $2.0\%\text{--}2.3\%$ on dense grids.
    \item \textbf{Moderate Perturbation ($25\%$ Jitter, $p = 0.84$):} Under $25\%$ random dislocation with $\alpha_s = 0.80$, mean errors drop monotonically from $12.67\%$ ($15\times 15$) to $\mathbf{3.95\%}$ ($49\times 49$) with zero spurious modes in the 10-mode operational band.
\end{itemize}
Crucially, even with superlinear divergence scaling \eqref{eq:s_div_scaling_global}, cutoff convergence remains sublinear ($p = 0.53$ regular, $p = 0.84$ jittered) and stalls above $2\%\text{--}4\%$ on $49\times 49$: penalizing $\nabla \cdot \vec{E}$ suppresses irrotational leakage but leaves the two harmonic linear gradient modes in $\ker(\mathbf{K}_I^{\mathrm{local}})$ uncontrolled (Section~\ref{subsubsec:local_rank_analysis}).

\subsubsection{Evaluation of Orthogonal Energy Stabilization and Coupled $(\beta, \gamma)$ Calibration}
\label{subsubsec:resolution_stagnation_stabilization}

To overcome this residual stagnation by restoring rank 4 over the local gradient subspace, we now activate Orthogonal Energy Stabilization \eqref{eq:k_stab_analytic}--\eqref{eq:alpha_stab_vms_scaling} to answer two questions: whether the stabilized operator breaks through the $2\%\text{--}4\%$ plateau to sustain convergence below $1\%$ on finer point clouds, and whether the divergence exponent $\beta = 1.5$ interacts adversely with $\alpha_{\mathrm{stab}}$. For this study, the refinement sequence is extended by two fine levels: $65 \times 65$ ($h_{\mathrm{nom}} = 0.0491\text{ m}$, $3,969$ active DoFs) and $81 \times 81$ ($h_{\mathrm{nom}} = 0.0393\text{ m}$, $6,241$ active DoFs). We compare the unstabilized operator ($c_{\mathrm{stab}} = 0$), quadratic stabilization ($\gamma = 2$, $c_2 = 70$), and linear stabilization ($\gamma = 1$, $c_{\mathrm{stab}} \in \{1.0, 2.0, 4.0\}$) in Table~\ref{tab:cavity_convergence_stabilization} and Fig.~\ref{fig:mlpg_cavity_stabilization}, while Table~\ref{tab:beta_gamma_fine_mesh} examines the coupled $(\beta, \gamma, c_{\mathrm{stab}})$ parameter space across all eight grids. Across all 240 cavity runs summarized in Tables~\ref{tab:cavity_convergence_stabilization}--\ref{tab:beta_gamma_fine_mesh} (and all 654 unique parameter-grid evaluations across the full $(\beta, \gamma, c_{\mathrm{stab}})$ sweeps), zero spurious modes enter the 10-mode operational band ($\mathrm{spur} = 0$).

\begin{table}[!t]
    \centering
    \small
    \caption{Extended cavity cutoff convergence ($15 \times 15$ to $81 \times 81$, $\alpha_s = 0.80$, $\text{seed}=42$, $\beta = 1.5$) comparing the unstabilized operator ($c_{\mathrm{stab}} = 0$), quadratic stabilization ($\gamma = 2, c_2 = 70$), and linear stabilization ($\gamma = 1$, $\alpha_{\mathrm{stab}} = c_{\mathrm{stab}}(h_{\mathrm{nom}}/L)$ for $c_{\mathrm{stab}} \in \{1.0, 2.0, 4.0\}$) on regular ($0\%$ jitter) and perturbed ($25\%$ jitter) clouds.}
    \label{tab:cavity_convergence_stabilization}
    \resizebox{\textwidth}{!}{
        \begin{tabular}{ccccccc}
            \toprule
            \textbf{Grid}                                                                                     & \textbf{Nominal Spacing}                     & \textbf{Unstabilized}                  & \textbf{Quadratic ($\gamma=2$)} & \multicolumn{3}{c}{\textbf{Linear Stabilization ($\gamma = 1$, $\boldsymbol{\alpha_{\mathrm{stab}} = c_{\mathrm{stab}}(h_{\mathrm{nom}}/L)}$)}} \\
            \cmidrule(lr){5-7}
            $\boldsymbol{N_x \times N_y}$                                                                     & $\boldsymbol{h_{\mathrm{nom}}}$ \textbf{[m]} & $\boldsymbol{c_{\mathrm{stab}} = 0}$   & $\boldsymbol{c_2 = 70}$         & $\boldsymbol{c_{\mathrm{stab}} = 1.0}$ & $\boldsymbol{c_{\mathrm{stab}} = 2.0}$ & $\boldsymbol{c_{\mathrm{stab}} = 4.0}$ \\
            \midrule
            \multicolumn{7}{l}{\textit{Regular Lattice ($0\%$ Jitter)}} \\
            $15 \times 15$                                                                                    & 0.2244                                       & $3.87\%$                               & $17.67\%$                       & $5.42\%$                               & $8.53\%$                               & $14.58\%$                              \\
            $22 \times 22$                                                                                    & 0.1496                                       & $3.55\%$                               & $7.11\%$                        & $3.33\%$                               & $4.56\%$                               & $8.35\%$                               \\
            $29 \times 29$                                                                                    & 0.1122                                       & $2.58\%$                               & $3.34\%$                        & $1.84\%$                               & $2.59\%$                               & $6.19\%$                               \\
            $36 \times 36$                                                                                    & 0.0898                                       & $2.28\%$                               & $2.87\%$                        & $2.17\%$                               & $2.87\%$                               & $5.64\%$                               \\
            $43 \times 43$                                                                                    & 0.0748                                       & $2.03\%$                               & $1.39\%$                        & $1.34\%$                               & $1.48\%$                               & $3.50\%$                               \\
            $49 \times 49$                                                                                    & 0.0654                                       & $2.32\%$                               & $1.29\%$                        & $1.50\%$                               & $\mathbf{1.16\%}$                      & $2.39\%$                               \\
            $65 \times 65$                                                                                    & 0.0491                                       & $2.60\%$                               & $1.69\%$                        & $1.75\%$                               & $1.09\%$                               & $\mathbf{1.03\%}$                      \\
            $81 \times 81$                                                                                    & 0.0393                                       & $2.12\%$                               & $1.52\%$                        & $1.44\%$                               & $0.89\%$                               & $\mathbf{0.78\%}$                      \\
            \midrule
            \multicolumn{2}{l}{\textbf{Rates ($0\%$ Jitter):} $p_{15\text{--}49}$ / $p_{15\text{--}81}$ / $p_{36\text{--}81}$} & $0.53 / 0.33 / -0.05$                  & $2.15 / 1.46 / 0.44$            & $1.09 / 0.73 / 0.24$                   & $1.55 / 1.33 / 1.22$                   & $1.35 / \mathbf{1.72} / \mathbf{2.48}$ \\
            \midrule
            \multicolumn{7}{l}{\textit{Perturbed Point Cloud ($25\%$ Jitter, $\text{seed}=42$)}} \\
            $15 \times 15$                                                                                    & 0.2244                                       & $12.67\%$                              & $11.88\%$                       & $5.78\%$                               & $\mathbf{3.68\%}$                      & $8.80\%$                               \\
            $22 \times 22$                                                                                    & 0.1496                                       & $5.70\%$                               & $6.82\%$                        & $\mathbf{3.75\%}$                      & $4.51\%$                               & $8.27\%$                               \\
            $29 \times 29$                                                                                    & 0.1122                                       & $5.61\%$                               & $3.77\%$                        & $4.13\%$                               & $\mathbf{3.67\%}$                      & $5.28\%$                               \\
            $36 \times 36$                                                                                    & 0.0898                                       & $5.15\%$                               & $\mathbf{2.36\%}$               & $3.34\%$                               & $\mathbf{2.36\%}$                      & $2.73\%$                               \\
            $43 \times 43$                                                                                    & 0.0748                                       & $4.27\%$                               & $1.77\%$                        & $2.60\%$                               & $\mathbf{1.57\%}$                      & $2.13\%$                               \\
            $49 \times 49$                                                                                    & 0.0654                                       & $3.95\%$                               & $1.97\%$                        & $2.58\%$                               & $1.50\%$                               & $\mathbf{1.46\%}$                      \\
            $65 \times 65$                                                                                    & 0.0491                                       & $3.48\%$                               & $2.43\%$                        & $2.52\%$                               & $1.58\%$                               & $\mathbf{0.85\%}$                      \\
            $81 \times 81$                                                                                    & 0.0393                                       & $3.48\%$                               & $2.82\%$                        & $2.72\%$                               & $1.99\%$                               & $\mathbf{0.79\%}$                      \\
            \midrule
            \multicolumn{2}{l}{\textbf{Rates ($25\%$ Jitter):} $p_{15\text{--}49}$ / $p_{15\text{--}81}$ / $p_{36\text{--}81}$} & $0.84 / 0.68 / 0.46$                   & $1.62 / 0.94 / -0.36$           & $0.63 / 0.46 / 0.19$                   & $0.87 / 0.62 / 0.11$                   & $1.54 / \mathbf{1.61} / \mathbf{1.62}$ \\
            \bottomrule
        \end{tabular}
    }
\end{table}

\begin{figure}[!t]
    \centering
    \includegraphics[width=\textwidth]{fig_mlpg_cavity_stabilization.pdf}
    \caption{Effect of Orthogonal Energy Stabilization on cavity cutoff convergence (data in Table~\ref{tab:cavity_convergence_stabilization}, $15\times 15$ to $81\times 81$, $\beta = 1.5$): (a) regular grids ($0\%$ jitter); (b) perturbed point clouds ($25\%$ jitter).}
    \label{fig:mlpg_cavity_stabilization}
\end{figure}

\begin{table}[!t]
    \centering
    \small
    \caption{Coupled sensitivity of the divergence exponent $\beta$ ($s_{\mathrm{div}} \propto (h_{\mathrm{nom}}/L)^{-\beta}$) and stabilization exponent $\gamma$ ($\alpha_{\mathrm{stab}} = \alpha_0 [h_{\mathrm{nom}} / (L\hat{h}_0)]^{\gamma}$, $\hat{h}_0 = 1/20$, i.e., $c_{\mathrm{stab}} = \alpha_0 \hat{h}_0^{-\gamma}$) across $15\times 15$ to $81\times 81$ ($\alpha_s = 0.80$, $\text{seed}=42$; at $\gamma = 1$, $\alpha_0 = 0.05, 0.10, 0.20$ give $c_{\mathrm{stab}} = 1.0, 2.0, 4.0$).}
    \label{tab:beta_gamma_fine_mesh}
    \resizebox{\textwidth}{!}{
        \begin{tabular}{ccccccccc}
            \toprule
            \textbf{Regime / Coupling} & $\boldsymbol{\beta}$ & $\boldsymbol{\gamma}$ & \textbf{Calibration} & \textbf{Regular} $\boldsymbol{81^2}$ & \textbf{Reg.} $\boldsymbol{p_{15\text{--}81} / p_{36\text{--}81}}$ & \textbf{Jittered} $\boldsymbol{81^2}$ & \textbf{Jit.} $\boldsymbol{p_{15\text{--}81} / p_{36\text{--}81}}$ & $\boldsymbol{\mathrm{cond}_{81^2}}$ \textbf{(Reg. / Jit.)} \\
            \midrule
            \multicolumn{9}{l}{\textit{Axis A: Lowering $\beta$ under Quadratic Stabilization ($c_2 = 70$)}} \\
            Under-penalized            & 0.0                  & 2.0                   & $c_2 = 70$           & $6.72\%$                             & $0.38 / -1.41$                                                    & $11.57\%$                              & $-0.02 / -0.64$                                                   & $1.13\times 10^6$ / $8.06\times 10^6$                      \\
            Weakly penalized           & 0.5                  & 2.0                   & $c_2 = 70$           & $4.69\%$                             & $0.66 / -1.60$                                                    & $7.51\%$                               & $0.23 / -0.49$                                                    & $2.16\times 10^6$ / $1.52\times 10^7$                      \\
            Sublinear $\beta$          & 1.0                  & 2.0                   & $c_2 = 70$           & $3.00\%$                             & $1.03 / -0.89$                                                    & $4.87\%$                               & $0.53 / -0.48$                                                    & $4.19\times 10^6$ / $2.90\times 10^7$                      \\
            Original quadratic         & 1.5                  & 2.0                   & $c_2 = 70$           & $1.52\%$                             & $1.46 / 0.44$                                                     & $2.82\%$                               & $0.94 / -0.36$                                                    & $8.17\times 10^6$ / $5.59\times 10^7$                      \\
            \midrule
            \multicolumn{9}{l}{\textit{Axis B: Raising $\beta$ above $1.5$ under Linear Stabilization ($\alpha_0 = 0.10$, i.e., $c_{\mathrm{stab}} = 2.0$ at $\gamma = 1$)}} \\
            Baseline $\beta = 1.5$     & 1.5                  & 1.0                   & $\alpha_0 = 0.10$    & $0.89\%$                             & $1.33 / 1.22$                                                     & $1.99\%$                               & $0.62 / 0.11$                                                     & $8.04\times 10^6$ / $5.50\times 10^7$                      \\
            Higher $\beta = 2.0$       & 2.0                  & 1.0                   & $\alpha_0 = 0.10$    & $1.04\%$                             & $1.16 / 1.58$                                                     & $0.99\%$                               & $0.94 / 0.82$                                                     & $1.58\times 10^7$ / $1.07\times 10^8$                      \\
            Over-penalized             & 2.5                  & 1.0                   & $\alpha_0 = 0.10$    & $2.09\%$                             & $0.61 / 1.05$                                                     & $1.88\%$                               & $0.51 / 0.39$                                                     & $3.10\times 10^7$ / $2.09\times 10^8$                      \\
            \midrule
            \multicolumn{9}{l}{\textit{Axis C: Stabilization Exponent $\gamma$ at Fixed $\beta = 1.5$}} \\
            Constant (biased)          & 1.5                  & 0.0                   & $\alpha_0 = 0.20$    & $8.06\%$                             & $0.16 / 0.24$                                                     & $7.31\%$                               & $-0.13 / -0.06$                                                   & $6.77\times 10^6$ / $4.58\times 10^7$                      \\
            Constant (small)           & 1.5                  & 0.0                   & $\alpha_0 = 0.05$    & $0.78\%$                             & $1.10 / 1.45$                                                     & $0.79\%$                               & $1.36 / 1.37$                                                     & $7.83\times 10^6$ / $5.34\times 10^7$                      \\
            Sublinear $\gamma = 0.5$   & 1.5                  & 0.5                   & $\alpha_0 = 0.10$    & $0.78\%$                             & $1.35 / 1.88$                                                     & $0.79\%$                               & $1.17 / 1.16$                                                     & $7.83\times 10^6$ / $5.34\times 10^7$                      \\
            \textbf{Linear Baseline}   & \textbf{1.5}         & \textbf{1.0}          & $\boldsymbol{\alpha_0 = 0.20}$ ($c_{\mathrm{stab}} = 4.0$) & $\mathbf{0.78\%}$ & $\mathbf{1.72} / \mathbf{2.48}$                & $\mathbf{0.79\%}$                      & $\mathbf{1.61} / \mathbf{1.62}$                                   & $\mathbf{7.83\times 10^6}$ / $\mathbf{5.34\times 10^7}$    \\
            Superlinear $\gamma = 1.5$ & 1.5                  & 1.5                   & $\alpha_0 = 0.20$    & $0.89\%$                             & $1.83 / 1.77$                                                     & $1.99\%$                               & $1.31 / 0.00$                                                     & $8.04\times 10^6$ / $5.50\times 10^7$                      \\
            \bottomrule
        \end{tabular}
    }
\end{table}

Tables~\ref{tab:cavity_convergence_stabilization}--\ref{tab:beta_gamma_fine_mesh} and Fig.~\ref{fig:mlpg_cavity_stabilization} resolve both questions raised by the extended refinement study:
\begin{enumerate}
    \item \textbf{Why Quadratic Scaling ($\gamma = 2$) Overshoots Past $49 \times 49$:} Under $\gamma = 2, c_2 = 70$, $\alpha_{\mathrm{stab}}$ drops rapidly from $0.357$ ($15^2$) to $0.030$ ($49^2$) and $0.011$ ($81^2$), while $s_{\mathrm{div}}$ grows to $120.00$ at $81^2$. Consequently, $\alpha_{\mathrm{stab}}/s_{\mathrm{div}} \propto h_{\mathrm{nom}}^{3.5}$ collapses on fine meshes, returning them toward the unstabilized softening regime: signed $TE_{10}$ flips from $+13.08\%$ ($15^2$) to $-2.75\%$ ($49^2$) and $-2.55\%$ ($81^2$) regularly ($-3.96\%$ under $25\%$ jitter), causing $\overline{\varepsilon}_{k_c}$ to rebound above $49^2$ (Table~\ref{tab:cavity_convergence_stabilization}).
    \item \textbf{Why $\beta = 1.5$ Remains Essential with $\mathbf{K}_{\mathrm{stab}}$ Active:} Lowering $\beta \in \{0.0, 0.5, 1.0\}$ improves $\mathrm{cond}(\mathbf{K}_{\mathrm{red}})$ at $81\times 81$ by up to $7.2\times$, but leaves the expanding gradient nullspace under-penalized ($3.00\%\text{--}6.72\%$ regular and $4.87\%\text{--}11.57\%$ jittered error at $81^2$, $p_{36\text{--}81} < 0$; Table~\ref{tab:beta_gamma_fine_mesh}). Conversely, raising $\beta = 2.0$ halves the jittered $81^2$ error of $\alpha_0 = 0.10$ at double the condition number ($1.07\times 10^8$), while $\beta = 2.5$ induces volumetric locking ($2.09\%$ regular, $\mathrm{cond} = 2.09\times 10^8$). Hence $\beta = 1.5$ provides the most favorable trade-off between divergence control and conditioning among the tested exponents.
    \item \textbf{Linear Stabilization ($\gamma = 1$, $c_{\mathrm{stab}} = 4.0$) Sustains Monotonic Decay to $0.78\%\text{--}0.79\%$:} Decreasing $\gamma$ to $1$ with $c_{\mathrm{stab}} = 4.0$ ($\alpha_{\mathrm{stab}} = 0.050$ at $81^2$) preserves harmonic rank control while decaying asymptotically as $\mathcal{O}(h_{\mathrm{nom}})$. Cutoff errors drop monotonically from $14.58\%$ ($15^2$) to $\mathbf{0.78\%}$ ($81^2$) regularly ($p_{15\text{--}81} = 1.72$, $p_{36\text{--}81} = 2.48$), and from $8.80\%$ to $\mathbf{0.79\%}$ under $25\%$ jitter ($p_{15\text{--}81} = 1.61$, $p_{36\text{--}81} = 1.62$). In contrast, $\gamma = 0$ either injects an $\mathcal{O}(1)$ bias ($8.06\%$ at $\alpha_0 = 0.20$) or sacrifices asymptotic decay of the stabilization bias (Term~(IIIb) in Eq.~\eqref{eq:spectral_error_decomposition}, Appendix~\ref{subsubsec:strang_framework}).
    \item \textbf{Fine-Mesh Behavior ($0.78\%\text{--}0.79\%$ at $81\times 81$):} Because $(h_{\mathrm{nom}}/L)/\hat{h}_0 = 1/4$ at $81\times 81$, the Axis~C triples $(\gamma, \alpha_0) \in \{(0, 0.05), (0.5, 0.10), (1.0, 0.20)\}$ evaluate to the identical terminal weight $\alpha_{\mathrm{stab}}(81^2) = 0.050$ (and $\{(1.0, 0.10), (1.5, 0.20)\}$ to $0.025$), explaining their matching $81^2$ errors in Table~\ref{tab:beta_gamma_fine_mesh}. Among these trajectories, $\gamma = 1$ ($c_{\mathrm{stab}} = 4.0$) achieves the cleanest refinement progression. Under $25\%$ jitter, the $0.79\%$ level at $81^2$ (local slope $p_{65\to 81} = 0.31$) aligns directly with the continuous multiplicity floor induced by $94.4\%$ coverage (Section~\ref{subsubsec:multiplicity_rayleigh_effect}). On regular grids, where the multiplicity shift vanishes to $0.024\%$, the error continues descending from $1.03\%$ ($65^2$) to $0.78\%$ ($81^2$, local slope $p_{65\to 81} = 1.21$), consistent with the formal $\mathcal{O}(h_{\mathrm{nom}})$ stabilization term paired with the unpaired internal-arc fluxes $\mathcal{R}_h^{\mathrm{eig}}$ in Eq.~\eqref{eq:spectral_error_decomposition}.
    \item \textbf{Conditioning Growth and Sparse Shift-and-Invert Solvability:} Under the baseline linear stabilization ($\gamma = 1$, $c_{\mathrm{stab}} = 4.0$), $\mathrm{cond}(\mathbf{K}_{\mathrm{red}})$ increases from $1.00 \times 10^4$ ($15\times 15$) through $1.22 \times 10^6$ ($49\times 49$) to $7.83 \times 10^6$ ($81\times 81$) on regular grids (empirical slope $p_{\mathrm{cond}} = -4.03$). Under $25\%$ jitter, it rises similarly from $2.37 \times 10^4$ ($15\times 15$) through $3.91 \times 10^6$ ($49\times 49$) to $5.34 \times 10^7$ ($81\times 81$, $p_{\mathrm{cond}} = -4.61$), nearly matching the unstabilized $c_{\mathrm{stab}} = 0$ growth of $8.27 \times 10^6$ and $5.67 \times 10^7$ at $81 \times 81$ (Table~\ref{tab:beta_gamma_fine_mesh}). The high-frequency penalty modes are mapped to small-magnitude eigenvalues by the sparse Shift-and-Invert transformation $(\mathbf{K}_{\mathrm{red}} - \sigma_{\mathrm{shift}} \mathbf{M}_{\mathrm{red}})^{-1}\mathbf{M}_{\mathrm{red}}$ (with spectral shift $\sigma_{\mathrm{shift}} < \lambda_{\min}$) \cite{saad2011numerical, lehoucq1998arpack} and are therefore naturally excluded from the target low-frequency spectrum by the Krylov subspace. All eight-grid eigensolutions in Tables~\ref{tab:cavity_convergence_stabilization}--\ref{tab:cavity_81x81_modes} are computed with this sparse solver and cross-validated against dense SciPy \texttt{eigh} \cite{virtanen2020scipy} on $22\times 22$ to $< 1.9 \times 10^{-13}$ cutoff discrepancy. In 2D, sparse direct LU factorizations handle these shifted systems cleanly ($\approx 0.032\text{ s}$ on $49\times 49$ and $\approx 0.126\text{ s}$ on $81\times 81$), while large-scale 3D iterative Krylov solvers would require dedicated preconditioning (Section~\ref{subsec:deficiencies_challenges}).
\end{enumerate}

Accordingly, linear Orthogonal Energy Stabilization ($\gamma = 1$, $c_{\mathrm{stab}} = 4.0$) paired with $\beta = 1.5$ is adopted as the calibrated VNMM-MLPG baseline for Stage~6. Table~\ref{tab:cavity_81x81_modes} verifies all 10 physical modes on the finest $81\times 81$ cloud: every cutoff error remains below $1.93\%$ ($0\%$ jitter) and $1.69\%$ ($25\%$ jitter), while $\max R_{\mathrm{div}} \le 0.06\%$ demonstrates nearly $60\times$ tighter discrete solenoidality than on the unstabilized $21\times 21$ mesh of Table~\ref{tab:cavity_eigenvalues}.

\begin{table}[!t]
    \centering
    \small
    \caption{Full 10-mode physical cavity spectrum on the finest $81 \times 81$ cloud ($6,241$ active DoFs, $h_{\mathrm{nom}} = 0.0393\text{ m}$) for the calibrated linear baseline ($\beta = 1.5$, $\gamma = 1$, $c_{\mathrm{stab}} = 4.0$, giving $s_{\mathrm{div}} = 120.00$ and $\alpha_{\mathrm{stab}} = 0.050$, $\alpha_s = 0.80$, $\text{seed}=42$) under $0\%$ and $25\%$ jitter.}
    \label{tab:cavity_81x81_modes}
    \resizebox{\textwidth}{!}{
        \begin{tabular}{cccccccc}
            \toprule
            \textbf{Mode} & $\boldsymbol{k_{c,\mathrm{ref}}}$ \textbf{[m$^{-1}$]} & \multicolumn{3}{c}{\textbf{Regular Lattice ($0\%$ Jitter)}} & \multicolumn{3}{c}{\textbf{Perturbed Cloud ($25\%$ Jitter)}} \\
            \cmidrule(lr){3-5} \cmidrule(lr){6-8}
                          &                                                        & $\boldsymbol{k_{c,\mathrm{num}}}$ & \textbf{Signed Err.} & $\boldsymbol{R_{\mathrm{div}}}$ & $\boldsymbol{k_{c,\mathrm{num}}}$ & \textbf{Signed Err.} & $\boldsymbol{R_{\mathrm{div}}}$ \\
            \midrule
            $TE_{10}$     & 1.0000                                                 & 0.9955                            & $-0.45\%$            & $0.03\%$                        & 0.9831                            & $-1.69\%$            & $0.04\%$                        \\
            $TE_{01}$     & 1.0000                                                 & 0.9977                            & $-0.23\%$            & $0.03\%$                        & 0.9872                            & $-1.28\%$            & $0.04\%$                        \\
            $TE_{11}$     & 1.4142                                                 & 1.4091                            & $-0.36\%$            & $0.03\%$                        & 1.3998                            & $-1.02\%$            & $0.04\%$                        \\
            $TE_{20}$     & 2.0000                                                 & 2.0091                            & $+0.45\%$            & $0.03\%$                        & 1.9673                            & $-1.63\%$            & $0.05\%$                        \\
            $TE_{02}$     & 2.0000                                                 & 2.0147                            & $+0.73\%$            & $0.04\%$                        & 1.9892                            & $-0.54\%$            & $0.05\%$                        \\
            $TE_{21}$     & 2.2361                                                 & 2.2454                            & $+0.42\%$            & $0.04\%$                        & 2.2239                            & $-0.54\%$            & $0.05\%$                        \\
            $TE_{12}$     & 2.2361                                                 & 2.2504                            & $+0.64\%$            & $0.04\%$                        & 2.2367                            & $+0.03\%$            & $0.05\%$                        \\
            $TE_{22}$     & 2.8284                                                 & 2.8627                            & $+1.21\%$            & $0.04\%$                        & 2.8316                            & $+0.11\%$            & $0.06\%$                        \\
            $TE_{30}$     & 3.0000                                                 & 3.0422                            & $+1.41\%$            & $0.04\%$                        & 3.0009                            & $+0.03\%$            & $0.06\%$                        \\
            $TE_{03}$     & 3.0000                                                 & 3.0580                            & $+1.93\%$            & $0.05\%$                        & 3.0315                            & $+1.05\%$            & $0.06\%$                        \\
            \midrule
            \multicolumn{2}{l}{\textbf{Mean Abs.\ Error (10 Modes):}} & --- & $\mathbf{0.78\%}$ & --- & --- & $\mathbf{0.79\%}$ & --- \\
            \multicolumn{2}{l}{\textbf{Max Abs.\ Error / Max $R_{\mathrm{div}}$:}} & --- & $1.93\%$ & $0.05\%$ & --- & $1.69\%$ & $0.06\%$ \\
            \multicolumn{2}{l}{\textbf{Spurious Modes in Band ($\le 10$):}} & --- & \textbf{0} & --- & --- & \textbf{0} & --- \\
            \bottomrule
        \end{tabular}
    }
\end{table}

\subsubsection{Spectral Effect of Spatially Varying Multiplicity \texorpdfstring{$M(\mathbf{x})$}{M(x)}}
\label{subsubsec:multiplicity_rayleigh_effect}

To isolate the geometric shift introduced by $\widetilde{M}(\mathbf{x})$ in \eqref{eq:multiplicity_rayleigh_quotient} independently of $\mathcal{P}^1$ interpolation errors, we evaluate $k_{c,nm}^{(M)} = \sqrt{\lambda^{(M)}(\vec{E}_{nm}^*)}$ on the first 10 analytical $TE_{nm}$ modes over the same $\alpha_s = 0.80$ disk covers across all eight grids ($15\times 15$ to $81\times 81$, $\text{seed}=42$). On regular grids ($0\%$ jitter), the mean relative shift $\overline{|k_{c,nm}^{(M)} - k_{c,nm}|/k_{c,nm}}$ decays monotonically from $0.635\%$ ($15^2$) $\to 0.319\%$ ($22^2$) $\to 0.187\%$ ($29^2$) $\to 0.122\%$ ($36^2$) $\to 0.086\%$ ($43^2$) $\to 0.066\%$ ($49^2$) $\to 0.037\%$ ($65^2$) to $0.024\%$ ($81^2$) at empirical rate $p = 1.89$ (matching Appendix~\ref{subsec:conformity_discussion}). Under $25\%$ jitter, irregular $M(\mathbf{x})$ fluctuations and the $5.6\%$ unintegrated voids yield mean shifts of $2.03\%$ ($15^2$), $1.49\%$ ($22^2$), $1.43\%$ ($29^2$), $1.21\%$ ($36^2$), $0.96\%$ ($43^2$), $1.33\%$ ($49^2$), $0.85\%$ ($65^2$), and $0.63\%$ ($81^2$), corresponding to an empirical rate of $p = 0.59$ (peak single-mode shifts $1.12\%\text{--}3.92\%$). This behavior quantitatively explains the $\sim 0.6\%\text{--}0.8\%$ fine-mesh floor observed on perturbed clouds in Table~\ref{tab:cavity_convergence_stabilization}.

\subsection{Stage 6: Computational Benchmark: VNMM-MLPG vs. VNMM-Gauss}
\label{subsec:stage6_benchmark}

We compare calibrated cell-free VNMM-MLPG (scaled $s_{\mathrm{div}}$ \eqref{eq:s_div_scaling_global} with $\beta = 1.5$, linear $\alpha_{\mathrm{stab}}$ \eqref{eq:alpha_stab_vms_scaling} with $\gamma = 1$, $c_{\mathrm{stab}} = 4.0$) against calibrated cell-based VNMM-Gauss ($3 \times 3$ cell quadrature, fixed $s_{\mathrm{div}} = 6.0$ \cite{stable_vnmm2d_2026}; Section~\ref{subsubsec:constraint_ratio_scaling}). Both methods are evaluated across the six-grid benchmark sequence inherited from \cite{stable_vnmm2d_2026} ($15 \times 15$ to $49 \times 49$ under $0\%$ and $25\%$ jitter, $\text{seed}=42$, single-thread AMD Ryzen / Linux). VNMM-Gauss stops at $49\times 49$ because its per-Gauss-point stencil lookups scale steeply with cell count, whereas cell-free VNMM-MLPG is further extended to $65 \times 65$ and $81 \times 81$.

Table~\ref{tab:benchmark_comparison} reports assembly runtime and stiffness sparsity (non-zero entries, NNZ) on the six matched grids, while Table~\ref{tab:jitter_benchmark_comparison} and Fig.~\ref{fig:benchmark} compare cutoff errors and empirical rates $p$. Here, \emph{Assembly Time} covers the full matrix pipeline from raw coordinates to CSR PEC-condensed operators---$k$-d tree setup ($< 0.4\text{ ms}$), stencil selection, local inversions $\boldsymbol{\beta}^I = (\mathbf{A}_I)^{-1}$ ($\sim 60\%\text{--}65\%$ of runtime), and closed-form matrix assembly ($\sim 35\%\text{--}40\%$)---excluding the sparse Shift-and-Invert Arnoldi solve (Section~\ref{subsubsec:resolution_stagnation_stabilization}).

\begin{table}[!t]
    \centering
    \small
    \caption{Computational benchmark on the 2D PEC resonant cavity across the six grid resolutions matching \cite{stable_vnmm2d_2026}: assembly runtime (scoped in the text, excluding eigensolver solves) and global stiffness non-zero entries (NNZ) for calibrated VNMM-MLPG ($\beta = 1.5$, $\gamma = 1$, $c_{\mathrm{stab}} = 4.0$) vs. VNMM-Gauss ($3 \times 3$ cell quadrature, $s_{\mathrm{div}} = 6.0$). On finer grids, VNMM-MLPG reaches $64,301$ NNZ ($3,969$ DoFs) at $65\times 65$ and $102,015$ NNZ ($6,241$ DoFs) at $81\times 81$.}
    \label{tab:benchmark_comparison}
    \resizebox{0.88\textwidth}{!}{
        \begin{tabular}{ccccccc}
            \toprule
            \textbf{Grid}                 & \textbf{Active} & \multicolumn{2}{c}{\textbf{Assembly Time [s]}} & \textbf{Assembly} & \multicolumn{2}{c}{\textbf{Stiffness NNZ}}                                                      \\
            \cmidrule(lr){3-4} \cmidrule(lr){6-7}
            $\boldsymbol{N_x \times N_y}$ & \textbf{DoFs}   & \textbf{Gauss}                                 & \textbf{MLPG}     & \textbf{Speedup}                           & \textbf{Gauss} & \textbf{MLPG}            \\
            \midrule
            $15 \times 15$            & 169             & 0.1239 s                                       & \textbf{0.0624 s} & \textbf{1.99}$\times$                      & 3,091          & \textbf{2,467} ($-20\%$) \\
            $22 \times 22$            & 400             & 0.2670 s                                       & \textbf{0.1060 s} & \textbf{2.52}$\times$                      & 7,802          & \textbf{6,140} ($-21\%$) \\
            $29 \times 29$            & 729             & 0.4529 s                                       & \textbf{0.1593 s} & \textbf{2.84}$\times$                      & 14,677         & \textbf{11,509} ($-22\%$) \\
            $36 \times 36$            & 1156            & 0.6813 s                                       & \textbf{0.2263 s} & \textbf{3.01}$\times$                      & 23,634         & \textbf{18,460} ($-22\%$) \\
            $43 \times 43$            & 1681            & 0.9754 s                                       & \textbf{0.2855 s} & \textbf{3.42}$\times$                      & 35,409         & \textbf{26,813} ($-24\%$) \\
            $49 \times 49$            & 2209            & 1.2750 s                                       & \textbf{0.3459 s} & \textbf{3.69}$\times$                      & 44,939         & \textbf{35,385} ($-21\%$) \\
            \bottomrule
        \end{tabular}
    }
\end{table}

\begin{figure}[!t]
    \centering
    \includegraphics[width=0.68\textwidth]{fig_mlpg_benchmark_comparison.pdf}
    \caption{Log--log cutoff convergence corresponding to Table~\ref{tab:jitter_benchmark_comparison}, comparing VNMM-Gauss (dashed, $15^2\text{--}49^2$) against calibrated VNMM-MLPG (solid, $15^2\text{--}81^2$) under $0\%$ (blue) and $25\%$ (green) nodal jitter.}
    \label{fig:benchmark}
\end{figure}

\begin{table}[!t]
    \centering
    \small
    \caption{Cavity eigenvalue cutoff wavenumber convergence under nodal perturbations ($0\%$ and $25\%$ jitter): VNMM-Gauss ($3 \times 3$ cell quadrature, $s_{\mathrm{div}} = 6.0$ across $15\times 15$ to $49\times 49$ \cite{stable_vnmm2d_2026}) vs. calibrated VNMM-MLPG ($\beta = 1.5$, $\gamma = 1$, $c_{\mathrm{stab}} = 4.0$, extended to $65 \times 65$ and $81 \times 81$). Bold entries highlight the lower error on matched resolutions ($15^2\text{--}49^2$); empirical rates $p$ are computed via log-log linear regression.}
    \label{tab:jitter_benchmark_comparison}
    \resizebox{0.85\textwidth}{!}{
        \begin{tabular}{cccccc}
            \toprule
            \textbf{Grid}                                                                                     & \textbf{Nominal}                          & \multicolumn{2}{c}{\textbf{0\% Jitter (Regular)}} & \multicolumn{2}{c}{\textbf{25\% Jitter}}                                           \\
            \cmidrule(lr){3-4} \cmidrule(lr){5-6}
            $\boldsymbol{N_x \times N_y}$                                                                     & $\boldsymbol{h_{\mathrm{nom}}}$ \textbf{[m]} & \textbf{Gauss}                                    & \textbf{MLPG}                            & \textbf{Gauss}    & \textbf{MLPG}      \\
            \midrule
            $15 \times 15$                                                                                    & 0.2244                                    & $25.25\%$                                         & $\mathbf{14.58\%}$                       & $24.67\%$         & $\mathbf{8.80\%}$  \\
            $22 \times 22$                                                                                    & 0.1496                                    & $13.46\%$                                         & $\mathbf{8.35\%}$                        & $11.49\%$         & $\mathbf{8.27\%}$  \\
            $29 \times 29$                                                                                    & 0.1122                                    & $\mathbf{6.18\%}$                                 & $6.19\%$                                 & $5.94\%$          & $\mathbf{5.28\%}$  \\
            $36 \times 36$                                                                                    & 0.0898                                    & $\mathbf{4.22\%}$                                 & $5.64\%$                                 & $3.10\%$          & $\mathbf{2.73\%}$  \\
            $43 \times 43$                                                                                    & 0.0748                                    & $\mathbf{2.46\%}$                                 & $3.50\%$                                 & $\mathbf{1.96\%}$ & $2.13\%$           \\
            $49 \times 49$                                                                                    & 0.0654                                    & $\mathbf{1.53\%}$                                 & $2.39\%$                                 & $1.63\%$          & $\mathbf{1.46\%}$  \\
            $65 \times 65$                                                                                    & 0.0491                                    & ---                                               & $1.03\%$                                 & ---               & $0.85\%$           \\
            $81 \times 81$                                                                                    & 0.0393                                    & ---                                               & $0.78\%$                                 & ---               & $0.79\%$           \\
            \midrule
            \multicolumn{2}{c}{\textbf{Rate} $\boldsymbol{p_{15\text{--}49}}$ ($\boldsymbol{\mathcal{O}(h_{\mathrm{nom}}^p)}$):} & $\mathbf{2.25}$                  & $1.35$                                            & $\mathbf{2.29}$   & $1.54$             \\
            \multicolumn{2}{c}{\textbf{MLPG Fine Rate} $\boldsymbol{p_{15\text{--}81}}$:}                                         & ---                              & $\mathbf{1.72}$                                   & ---               & $\mathbf{1.61}$    \\
            \bottomrule
        \end{tabular}
    }
\end{table}

Tables~\ref{tab:benchmark_comparison}--\ref{tab:jitter_benchmark_comparison} and Fig.~\ref{fig:benchmark} highlight four comparative features:
\begin{enumerate}
    \item \textbf{Assembly Speedup (2.0$\times$--3.7$\times$) and End-to-End Runtime (3.35$\times$--3.94$\times$):} Evaluating closed-form formulas \eqref{eq:k_curl_analytic}--\eqref{eq:m_local_analytic} over $N$ node-centered disks replaces thousands of Gauss-point queries, cutting assembly time from $0.124\text{--}1.275\text{ s}$ (Gauss) to $0.062\text{--}0.346\text{ s}$ ($15^2$ to $49^2$). Including the sparse Shift-and-Invert Arnoldi eigensolve ($k = 25, \sigma_{\mathrm{shift}} = -0.05$), the $49\times 49$ solve stage takes $\approx 0.032\text{ s}$ for VNMM-MLPG vs.\ $\approx 0.038\text{ s}$ for VNMM-Gauss ($0\%$ jitter; $0.031\text{ s}$ vs.\ $0.039\text{ s}$ under $25\%$ jitter). MLPG's $21\%$ lower NNZ offsets its higher condition number ($1.22\times 10^6$ vs.\ $4.95\times 10^5$ regularly; $3.91\times 10^6$ vs.\ $4.17\times 10^5$ at $25\%$ jitter), preserving an end-to-end (assembly plus eigensolve) speedup of $3.35\times$ ($0\%$ jitter, $0.393\text{ s}$ vs.\ $1.316\text{ s}$) to $3.94\times$ ($25\%$ jitter, $0.321\text{ s}$ vs.\ $1.267\text{ s}$).
    \item \textbf{$20\%\text{--}24\%$ Sparser Global Operators:} Because each disk $D_I$ couples only the 6 nodes of $\Omega_u^I$ rather than the merged stencils of 9 Gauss points per cell, stiffness NNZ drops by $20\%\text{--}24\%$ (e.g., $35,385$ vs.\ $44,939$ on $49\times 49$), reducing memory footprint and improving cache locality.
    \item \textbf{Regular-Grid Accuracy and Fine-Mesh Continuation ($0\%$ Jitter):} Under linear stabilization ($\gamma = 1, c_{\mathrm{stab}} = 4.0$), VNMM-MLPG achieves substantially lower coarse-grid errors than VNMM-Gauss ($14.58\%$ vs.\ $25.25\%$ at $15^2$, $8.35\%$ vs.\ $13.46\%$ at $22^2$), while incurring slightly higher intermediate residuals ($2.39\%$ vs.\ $1.53\%$ at $49^2$; note $c_{\mathrm{stab}} = 2.0$ achieves $1.16\%$ at $49^2$). On finer clouds, however, it continues monotonically past $49\times 49$ down to $1.03\%$ ($65^2$) and $\mathbf{0.78\%}$ ($81^2$), giving an overall rate $p_{15\text{--}81} = 1.72$ and accelerating to $p_{36\text{--}81} = 2.48$.
    \item \textbf{Perturbed-Cloud Superiority ($25\%$ Jitter):} Even with $94.4\%$ subdomain coverage, linearly stabilized VNMM-MLPG outperforms VNMM-Gauss on five of the six matched grids ($15^2$ to $36^2$ and $49^2$), achieving $2.80\times$ lower error at $15^2$ (where intra-cell Gauss stencil switching degrades accuracy) and $1.46\%$ vs.\ $1.63\%$ at $49^2$. Beyond $49\times 49$, the error continues descending cleanly to $0.85\%$ ($65^2$) and $0.79\%$ ($81^2$) at $p_{15\text{--}81} = 1.61$, reaching the multiplicity floor of Section~\ref{subsubsec:multiplicity_rayleigh_effect}.
\end{enumerate}

%% =========================================================================
%% SECTION 4: DISCUSSION
%% =========================================================================
\section{Discussion: Advantages, Deficiencies, and Open Questions}
\label{sec:discussion}

\subsection{Synthesis of Advantages}
\label{subsec:advantages_synthesis}

Table~\ref{tab:method_comparison} summarizes the structural contrast among conforming edge FEM (N\'{e}d\'{e}lec $\mathcal{ND}_1$), cell-based VNMM-Gauss \cite{stable_vnmm2d_2026}, and cell-free VNMM-MLPG. Beyond the benchmark summary in Table~\ref{tab:method_comparison}, the proposed formulation offers five key methodological and computational advantages:
\begin{enumerate}
    \item \textbf{True Cell-Free Autonomy and Geometric Robustness:} Replacing conforming triangulations, Voronoi dual diagrams, and Cartesian background cells with autonomous node-centered circular disks eliminates mesh generation, background-cell trimming, and intra-cell stencil discontinuities. Coupled with the scale-invariant determinant stencil qualification \eqref{eq:scale_invariant_det_criterion}, each node independently builds a well-conditioned $\mathcal{P}^1$ support cloud. This autonomy translates directly into robustness under disorder: on $25\%$-jittered point clouds, VNMM-MLPG achieves lower cutoff errors than cell-based VNMM-Gauss on five of the six matched grids (up to $2.80\times$ lower at $15\times 15$; Table~\ref{tab:jitter_benchmark_comparison}).
    \item \textbf{Exact Closed-Form Analytical Integration of the Local Weak Form:} Because $\mathcal{P}^1$ curl, divergence, and orthogonal stabilization tensors are spatially constant over each disk while mass integrands are quadratic, all $6\times 6$ local stiffness and mass blocks evaluate in exact closed form \eqref{eq:k_curl_analytic}--\eqref{eq:m_local_analytic}. Near rectilinear boundaries and convex corners, exact cap-area and geometric-moment formulas \eqref{eq:k_truncated_scaled}--\eqref{eq:m_truncated_moments} preserve machine-precision integration without numerical quadrature, boundary-clipping errors, or pairwise lens-intersection cubatures (Table~\ref{tab:stage1_element_verification}).
    \item \textbf{Direct, Penalty-Free Boundary Condition Enforcement:} Unlike moving least-squares (MLS) meshfree approximations that lose the Kronecker-delta property and require Lagrange multipliers, Nitsche terms, or penalty weights for essential boundaries, directional $\mathcal{P}^1$ VNMM interpolation satisfies $e_I = \vec{E}(\mathbf{x}_I)\cdot\vec{t}_I$ directly at every boundary node. Essential Dirichlet (PEC) conditions ($E_t = 0$) are therefore enforced exactly by standard row/column elimination (Section~\ref{subsubsec:dirichlet_bc}) while preserving operator symmetry and sparsity. In turn, natural Neumann (PMC) conditions enter cleanly through 1D boundary-segment integrals that are exact with 2-point Gauss--Legendre quadrature for affine fluxes (Table~\ref{tab:stage1_element_verification}(b)).
    \item \textbf{Guaranteed Algebraic Symmetry, Real Spectra, and Operator Efficiency:} The additive Symmetric Subdomain Galerkin (SSG) assembly \eqref{eq:global_assembly} yields a symmetric positive-semidefinite stiffness $\mathbf{K}$ and a strictly symmetric positive-definite mass $\mathbf{M}$ by construction. This guarantees strictly real eigenvalues and avoids the complex-conjugate spectral pollution and refinement instabilities observed in classical unsymmetric row-by-row MLPG/LPIM assemblies (Section~\ref{subsubsec:mms_assembly_comparison}). Simultaneously, evaluating one closed-form disk per node rather than $3\times 3$ Gauss points per background cell reduces global stiffness NNZ by $20\%\text{--}24\%$ and delivers 2.0$\times$--3.7$\times$ faster assembly (3.35$\times$--3.94$\times$ end-to-end speedup at $49\times 49$; Section~\ref{subsec:stage6_benchmark}).
    \item \textbf{Spurious-Mode Shifting, Spectral Fidelity, and Clean Field Reconstruction:} Moving from the $9:1$ over-constraint ratio of $3\times 3$ cell quadrature to the leaner $1.1:1$ single-disk constraint ratio makes pairing superlinear divergence scaling with VMS-inspired orthogonal energy stabilization essential to restore local rank over the harmonic gradient subspace. Once paired, the stabilized cell-free operator keeps spurious gradient modes shifted outside the 10-mode operational band, sustains monotonic fine-mesh cutoff convergence down to $0.78\%\text{--}0.79\%$ at $81\times 81$ (Section~\ref{subsubsec:resolution_stagnation_stabilization}), and produces faithful, artifact-free global modal field reconstructions $\vec{E}_{\mathrm{global}}$ (Fig.~\ref{fig:modal_fields_21x21}).
\end{enumerate}

\begin{table}[!t]
    \centering
    \small
    \caption{Comparative synthesis of methodological and benchmark features among CEM formulations (sparsity and assembly speedups of VNMM-MLPG are benchmarked relative to calibrated VNMM-Gauss).}
    \label{tab:method_comparison}
    \resizebox{\textwidth}{!}{
        \begin{tabular}{lccc}
            \toprule
            \textbf{Methodological Feature}      & \textbf{Edge FEM (N\'{e}d\'{e}lec)} & \textbf{VNMM-Gauss (Previous)} & \textbf{VNMM-MLPG (Present)}                 \\
            \midrule
            Mesh Dependency                      & Conforming Triangulation    & Cartesian Background Cells     & \textbf{Truly Cell-Free}                     \\
            Integration Scheme                   & Numerical Simplex Gauss     & $3 \times 3$ Cell Quadrature   & \textbf{Closed-Form Analytic}                \\
            Intra-Cell Stencil Switching         & None                        & Present (requires dense Gauss) & \textbf{None (Autonomous Disks)}             \\
            Matrix Symmetry                      & Symmetric                   & Symmetric                      & \textbf{Symmetric}                           \\
            Matrix Sparsity (Stiffness NNZ)      & Not benchmarked ($\approx 5$ NNZ/row) & Benchmark baseline   & \textbf{20\%--24\% Sparser vs.\ Gauss}       \\
            Assembly Time Speedup                & Not benchmarked             & Benchmark baseline             & \textbf{2.0$\times$--3.7$\times$ Faster vs.\ Gauss} \\
            Patch Test Satisfaction              & Machine Precision           & Machine Precision              & \textbf{Machine Precision ($\sim 10^{-13}\text{--}10^{-15}$, local eq.)} \\
            Spurious Mode Treatment              & Discrete de Rham complex    & Div-Curl Regularization        & \textbf{Superlinear Div-Curl Scaling}        \\
            Local Operator Rank Stability        & Conforming exact sequence   & Dense $3 \times 3$ cell quadrature & \textbf{Orthogonal Energy Stabilization} \\
            \bottomrule
        \end{tabular}
    }
\end{table}

\subsection{Current Deficiencies and Open Challenges}
\label{subsec:deficiencies_challenges}

Several theoretical and practical limitations remain for future investigation:
\begin{enumerate}
    \item \textbf{Open Spectral Convergence Proof and Scope Restriction:} Because $\mathcal{V}_h \not\subset H(\mathbf{curl};\Omega)$, the consistency residual $\mathcal{R}_h$ \eqref{eq:global_consistency_residual} involves unpaired fluxes across internal arcs $L_{sI}$. Establishing uniform discrete coercivity, discrete compactness \cite{boffi2006finite, boffi2013mixed}, and asymptotic residual decay for general point clouds (Appendix~\ref{subsubsec:strang_framework}) remains open. Likewise, extending single-disk assembly to general source-driven BVPs without analytical $L_{sI}$ fluxes would require inter-subdomain coupling (e.g., partition-of-unity weightings or jump penalties).
    \item \textbf{Partial Coverage, Multiplicity Shift, and Solenoidal Boundary Assumption:} Under $25\%$ jitter with fixed $\alpha_s = 0.80$ and nearest-neighbor $h_I$, coverage drops to $94.4\%$, restricting norm equivalence to $\Omega_{\mathrm{cov}}$, breaking periodic $\mathcal{O}(h_{\mathrm{nom}}^2)$ cancellation, and leaving a slowly decaying multiplicity floor (from $1.33\%$ on $49\times 49$ to $0.63\%$ on $81\times 81$, Section~\ref{subsubsec:multiplicity_rayleigh_effect}). Adaptive $k$-NN radii restore full coverage. In addition, omitting the normal traction $\tau_n(\vec{E})$ on $\Gamma_N$ (Section~\ref{subsubsec:neumann_divergence_treatment}) assumes exact solenoidality ($\nabla \cdot \vec{E} = 0$).
    \item \textbf{Higher-Order Dispersion, Conditioning, and Empirical Calibration:} Higher-order modes exhibit larger divergence ratios due to linear $\mathcal{P}^1$ capacity (rising with mode index in Tables~\ref{tab:cavity_eigenvalues} and~\ref{tab:cavity_81x81_modes}, though suppressed to $\le 0.06\%$ at $81\times 81$ under scaled $s_{\mathrm{div}}$). Moreover, the calibrated scalings ($s_{\mathrm{div}} \propto (h_{\mathrm{char}}/L)^{-1.5}$, linear $\alpha_{\mathrm{stab}} \propto h_{\mathrm{char}}/L$) rely on empirical constants ($s_{\mathrm{div}}^{(0)} = 15.0$, $c_{\mathrm{stab}} = 4.0$) and introduce coarse-grid stiffening. They also raise $\mathrm{cond}(\mathbf{K}_{\mathrm{red}})$ with empirical slopes $p_{\mathrm{cond}} = -4.03$ (regular) to $-4.61$ ($25\%$ jitter; Section~\ref{subsubsec:resolution_stagnation_stabilization}), necessitating preconditioning for large-scale 3D iterative Krylov solvers. Exploring quadratic ($\mathcal{P}^2$) vector bases and parameter-free mixed formulations (item~4) will mitigate these limitations.
    \item \textbf{Curved Boundaries, Non-Convex Density, Material Interfaces, and 3D Extension:} The closed-form truncation formulas of Section~\ref{subsubsec:truncated_subdomains_analytic} and Appendix~\ref{app:truncated_geometry} assume rectilinear walls and orthogonal convex corners, and all spectral benchmarks herein are restricted to convex rectangular cavities where solutions lie in $\mathbf{H}^1(\Omega)$. Generalizing integration to curved CAD boundaries, multi-dielectric interfaces with normal jumps ($[E_n] \ne 0$, as in the Method of Finite Spheres \cite{nicomedes2021acoustic}), and non-convex domains with singular re-entrant corners ($\vec{E} \sim r^{\pi/\theta_{\mathrm{corner}} - 1} \in H(\mathbf{curl},\Omega) \setminus \mathbf{H}^1(\Omega)$) is critical. In particular, because $\mathbf{H}^1(\Omega)$ is not dense in $H_0(\mathbf{curl},\Omega) \cap H(\operatorname{div},\Omega)$ on non-convex geometries, unweighted divergence penalties risk Costabel--Dauge blocking onto smooth subspaces unless paired with singular-weight regularization \cite{costabel2002weighted} or a mixed saddle-point formulation with a scalar Lagrange multiplier. Validating $H(\mathbf{curl},\Omega)$ approximability on non-convex benchmarks (such as L-shaped cavities) and extending closed-form assembly to 3D spherical subdomains ($V_0 = \frac{4}{3}\pi r_s^3$, $I_2 = \frac{4}{15}\pi r_s^5$ with 12-term $\mathcal{P}_1^3$ bases) constitute primary directions for future work.
\end{enumerate}

%% =========================================================================
%% SECTION 5: CONCLUSIONS AND FUTURE WORK
%% =========================================================================
\section{Conclusions and Future Work}
\label{sec:conclusions}

This paper developed and validated VNMM-MLPG, a truly cell-free vector nodal meshless method for source-free 2D electromagnetic eigenvalue problems in $H(\mathbf{curl})$. Starting from directional $\mathcal{P}^1$ vector point interpolation on autonomous circular subdomains, we first showed via Stokes' and Gauss' theorems that boundary-vanishing weight functions (MLPG1) collapse the stiffness of constant-derivative $\mathcal{P}^1$ stencils to zero, motivating a Local Galerkin (MLPG6) weak form. Evaluating all local curl, divergence, stabilization, and mass integrals in exact closed form---including cap-area and moment formulas for boundary-truncated disks---eliminates background integration cells and lens-intersection cubatures alike. Combining these local blocks through an additive Symmetric Subdomain Galerkin (SSG) assembly produces a symmetric positive-semidefinite stiffness $\mathbf{K}$ and a strictly positive-definite mass $\mathbf{M}$, whose multiplicity-weighted Rayleigh quotient homogenizes at second order on regular lattices. In addition, directional Kronecker-delta collocation enables direct, penalty-free enforcement of PEC boundaries alongside natural 1D Neumann integrals.

A six-stage numerical verification established both the algebraic reliability and the spectral accuracy of the formulation. Element-level checks confirmed machine-precision agreement of the analytical interior, truncated-disk, and 1D Neumann operators, while affine patch tests and auxiliary local-equilibrium MMS studies verified linear completeness and confirmed that SSG preserves positive definiteness where unsymmetric row-by-row MLPG/LPIM loses stability under refinement. Subdomain sensitivity sweeps further identified $\alpha_s \in [0.70, 1.20]$ as the operational band (with baseline $\alpha_s = 0.80$) and isolated the geometric multiplicity floor induced by partial coverage on disordered clouds. On 2D PEC cavity benchmarks, diagnosing the rank-2 local stiffness kernel led to pairing a VMS-inspired linear orthogonal energy stabilization ($\alpha_{\mathrm{stab}} = c_{\mathrm{stab}}\,h_{\mathrm{char}}/L$, $\gamma = 1$, $c_{\mathrm{stab}} = 4.0$) with superlinear divergence penalization ($s_{\mathrm{div}} \propto (h_{\mathrm{char}}/L)^{-1.5}$, $\beta = 1.5$). Extended refinement studies up to $81\times 81$ showed that this linear scaling avoids both the fine-mesh rebound of quadratic decay ($\gamma = 2$) and the nullspace under-penalization of $\beta < 1.5$. It sustains monotonic cutoff convergence down to $0.78\%$ ($0\%$ jitter, $p_{15\text{--}81} = 1.72$) and $0.79\%$ ($25\%$ jitter, $p_{15\text{--}81} = 1.61$), with zero spurious modes in the 10-mode operational band and faithful modal field reconstructions. Across matched grids ($15^2\text{--}49^2$), VNMM-MLPG also delivered 2.0$\times$--3.7$\times$ faster assembly (3.35$\times$--3.94$\times$ end-to-end speedup at $49\times 49$), $20\%\text{--}24\%$ sparser stiffness operators, and lower cutoff errors on most perturbed clouds relative to calibrated cell-based VNMM-Gauss.

Building on the open questions analyzed in Section~\ref{subsec:deficiencies_challenges}, future research will proceed along three main directions. First, on the theoretical front, establishing uniform discrete coercivity, discrete compactness, and rigorous spectral convergence bounds for non-conforming SSG spaces---as well as inter-subdomain coupling mechanisms for general source-driven boundary-value problems---remains an important mathematical milestone. Second, to broaden geometric and physical applicability, the method will be extended to curved boundaries, adaptive $k$-NN subdomain radii that remove the irregular-cloud coverage floor, multi-dielectric interfaces with normal-component jumps, and non-convex domains with re-entrant corner singularities using singular-weight regularization or mixed saddle-point formulations. Finally, investigating quadratic ($\mathcal{P}^2$) vector bases, closed-form 3D spherical subdomains, and preconditioned iterative eigensolvers will pave the way toward high-accuracy, large-scale 3D cell-free electromagnetic cavity and waveguide modeling.

\appendix
\section{Closed-Form Geometry, Active Areas, and Moments for Truncated Subdomains}
\label{app:truncated_geometry}

This appendix provides analytical closed-form expressions for the active physical area $S_0^I = \iint_{D_I} d\Omega$, area fraction $\eta_I = S_0^I / S_0^{\mathrm{full}}$, and geometric moments ($S_{1x}^I, S_{1y}^I, S_{2xx}^I, S_{2yy}^I, S_{2xy}^I$) required by local truncated matrices in Section~\ref{subsubsec:truncated_subdomains_analytic}.

\subsection{Active Area and Area Fraction Formulas}
For a chord at normal distance $d \in [0, r_s]$ from the disk center, the circular cap area is:
\begin{equation}
    A_{\mathrm{cap}}(d, r_s) = r_s^2 \left[ \arccos\left(\frac{d}{r_s}\right) - \frac{d}{r_s} \sqrt{1 - \left(\frac{d}{r_s}\right)^2} \right].
    \label{eq:cap_area_formula}
\end{equation}
Canonical boundary configurations evaluate as follows:
\begin{itemize}
    \item \textbf{Flat Boundary Half-Disk (Case 3):} $\eta_B = 1/2$, $S_0^B = \frac{1}{2} \pi (r_s^B)^2$. At an orthogonal convex corner vertex, $\eta_B = 1/4$.
    \item \textbf{Boundary Node Near Corner (Case 3B):} At tangential distance $d_{\parallel} < r_s^B$ along the wall from a perpendicular corner wall:
          \begin{equation}
              \eta_B(d_{\parallel}) = \frac{1}{2} \left( 1 - \frac{A_{\mathrm{cap}}(d_{\parallel}, r_s^B)}{\pi (r_s^B)^2} \right) \in [1/4, 1/2).
              \label{eq:eta_boundary_near_corner}
          \end{equation}
    \item \textbf{Interior Node Crossing a Single Wall (Case 2):} At wall distance $d_I < r_s^I$:
          \begin{equation}
              \eta_{I_2} = 1 - \frac{A_{\mathrm{cap}}(d_I, r_s^{I_2})}{\pi (r_s^{I_2})^2} \in (1/2, 1).
              \label{eq:eta_circular_cap}
          \end{equation}
    \item \textbf{Interior Node Near Corner (Case 2B):} Let $d_x, d_y < r_s^I$ denote the normal distances from node $\mathbf{x}_I$ to the two orthogonal corner walls. If the corner vertex lies outside the disk ($d_x^2 + d_y^2 \ge (r_s^I)^2$), then $S_0^I = \pi (r_s^I)^2 - A_{\mathrm{cap}}(d_x, r_s^I) - A_{\mathrm{cap}}(d_y, r_s^I)$. If it lies inside ($d_x^2 + d_y^2 < (r_s^I)^2$), inclusion--exclusion adds the exact corner-quadrant overlap area $S_0^I = \pi (r_s^I)^2 - A_{\mathrm{cap}}(d_x, r_s^I) - A_{\mathrm{cap}}(d_y, r_s^I) + A_{\mathrm{overlap}}^I$, where:
          \begin{equation}
              A_{\mathrm{overlap}}^I = \frac{(r_s^I)^2}{2} \left[ \frac{\pi}{2} - \arcsin\left(\frac{d_x}{r_s^I}\right) - \arcsin\left(\frac{d_y}{r_s^I}\right) \right] - \frac{d_x}{2}\sqrt{(r_s^I)^2 - d_x^2} - \frac{d_y}{2}\sqrt{(r_s^I)^2 - d_y^2} + d_x d_y.
              \label{eq:corner_overlap_formula}
          \end{equation}
\end{itemize}

\subsection{Geometric Moments for Truncated Mass Matrices}
The Cartesian first- and second-order moments over $D_I = \Omega_s^I \cap \Omega$ relative to $\mathbf{x}_I$ required in \eqref{eq:m_truncated_moments} evaluate in closed form as follows:
\begin{itemize}
    \item \textbf{Flat Half-Disk (Case 3) and Orthogonal Corner Quarter-Disk:} For a straight wall with inward unit normal $\vec{n}_{\mathrm{in}} = [n_x, n_y]^T$, $S_{1x}^B = \frac{2}{3} n_x (r_s^B)^3$, $S_{1y}^B = \frac{2}{3} n_y (r_s^B)^3$, $S_{2xx}^B = S_{2yy}^B = \frac{\pi}{8} (r_s^B)^4$, and $S_{2xy}^B = 0$. At an orthogonal convex corner vertex with inward quadrant signs $\sigma_x, \sigma_y \in \{-1, +1\}$, $S_{1x}^B = \frac{1}{3} \sigma_x (r_s^B)^3$, $S_{1y}^B = \frac{1}{3} \sigma_y (r_s^B)^3$, $S_{2xx}^B = S_{2yy}^B = \frac{\pi}{16} (r_s^B)^4$, and $S_{2xy}^B = \frac{1}{8} \sigma_x \sigma_y (r_s^B)^4$.
    \item \textbf{Single-Wall Cap Truncation (Case 2):} For an interior node at distance $d_I \in (0, r_s^I)$ from a straight wall with inward unit normal $\vec{n}_{\mathrm{in}} = [n_x, n_y]^T$ and unit tangent $\vec{t}_w = [-n_y, n_x]^T$, let $\xi_{\parallel} = (\mathbf{x} - \mathbf{x}_I) \cdot \vec{t}_w$ and $\xi_{\perp} = (\mathbf{x} - \mathbf{x}_I) \cdot \vec{n}_{\mathrm{in}}$ denote the local coordinates parallel (tangential) and perpendicular (normal) to the wall, respectively. Define $a_I = \sqrt{(r_s^I)^2 - d_I^2}$ and $\theta_0 = \arccos(d_I / r_s^I)$. In the wall-aligned $(\xi_{\parallel}, \xi_{\perp})$ frame, the first-order moments are $S_{1\parallel}^I = \iint_{D_I} \xi_{\parallel}\,d\Omega = 0$ and $S_{1\perp}^I = \iint_{D_I} \xi_{\perp}\,d\Omega = \frac{2}{3} a_I^3$, while the second-order moments $S_{2\parallel\parallel}^I = \iint_{D_I} \xi_{\parallel}^2\,d\Omega$, $S_{2\perp\perp}^I = \iint_{D_I} \xi_{\perp}^2\,d\Omega$, and $S_{2\parallel\perp}^I = \iint_{D_I} \xi_{\parallel}\xi_{\perp}\,d\Omega = 0$ satisfy:
          \begin{equation}
              S_{2\parallel\parallel}^I = \frac{\pi}{4}(r_s^I)^4 - \left[ \frac{(r_s^I)^4}{4}\theta_0 - \frac{d_I a_I}{12}\left(5(r_s^I)^2 - 2d_I^2\right) \right], \qquad S_{2\perp\perp}^I = \frac{\pi}{4}(r_s^I)^4 - \left[ \frac{(r_s^I)^4}{4}\theta_0 + \frac{d_I a_I}{4}\left((r_s^I)^2 - 2d_I^2\right) \right].
          \end{equation}
          Rotating from $(\xi_{\parallel}, \xi_{\perp})$ back to $(x, y)$ yields the Cartesian moments directly:
          \begin{equation}
              S_{1x}^I = n_x S_{1\perp}^I, \quad S_{1y}^I = n_y S_{1\perp}^I, \quad S_{2xx}^I = n_x^2 S_{2\perp\perp}^I + n_y^2 S_{2\parallel\parallel}^I, \quad S_{2yy}^I = n_y^2 S_{2\perp\perp}^I + n_x^2 S_{2\parallel\parallel}^I, \quad S_{2xy}^I = n_x n_y \left(S_{2\perp\perp}^I - S_{2\parallel\parallel}^I\right).
          \end{equation}
\end{itemize}
For doubly truncated near-corner subdomains (Cases~2B and 3B), these five Cartesian moments are evaluated via autonomous $16 \times 64$ polar Gauss--Legendre quadrature clipped to the active region $D_I = \Omega_s^I \cap \Omega$, without requiring background grids. On general non-rectilinear boundaries such as curved walls or non-orthogonal corners, $S_0^I$ and moments can similarly be evaluated via autonomous disk-confined polar quadrature.

\section{Broken Norms, Variational Consistency, and Convergence Characterization}
\label{subsec:variational_consistency_strang}

Because local $\mathcal{P}^1$ reconstructions reside in $\mathcal{V}_h \not\subset H(\mathbf{curl}, \Omega)$, this appendix formalizes the multiplicity-weighted broken Sobolev norms (Appendix~\ref{subsec:conformity_discussion}), consistency residual (Appendix~\ref{subsubsec:consistency_residual}), algebraic and approximation properties (Appendix~\ref{subsubsec:demonstrated_properties}), and conditional Strang spectral error decomposition (Appendix~\ref{subsubsec:strang_framework}).

\subsection{Multiplicity-Weighted Broken Sobolev Norm and Homogenization of \texorpdfstring{$M(\mathbf{x})$}{M(x)}}
\label{subsec:conformity_discussion}

Discretization errors over the active subdomains $\{D_I\}_{I=1}^N$ are quantified by the \textbf{multiplicity-weighted broken Sobolev norm}:
\begin{equation}
    \|\vec{w}_h\|_{h, \mathbf{curl}}^2 \equiv \frac{1}{\bar{M}} \sum_{I=1}^N \|\vec{w}_{I,h}\|_{\boldsymbol{H}(\mathbf{curl}, D_I)}^2 = \frac{1}{\bar{M}} \sum_{I=1}^N \iint_{D_I} \left[ |\vec{w}_{I,h}|^2 + |\nabla \times \vec{w}_{I,h}|^2 \right] d\Omega,
    \label{eq:rigorous_broken_norm}
\end{equation}
where $M(\mathbf{x})$ and $\bar{M} \approx \pi \alpha_s^2$ are defined in Section~\ref{subsubsec:discrete_space_definition}.

\paragraph*{Asymptotic Homogenization of the Weighted Rayleigh Quotient:}
Because $M(\mathbf{x})/\bar{M} = 1 + \widetilde{M}(\mathbf{x})/\bar{M} \ne 1$, the continuous weighted Rayleigh quotient \eqref{eq:multiplicity_rayleigh_quotient} differs from the unweighted physical quotient $\lambda(\vec{E}) = \frac{\iint_\Omega \mu_r^{-1}|\nabla\times\vec{E}|^2\,d\Omega}{\iint_\Omega \epsilon_r|\vec{E}|^2\,d\Omega}$ whenever $\widetilde{M}(\mathbf{x}) \not\equiv 0$. On rectangular domains discretized by regular Cartesian grids of cell spacing $h_{\mathrm{nom}}$, vertex-centered disks of radius $r_s = \alpha_s h_{\mathrm{nom}} \le h_{\mathrm{nom}}$ yield full interior disks, edge half-disks, and corner quarter-disks with no cross-wall spillover. Every Cartesian cell $Q_{ij} \subset \Omega$ is then covered by four identical quarter-disks at its vertices, so $\bar{M} = \pi\alpha_s^2$ holds exactly and $\widetilde{M}(\mathbf{x}) = M(\mathbf{x}) - \bar{M}$ has vanishing zero-order and reflection-antisymmetric first-order moments ($\iint_{Q_{ij}} \widetilde{M}\,d\Omega = 0$, $\iint_{Q_{ij}} (\mathbf{x}-\mathbf{x}_{ij}^c)\widetilde{M}\,d\Omega = \mathbf{0}$). Expanding any smooth eigenmode density $f \in W^{2,\infty}(\Omega)$ around each cell center $\mathbf{x}_{ij}^c$ yields (provided the normalized denominator stays uniformly bounded away from zero):
\begin{equation}
    \left| \iint_{\Omega} \frac{\widetilde{M}(\mathbf{x})}{\bar{M}} f(\mathbf{x}) \, d\Omega \right| \le C_{\mathrm{hom}} h_{\mathrm{nom}}^2 \|f\|_{W^{2,\infty}(\Omega)} \implies \left| \lambda^{(M)}(\vec{E}) - \lambda(\vec{E}) \right| = \mathcal{O}(h_{\mathrm{nom}}^2),
    \label{eq:multiplicity_homogenization}
\end{equation}
consistent with Section~\ref{subsubsec:multiplicity_rayleigh_effect}. On disordered clouds, deterministic periodic cell cancellation breaks down (worst-case $\mathcal{O}(1)$ shift, decaying sublinearly under i.i.d.\ jitter as quantified in Section~\ref{subsubsec:multiplicity_rayleigh_effect}) unless partition-of-unity normalization $w_I(\mathbf{x}) = \chi_{D_I}(\mathbf{x})/M(\mathbf{x})$ is introduced via lens-intersection quadratures. Note that \eqref{eq:multiplicity_homogenization} bounds only the continuous quotient $\lambda^{(M)}(\vec{E})$; its role in discrete spectral error is analyzed in Appendix~\ref{subsubsec:strang_framework}.

\subsection{Local Green Identity and Global Consistency Residual}
\label{subsubsec:consistency_residual}

For an exact field $\vec{E}$ satisfying the governing differential equation \eqref{eq:helmholtz_pde} in $\Omega$, Green's vector identity \eqref{eq:local_weak_form} on each $D_I$ gives, for any local test function $\vec{W}_{I,h}$:
\begin{equation}
    a_{D_I}(\vec{E}, \vec{W}_{I,h}) - k_0^2 m_{D_I}(\vec{E}, \vec{W}_{I,h}) - \ell_{D_I}^{\mathrm{vol}}(\vec{W}_{I,h}) = \int_{L_{sI}} \vec{\tau}(\vec{E}) \cdot \vec{W}_{I,h} \, ds + \int_{\Gamma_s^I} \vec{\tau}(\vec{E}) \cdot \vec{W}_{I,h} \, ds,
    \label{eq:local_green_flux_balance}
\end{equation}
where $\ell_{D_I}^{\mathrm{vol}}$ contains the prescribed volumetric source contribution. The second integral represents the physical boundary contribution, whereas the first arises from the artificial boundary introduced by the local integration subdomain.

Unlike element-based discretizations, the circular subdomains $\{D_I\}_{I=1}^N$ do not form a disjoint partition of $\Omega$. Their artificial boundaries $L_{sI}$ are generally non-coincident, and the associated local test functions are independently reconstructed from different directional stencils. Consequently, the artificial boundary contributions cannot, in general, be assumed to cancel pairwise.

For the assembled non-conforming formulation, the consistency residual therefore contains the aggregate contribution:
\begin{equation}
    \mathcal{R}_h(\vec{w}_h) = \sum_{I=1}^N \int_{L_{sI}} \vec{\tau}(\vec{E}) \cdot \vec{w}_{I,h} \, ds,
    \label{eq:global_consistency_residual}
\end{equation}
after accounting for the prescribed physical boundary terms. The magnitude and asymptotic behavior of this residual depend on the interaction between the local test functions, the overlap geometry, and the inter-stencil discrepancies. In particular, a scaled trace inequality alone does not establish a vanishing order for the aggregate flux contribution; an additional estimate controlling these interactions would be required.

\subsection{Demonstrated Algebraic, Approximation, and Definiteness Properties}
\label{subsubsec:demonstrated_properties}

Distinguishing the exact local Green identity and the algebraic properties of the assembled operators from an open global asymptotic consistency estimate, four foundational properties are established directly:
\begin{itemize}
    \item \textbf{Exact Affine Reproduction ($\vec{E}^* \in \mathcal{P}^1$):} Defining the exact nodal DoF vector $\mathbf{e}^* \in \mathbb{R}^N$ with entries $e_k^* = \vec{E}^*(\mathbf{x}_k) \cdot \vec{t}_k$, Green's identity \eqref{eq:local_green_flux_balance} ensures $a_{D_I}(\vec{E}^*, \vec{W}_{I,h}) - k_0^2 m_{D_I}(\vec{E}^*, \vec{W}_{I,h}) = \ell_{D_I}(\vec{W}_{I,h})$ on each subdomain when the local load $\ell_{D_I}$ incorporates the complete subdomain boundary flux. Since $\mathcal{P}^1$ reproduces linear fields exactly ($\vec{E}_{I,h} \equiv \vec{E}^*$), every local residual vanishes to machine precision:
          \begin{equation}
              \mathbf{r}_I(\mathbf{e}^*) \equiv (\mathbf{K}_I^{\mathrm{local}} - k_0^2 \mathbf{M}_I^{\mathrm{local}}) \mathbf{L}_I \mathbf{e}^* - \mathbf{f}_I^{\mathrm{local}} \equiv \mathbf{0} \implies \mathbf{R}(\mathbf{e}^*) = \sum_{I=1}^N \mathbf{L}_I^T \mathbf{r}_I(\mathbf{e}^*) \equiv \mathbf{0},
              \label{eq:subdomain_residual_zero}
          \end{equation}
          verifying the local Green identity, polynomial reproduction, and algebraic assembly implementation regardless of overlap multiplicity $M(\mathbf{x})$.
    \item \textbf{Local Taylor Approximation and Overlap Jumps:} For smooth fields $\vec{E} \in [H^2(\Omega)]^2$, suppose each stencil $\Omega_u^I$ has a uniformly bounded dimensionless collocation condition number $\kappa_2(\widehat{\mathbf{A}}_I) \le C_A$ ($\kappa_2(\widehat{\mathbf{A}}_I) \le 5.9$ on the tested clouds) and a controlled support-to-spacing ratio $r_u^I \equiv \max_{k \in \Omega_u^I} \|\mathbf{x}_k - \mathbf{x}_I\|_2 \le C_u h_I$. Then local polynomial interpolation on $D_I$ ($r_s^I = \alpha_s h_I$) satisfies the conditional Taylor interpolation bounds:
          \begin{equation}
              \|\vec{E} - \vec{E}_{I,h}\|_{\boldsymbol{L}^2(D_I)} \le C_1 h_I^2 \|\vec{E}\|_{\boldsymbol{H}^2(\Omega_u^I)}, \qquad
              \|\nabla(\vec{E} - \vec{E}_{I,h})\|_{\boldsymbol{L}^2(D_I)} \le C_2 h_I \|\vec{E}\|_{\boldsymbol{H}^2(\Omega_u^I)},
              \label{eq:local_taylor_bounds}
          \end{equation}
          with $C_1, C_2$ depending only on $\alpha_s, C_A,$ and $C_u$. By the triangle inequality across any overlap $\Omega_s^{IJ} \equiv D_I \cap D_J \ne \emptyset$, the inter-stencil discrepancy is governed locally by $h_{IJ} \equiv \max(h_I, h_J)$. Under the additional \textbf{quasi-uniformity hypothesis} that local spacings and stencil support radii remain comparable to $h_{\mathrm{char}}$ across the refinement family,
          \begin{equation}
              c_0 h_{\mathrm{char}} \le h_I \le r_u^I \le C_0 h_{\mathrm{char}} \qquad \forall I \in \{1, \dots, N\},
              \label{eq:quasi_uniformity_hypothesis}
          \end{equation}
          with refinement-independent constants $0 < c_0 \le C_0 < \infty$ (whereas on non-quasi-uniform clouds with strongly clustered or rarefied regions, $h_{\mathrm{char}}$ below is replaced by the local scale $h_{IJ} = \max(h_I, h_J)$), the overlap jumps scale globally as:
          \begin{align}
              \|\vec{E}_{I,h} - \vec{E}_{J,h}\|_{\boldsymbol{L}^2(\Omega_s^{IJ})} & \le C_3 h_{\mathrm{char}}^2 \|\vec{E}\|_{\boldsymbol{H}^2(\Omega_u^I \cup \Omega_u^J)}, \label{eq:stencil_jump_L2} \\
              \|\operatorname{curl}\vec{E}_{I,h} - \operatorname{curl}\vec{E}_{J,h}\|_{\boldsymbol{L}^2(\Omega_s^{IJ})} & \le C_4 h_{\mathrm{char}} \|\vec{E}\|_{\boldsymbol{H}^2(\Omega_u^I \cup \Omega_u^J)}. \label{eq:stencil_jump_curl}
          \end{align}
    \item \textbf{Norm Equivalence on the Covered Region $\Omega_{\mathrm{cov}}$:} For a locally finite cover whose pointwise overlap multiplicity satisfies $1 \le M_{\min} \le M(\mathbf{x}) \le M_{\max} < \infty$ almost everywhere in $\Omega_{\mathrm{cov}}$ (with $M_{\max}$ denoting the maximum overlap multiplicity of the discretization family), the subdomain-accumulated $L^2$ form satisfies two-sided equivalence on $\Omega_{\mathrm{cov}}$:
          \begin{equation}
              M_{\min} \|\vec{v}\|_{\boldsymbol{L}^2(\Omega_{\mathrm{cov}})}^2 \le \sum_{I=1}^N \|\vec{v}\|_{\boldsymbol{L}^2(D_I)}^2 = \iint_{\Omega_{\mathrm{cov}}} M(\mathbf{x}) |\vec{v}(\mathbf{x})|^2 \, d\Omega \le M_{\max} \|\vec{v}\|_{\boldsymbol{L}^2(\Omega_{\mathrm{cov}})}^2 \quad \forall \vec{v} \in \boldsymbol{L}^2(\Omega_{\mathrm{cov}}).
              \label{eq:norm_equivalence_multiplicity}
          \end{equation}
          On configurations with full coverage ($\Omega_{\mathrm{cov}} = \Omega$, such as regular grids where $\bar{M} \approx \pi\alpha_s^2 \approx 2.01$), Eq.~\eqref{eq:norm_equivalence_multiplicity} establishes global norm equivalence on $\boldsymbol{L}^2(\Omega)$; on perturbed clouds with partial coverage ($\Omega_{\mathrm{cov}} \subset \Omega$), it holds over the covered region $\Omega_{\mathrm{cov}}$.
    \item \textbf{Positive Definiteness on a Fixed Mesh after Stabilization and PEC Reduction ($\mathbf{K}_{\mathrm{red}} > 0$):}
          For $s_{\mathrm{div}} > 0$ and $\alpha_{\mathrm{stab}} > 0$, the combined div-curl and orthogonal stabilization operator $\mathbf{K}_I^{\mathrm{local, stab}}$ \eqref{eq:k_stab_analytic} has a local kernel consisting exclusively of constant vector fields ($\vec{E}_{I, h} = \vec{c}_I \in \mathbb{R}^2$). Indeed, the curl term penalizes $(\nabla\times\vec{E}_{I,h})_z$, the divergence term penalizes $\nabla\cdot\vec{E}_{I,h}$, and $\mathbf{K}_{I, \mathrm{stab}}^{\mathrm{local}}$ penalizes both curl- and divergence-free harmonic gradient modes. Hence $\mathbf{e}^T \mathbf{K} \mathbf{e} = 0$ requires $\mathbf{L}_I \mathbf{e} \in \ker(\mathbf{K}_I^{\mathrm{local, stab}})$ for all $I$, restricting each local field to $\vec{E}_{I, h} = \vec{c}_I$.

          Each shared node $k \in \Omega_u^I \cap \Omega_u^J$ imposes the scalar directional projection $e_k = \vec{c}_I \cdot \vec{t}_k = \vec{c}_J \cdot \vec{t}_k$; stacking the shared unit directors $\{\vec{t}_k^T : k \in \Omega_u^I \cap \Omega_u^J\}$ as the rows of $\mathbf{T}_{IJ} \in \mathbb{R}^{|\Omega_u^I \cap \Omega_u^J| \times 2}$ gives $\mathbf{T}_{IJ}(\vec{c}_I - \vec{c}_J) = \mathbf{0}$. Call $(I, J)$ a full-rank edge when its shared-director matrix has full column rank $\operatorname{rank}(\mathbf{T}_{IJ}) = 2$, equivalently a strictly positive second (smallest) singular value $\sigma_2(\mathbf{T}_{IJ}) > 0$. Under the geometric connectivity condition that the stencil-overlap graph contains a connected spanning subgraph of such full-rank edges, adjacent constant vectors coincide ($\vec{c}_I = \vec{c}_J$) and therefore match globally ($\vec{c}_I = \vec{c} \in \mathbb{R}^2$), yielding $\dim(\ker(\mathbf{K})) = 2$. Along PEC boundaries with non-parallel walls, enforcing $\vec{c} \cdot \vec{t}_1 = \vec{c} \cdot \vec{t}_2 = 0$ forces $\vec{c} = \mathbf{0}$, so eliminating PEC DoFs via \eqref{eq:dirichlet_condensed_system} yields:
          \begin{equation}
              \ker(\mathbf{K}_{\mathrm{red}}) = \{ \mathbf{0} \} \implies \mathbf{K}_{\mathrm{red}} = \mathbf{K}_{\mathrm{red}}^T > 0 \quad \text{(on a fixed connected mesh)}.
              \label{eq:condensed_positive_definite_proof}
          \end{equation}
          We verified this geometric condition directly across all sixteen benchmark point clouds ($15\times 15$ to $81\times 81$ under both $0\%$ and $25\%$ jitter with isotropic random directors, Section~\ref{sec:results}). On every tested cloud, the subgraph formed by full-rank edges connects all $N$ nodes into a single component with $\sigma_2(\mathbf{T}_{IJ}) \ge 1.0 \times 10^{-3}$ (maximum-spanning-tree bottleneck $\sigma_2(\mathbf{T}_{IJ}) \ge 0.395$ across all 16 clouds).
\end{itemize}

\subsection{Conditional Strang Framework and Spectral Error Decomposition}
\label{subsubsec:strang_framework}

The assembled stiffness and mass matrices retain the symmetry and semidefiniteness properties established in Section~\ref{subsubsec:symmetric_assembly}, together with the fixed-mesh positive definiteness $\mathbf{K}_{\mathrm{red}} > 0$ proved in \eqref{eq:condensed_positive_definite_proof}. To formulate the abstract error estimate, let $\mathcal{V}_{h,0} \equiv \{ \vec{v}_h(\mathbf{x}; \mathbf{e}) \in \mathcal{V}_h : \mathbf{e}_{\mathrm{bnd}} = \mathbf{0} \}$ denote the PEC-reduced discrete space parameterized by the interior DoF vector $\mathbf{e}_{\mathrm{red}} \in \mathbb{R}^{N_{\mathrm{int}}}$. On $\mathcal{V}_{h,0}$, the stabilized broken bilinear form $a_h^{\mathrm{stab}}(\vec{v}_h, \vec{w}_h) \equiv a_h(\vec{v}_h, \vec{w}_h) + s_h^{\mathrm{stab}}(\vec{v}_h, \vec{w}_h) = \sum_{I=1}^N [a_{D_I}(\vec{v}_{I,h}, \vec{w}_{I,h}) + a_{D_I}^{\mathrm{stab}}(\vec{v}_{I,h}, \vec{w}_{I,h})]$ induces the \textbf{broken energy norm}:
\begin{equation}
    \begin{aligned}
        \|\vec{v}_h\|_{h, a}^2 \equiv a_h^{\mathrm{stab}}(\vec{v}_h, \vec{v}_h) & = \sum_{I=1}^N \iint_{D_I} \frac{1}{\mu_r} \Big[ |(\nabla \times \vec{v}_{I,h})_z|^2 + s_{\mathrm{div}} |\nabla \cdot \vec{v}_{I,h}|^2 \\
                                                                                & \qquad + 2\alpha_{\mathrm{stab}} \operatorname{tr}\!\big((\nabla_{\mathrm{sym},0}\vec{v}_{I,h})^T \nabla_{\mathrm{sym},0}\vec{v}_{I,h}\big) \Big] d\Omega = \mathbf{e}_{\mathrm{red}}^T \mathbf{K}_{\mathrm{red}} \mathbf{e}_{\mathrm{red}},
    \end{aligned}
    \label{eq:broken_energy_norm_def}
\end{equation}
which combines the curl, penalized divergence, and orthogonal stabilization energies over $\mathcal{V}_{h,0}$. Because $\mathbf{K}_{\mathrm{red}} > 0$ by \eqref{eq:condensed_positive_definite_proof}, $\|\cdot\|_{h,a}$ defines a strict norm on $\mathcal{V}_{h,0}$ on each fixed connected mesh. For the shifted resolvent problem of Section~\ref{subsec:stage3_mms} with $-k_0^2 \mapsto +\kappa > 0$, adding the positive mass term $\kappa m_h(\vec{v}_h, \vec{v}_h) = \kappa \mathbf{e}_{\mathrm{int}}^T \mathbf{M}_{\mathrm{int},\mathrm{int}} \mathbf{e}_{\mathrm{int}} > 0$ defines the shifted energy norm $\|\vec{v}_h\|_{h,a,\kappa}^2 \equiv a_h(\vec{v}_h, \vec{v}_h) + \kappa m_h(\vec{v}_h, \vec{v}_h)$, which is positive-definite on the active discrete space even for $\alpha_{\mathrm{stab}} = 0$ and mixed Dirichlet--Neumann boundaries.

In this non-conforming broken setting, the abstract error for the auxiliary linear problem on $\mathcal{V}_{h,0}$ is governed by Strang's First Lemma \cite{ciarlet2002finite, brenner2008mathematical}:
\begin{equation}
    \|\vec{E} - \vec{E}_h\|_{h, a} \le \left( 1 + \frac{\gamma_a}{\alpha_h} \right) \inf_{\vec{v}_h \in \mathcal{V}_{h,0}} \|\vec{E} - \vec{v}_h\|_{h, a} + \frac{1}{\alpha_h} \sup_{\vec{w}_h \in \mathcal{V}_{h,0} \setminus \{\mathbf{0}\}} \frac{|\mathcal{R}_h(\vec{w}_h)|}{\|\vec{w}_h\|_{h, a}},
    \label{eq:strang_lemma}
\end{equation}
where $\gamma_a > 0$ and $\alpha_h > 0$ denote the continuity and discrete coercivity constants of the active bilinear form.

\paragraph*{Relation Between the SSG Discrete Eigenproblem and the Continuous Maxwell Eigenproblem:}
Let $(\lambda, \vec{E})$ denote a physical continuous eigenpair ($\nabla\times(\mu_r^{-1}\nabla\times\vec{E}) = \lambda\epsilon_r\vec{E}$, $\nabla\cdot\vec{E} = 0$, $\vec{E}\cdot\vec{t} = 0$ on $\Gamma_{\mathrm{PEC}}$). To clarify why the discrete eigenvalues $\lambda_h$ of the stabilized SSG operator $a_h^{\mathrm{stab}}(\vec{E}_h, \vec{w}_h) = \lambda_h m_h(\vec{E}_h, \vec{w}_h)$ can approximate $(\lambda, \vec{E})$---and what remains unproven---we expand the broken bilinear error $a_h(\vec{E}_h - \vec{E}, \vec{E}_h - \vec{E}) - \lambda m_h(\vec{E}_h - \vec{E}, \vec{E}_h - \vec{E})$ using $a_h^{\mathrm{stab}}(\vec{E}_h, \vec{E}_h) = \lambda_h m_h(\vec{E}_h, \vec{E}_h)$ and $a_h(\vec{E}, \vec{E}) = \lambda^{(M)}(\vec{E})\,m_h(\vec{E}, \vec{E})$. This yields the exact algebraic identity:
\begin{equation}
    \begin{aligned}
        (\lambda_h - \lambda)\,m_h(\vec{E}_h, \vec{E}_h)
        & = \underbrace{a_h(\vec{E}_h - \vec{E}, \vec{E}_h - \vec{E}) - \lambda m_h(\vec{E}_h - \vec{E}, \vec{E}_h - \vec{E})}_{\text{(I) } \mathcal{P}^1 \text{ field/curl approximation error}}
        - \underbrace{\big(\lambda^{(M)}(\vec{E}) - \lambda\big)\,m_h(\vec{E}, \vec{E})}_{\text{(II) spatial multiplicity shift}} \\
        & \quad + \underbrace{2\,\mathcal{R}_h^{\mathrm{eig}}(\vec{E}, \vec{E}_h)}_{\text{(IIIa) non-conformity over } L_{sI}}
        + \underbrace{s_h^{\mathrm{stab}}(\vec{E}_h, \vec{E}_h)}_{\text{(IIIb) stabilization bias } \mathcal{O}(\alpha_{\mathrm{stab}})},
    \end{aligned}
    \label{eq:spectral_error_decomposition}
\end{equation}
where $\mathcal{R}_h^{\mathrm{eig}}(\vec{E}, \vec{E}_h) \equiv a_h(\vec{E}, \vec{E}_h) - \lambda m_h(\vec{E}, \vec{E}_h) = \sum_{I=1}^N \int_{\partial D_I} \vec{\tau}(\vec{E}) \cdot \vec{E}_{I,h} \, ds$ is the non-conforming boundary consistency pairing evaluated on the discrete eigenmode $\vec{E}_h$ (with $\int_{\Gamma_s^I} \vec{\tau}(\vec{E})\cdot\vec{E}_{I,h}\,ds$ reflecting any non-conforming trace discrepancy along $\Gamma_{\mathrm{PEC}}$ in addition to the internal arcs $L_{sI}$). Equation~\eqref{eq:spectral_error_decomposition} separates the distinct contributions to the discrete eigenvalue error, although their asymptotic orders are not established uniformly by the identity itself:
\begin{itemize}
    \item \emph{Term (I) (Approximation and Penalized Energy Error):} The local interpolation estimates \eqref{eq:local_taylor_bounds} provide an $\mathcal{O}(h_{\mathrm{char}}^2)$ approximation bound in $L^2$ and an $\mathcal{O}(h_{\mathrm{char}})$ bound for first derivatives under the stated regularity, stencil-conditioning, and quasi-uniformity hypotheses. However, because the divergence penalty is resolution-dependent ($s_{\mathrm{div}} \propto (h_{\mathrm{char}}/L)^{-1.5}$), \eqref{eq:local_taylor_bounds} alone yields at best $s_{\mathrm{div}} \sum_I \|\nabla \cdot (\vec{E}_{I,h} - \vec{E})\|_{L^2(D_I)}^2 \lesssim (h_{\mathrm{char}}/L)^{-1.5} h_{\mathrm{char}}^2 \propto h_{\mathrm{char}}^{0.5}$ and does not establish an $\mathcal{O}(h_{\mathrm{char}}^2)$ bound for the penalized divergence contribution in Term~(I). A sharper estimate controlling the divergence of the discrete eigenmodes under the penalty scaling would be required to quantify this contribution; accordingly, quadratic Rayleigh-quotient error behavior associated with Term~(I) is interpreted conditionally on the convergence and additional divergence control of the discrete eigenfunctions.
    \item \emph{Term (II) (Spatial Multiplicity Shift):} Controlled at $\mathcal{O}(h_{\mathrm{nom}}^2)$ on regular Cartesian grids via \eqref{eq:multiplicity_homogenization}, whereas deterministic periodic cancellation fails under non-periodic jitter (leaving the slowly decaying $0.63\%\text{--}1.33\%$ floor across $49^2\text{--}81^2$, Section~\ref{subsubsec:multiplicity_rayleigh_effect}).
    \item \emph{Terms (IIIa)--(IIIb) (Non-Conformity and Orthogonal Stabilization):} The non-conforming pairing $\mathcal{R}_h^{\mathrm{eig}}(\vec{E}, \vec{E}_h)$ lacks a proven asymptotic bound on general overlapping clouds, whereas Term~(IIIb) scales as $\mathcal{O}(h_{\mathrm{char}}^{\gamma})$ for bounded discrete modes under \eqref{eq:alpha_stab_vms_scaling}---i.e., $\mathcal{O}(h_{\mathrm{char}})$ for the adopted linear calibration ($\gamma = 1$, $c_{\mathrm{stab}} = 4.0$). Consequently, while linear stabilization yields steeper effective slopes in the tested engineering range ($p_{15\text{--}81} = 1.72$, $p_{36\text{--}81} = 2.48$ on regular grids), the formal asymptotic eigenvalue error order as $h_{\mathrm{char}} \to 0$ is bounded by $\mathcal{O}(h_{\mathrm{char}})$ unless $\gamma > 1$ or a parameter-free mixed formulation is adopted.
\end{itemize}

\paragraph*{Hypotheses Governing the Numerical Interpretation of Spectral Convergence:}
Accordingly, the spectral convergence behavior established in Section~\ref{subsec:stage5_cavity} ($p = 1.72$ regular, $p = 1.61$ at $25\%$ jitter) constitutes \textbf{systematic numerical evidence} interpreted under four working conditions. Condition (H1) assumes $H^2$-regularity of physical cavity modes (satisfied on convex rectilinear domains without re-entrant corners). Condition (H2) assumes uniformly well-conditioned stencils ($\kappa_2(\widehat{\mathbf{A}}_I) \le C_A$) and quasi-uniform spacing \eqref{eq:quasi_uniformity_hypothesis}. Condition (H3) relies on periodic disk coverage on regular grids ($\Omega_{\mathrm{cov}} = \Omega$, $\alpha_s \in [\sqrt{2}/2, 1]$) supporting \eqref{eq:multiplicity_homogenization} (inactive under $25\%$ jitter). Finally, (H4) invokes full-rank stencil-overlap connectivity verified numerically via \eqref{eq:condensed_positive_definite_proof} for fixed-mesh $\mathbf{K}_{\mathrm{red}} > 0$, alongside the empirical shifting of spurious gradient modes outside the 10-mode operational band under the calibrated scalings ($s_{\mathrm{div}} \propto (h_{\mathrm{char}}/L)^{-1.5}$, $\alpha_{\mathrm{stab}} = c_{\mathrm{stab}}\,h_{\mathrm{char}}/L$). Establishing uniform discrete coercivity, asymptotic smallness of $\mathcal{R}_h^{\mathrm{eig}}$, and discrete compactness in $H(\mathbf{curl})$ \cite{boffi2006finite, boffi2013mixed} remains open. Whereas Stage~3 verifies local approximation under exact local equilibrium ($\mathcal{R}_h \equiv 0$), the target eigenproblem convergence rests on the numerical evidence in Section~\ref{subsec:stage5_cavity}.

\bmsection*{Conflict of Interest Statement}
The authors declare that they have no conflict of interest.

\bmsection*{Data Availability Statement}
The complete implementation code, generation scripts for nodal point clouds and vector directors, parametric sweep routines for subdomain sizing and orthogonal energy stabilization, as well as reproduction scripts for all tables and figures presented in this study, are openly available in the GitHub repository: \url{https://github.com/rencmbr/VNMM-MLPG2d-IJNME}.

\bmsection*{Declaration of Generative AI and AI-Assisted Technologies}
During the preparation and revision of this work, the authors used Google Antigravity and Google Gemini as AI-assisted development and writing tools across the following activities: (i) drafting, condensing, and refining the technical text, LaTeX formatting, tables, and bibliographic consistency; and (ii) assisting in the development, refactoring, and verification of Python implementation scripts, numerical benchmarks, and figure-generation routines. All the text, numerical codes, simulation results and citations were independently verified, edited, and validated by the authors, who take full and sole responsibility for the originality, integrity, accuracy, and content of this publication.

%% =========================================================================
%% REFERENCES
%% =========================================================================
\bibliography{references}

\end{document}
